% Les composés aromatiques — Chimie 1ère C — Cours signé OnBuch+
% Programme officiel MINESEC, Première C, D, E, TI. Compiler avec : tectonic N04-les-composes-aromatiques.tex
\newcommand{\DOCMATIERE}{Chimie}
\newcommand{\DOCNIVEAU}{1ère C}
\newcommand{\DOCMODULE}{Module 1 : Chimie organique}
\newcommand{\DOCLECON}{04}
\newcommand{\DOCTITRE}{Les composés aromatiques}
% OnBuch+ — préambule « cours signés » ENRICHI (compilé avec tectonic / XeLaTeX)
% Système visuel « Pop » d'OnBuch+ : fond crème (jamais blanc), bordures encre
% épaisses, ombres « dures » décalées, titres Archivo Black en capitales.
% Version enrichie pour les cours de chimie : chimie (mhchem, chemfig),
% graphes (pgfplots), boîtes pédagogiques supplémentaires (définition,
% propriété, méthode, attention, le savais-tu, exercice résolu, exercice,
% TP, bilan), page de garde refaite et signature OnBuch+ ancrée en bas.
%
% Macros attendues dans main.tex AVANT \input{preamble} :
%   \DOCMATIERE  \DOCNIVEAU  \DOCMODULE  \DOCLECON (numéro)  \DOCTITRE
\documentclass[11pt]{article}

\usepackage{fontspec}
\usepackage[french]{babel}
\usepackage[a4paper,margin=2cm,top=3.4cm,bottom=2.8cm,headheight=1.4cm,footskip=1.3cm]{geometry}
\usepackage{amsmath,amssymb,amsfonts}
\usepackage[table]{xcolor} % [table] : fournit aussi \rowcolors (alternance de lignes)
\usepackage{pagecolor}
\usepackage{tikz}
\usetikzlibrary{arrows.meta,positioning,calc,shapes.geometric,shapes.misc,shapes.arrows,decorations.pathmorphing,decorations.markings,patterns,shadows,mindmap,trees,fit,backgrounds,matrix,intersections}
\usepackage{pgfplots}
\pgfplotsset{compat=1.18}
\usepackage[version=4]{mhchem}
\usepackage{chemfig}
\usepackage{enumitem}
\usepackage{fancyhdr}
% [most] charge toutes les bibliothèques tcolorbox (listings, minted, raster,
% documentation...) et gonfle inutilement la mémoire TeX sur un document long
% (~300 boîtes) : on ne charge que ce qui est réellement utilisé.
\usepackage[skins,breakable]{tcolorbox}
\usepackage{graphicx}
\usepackage{colortbl}
\usepackage{booktabs}
\usepackage{tabularx}
\usepackage{array}
\usepackage{multicol}
\usepackage{siunitx}
% parse-numbers=false : accepte aussi \SI{2,5 \times 10^{-3}}{...}
\sisetup{locale=FR,output-decimal-marker={,},group-minimum-digits=4,parse-numbers=false}
\DeclareSIUnit\ohm{\ensuremath{\symup{\Omega}}} % U+2126 absent de la police
\usepackage{qrcode}
% \degree : commande fréquemment utilisée par les modèles pour un angle ou une
% température en dehors de \SI{}{} (non fournie par les paquets chargés).
\providecommand{\degree}{^{\circ}}
\usepackage{lastpage}
\usepackage{hyperref}
\hypersetup{hidelinks}

% ============================================================
% Palette « Pop » OnBuch+ — mêmes hex que la classe P (plus_ui.dart)
% ============================================================
\definecolor{poporange}{HTML}{F59321}
\definecolor{popdark}{HTML}{E07A0C}
\definecolor{popcream}{HTML}{FFF6E8}
\definecolor{popink}{HTML}{1C1714}
\definecolor{poppurple}{HTML}{7A5AE0}
\definecolor{popgreen}{HTML}{2E9E6A}
\definecolor{poppink}{HTML}{E0457B}
\definecolor{popblue}{HTML}{2F7DE1}
\definecolor{popgold}{HTML}{C98A12}
\definecolor{popmuted}{HTML}{8A7F76}
\definecolor{popline}{HTML}{EADFD2}
\definecolor{popgray}{HTML}{8A7F76}
\colorlet{popgrey}{popgray}
% Alias tolérants (le modèle nomme parfois les couleurs différemment)
\colorlet{popwhite}{white}
\colorlet{popblack}{popink}
\colorlet{popred}{poppink}
\colorlet{popyellow}{popgold}
% Teintes claires pour fonds de tableaux / zones
\colorlet{popblueL}{popblue!10!white}
\colorlet{poporangeL}{poporange!14!white}
\colorlet{popgreenL}{popgreen!12!white}
\colorlet{poppurpleL}{poppurple!12!white}
\colorlet{poppinkL}{poppink!12!white}
\colorlet{popgoldL}{popgold!14!white}

\pagecolor{popcream}

% ============================================================
% Typographie
% ============================================================
\defaultfontfeatures{Ligatures=TeX}
\setmainfont{PlusJakartaSans-Medium.ttf}[Path=../../../fonts/,BoldFont=PlusJakartaSans-Bold.ttf]
% Maths en sans-serif (Fira Math) avec chiffres et lettres droites de Plus
% Jakarta, pour que les formules se fondent dans le texte.
\usepackage[math-style=ISO,bold-style=ISO]{unicode-math}
\setmathfont{FiraMath-Regular.otf}[Path=../../../fonts/]
\setmathfont{PlusJakartaSans-Medium.ttf}[Path=../../../fonts/,range={up/num,up/{Latin,latin},\mathup}]
\usepackage{icomma} % virgule décimale sans espace en mode math
% Caractères absents de la police de texte (sinon carré vide) : rendus par
% leur équivalent mathématique, en texte comme en formule.
\usepackage{newunicodechar}
\newunicodechar{α}{\ensuremath{\alpha}}
\newunicodechar{Α}{\ensuremath{\alpha}}
\newunicodechar{β}{\ensuremath{\beta}}
\newunicodechar{δ}{\ensuremath{\delta}}
\newunicodechar{✓}{\popcheck{popgreen}}
\newfontfamily\popdisplay{ArchivoBlack-Regular.ttf}[Path=../../../fonts/]
\newfontfamily\poplogo{SpaceGrotesk-Bold.ttf}[Path=../../../fonts/]
\newfontfamily\popsemi{PlusJakartaSans-SemiBold.ttf}[Path=../../../fonts/]
\newfontfamily\popxbold{PlusJakartaSans-ExtraBold.ttf}[Path=../../../fonts/]
\newfontfamily\popxboldls{PlusJakartaSans-ExtraBold.ttf}[Path=../../../fonts/,LetterSpace=3.0]

\setlist[enumerate]{leftmargin=1.7em,itemsep=5pt,topsep=4pt}
\setlist[itemize]{leftmargin=1.7em,itemsep=4pt,topsep=4pt,label={\color{poporange}\textbullet}}
\setlength{\parindent}{0pt}
\setlength{\parskip}{5pt}
\renewcommand{\arraystretch}{1.25}

% chemfig : liaisons un peu plus épaisses, encre Pop
\setchemfig{atom sep=2em,bond style={line width=0.9pt,popink},cram width=3pt}

% pgfplots : style Pop par défaut
\pgfplotsset{
  every axis/.append style={axis line style={popink,line width=1pt},tick style={popink},
    grid style={popline,line width=0.5pt},label style={font=\small},tick label style={font=\footnotesize}},
  popcurve/.style={poporange,line width=1.8pt,no markers,smooth},
}

% ============================================================
% Pill / squiggle
% ============================================================
\newcommand{\poppill}[3]{%
  \tikz[baseline=-0.55ex]\node[draw=popink,line width=1.2pt,rounded corners=2.6pt,%
    fill=#2,inner xsep=6.5pt,inner ysep=3pt]{%
    \popxboldls\fontsize{7}{7}\selectfont\color{#3}\MakeUppercase{#1}};%
}
\newcommand{\popsquiggle}[1]{%
  \tikz[baseline=-0.4ex]{
    \draw[#1,line width=2.6pt,line cap=round]
      plot [smooth] coordinates {(0,0.09) (0.34,0.01) (0.68,0.21) (1.02,0.01) (1.36,0.10)};
  }%
}
% Mot-clé surligné (marqueur orange sous le texte)
\newcommand{\cle}[1]{{\popxbold\color{popdark}#1}}

% ============================================================
% En-tête / pied
% ============================================================
\pagestyle{fancy}
\fancyhf{}
\renewcommand{\headrulewidth}{0pt}
\renewcommand{\footrulewidth}{0pt}
\lhead{\raisebox{-0.15ex}{\poplogo\fontsize{13}{13}\selectfont\color{popink}OnBuch\color{poporange}+}}
\rhead{\poppill{\DOCMATIERE\ \textbullet\ \DOCNIVEAU\ \textbullet\ Leçon \DOCLECON}{white}{popink}}
\lfoot{\popxbold\fontsize{7.6}{7.6}\selectfont\color{popmuted}\MakeUppercase{Cours signé OnBuch+}\ \textbullet\ {\popsemi\DOCTITRE}}
\rfoot{\tikz[baseline=-0.6ex]\node[rounded corners=8.5pt,fill=popink,minimum height=17pt,minimum width=17pt,inner xsep=5pt,inner sep=0]{\popxbold\fontsize{8}{8}\selectfont\color{popcream}\thepage\,/\,\pageref*{LastPage}};}

% ============================================================
% Titres
% ============================================================
\newcommand{\chaptitre}[2]{%
  \begin{tcolorbox}[enhanced,colback=poporange,colframe=popink,boxrule=2.6pt,arc=11pt,
      drop shadow=popink,left=15pt,right=15pt,top=12pt,bottom=11pt,before skip=3pt,after skip=18pt]
    \poppill{#2}{popink}{popcream}\par\vspace{9pt}
    {\raggedright\hyphenpenalty=10000\exhyphenpenalty=10000\popdisplay\fontsize{22}{24}\selectfont\color{popcream}\MakeUppercase{#1}\par}
  \end{tcolorbox}
}
% Section numérotée : pastille orange avec le numéro + titre Archivo
\newcounter{popsec}
\newcommand{\coursec}[1]{%
  \stepcounter{popsec}\setcounter{popsub}{0}%
  \par\vspace{16pt}\noindent
  \begin{minipage}{\linewidth}
  \tikz[baseline=-0.7ex]\node[draw=popink,line width=1.6pt,fill=poporange,rounded corners=4pt,minimum size=22pt,inner sep=0]{\popdisplay\fontsize{12}{12}\selectfont\color{popcream}\thepopsec};%
  \hspace{8pt}{\popdisplay\fontsize{15}{17}\selectfont\color{popink}\MakeUppercase{#1}}\par
  \vspace{4pt}\noindent{\color{poporange}\rule{36pt}{3.6pt}}
  \end{minipage}\par\nopagebreak\vspace{8pt}
}
\newcounter{popsub}
\newcommand{\courssub}[1]{%
  \stepcounter{popsub}%
  \par\vspace{9pt}\noindent{\popxbold\fontsize{11.5}{13}\selectfont\color{popink}{\color{poporange}\thepopsec.\thepopsub}\hspace{6pt}#1}\par\nopagebreak\vspace{4pt}
}

% ============================================================
% Boîtes pédagogiques — même recette : fond blanc, bordure épaisse colorée,
% ombre dure encre, badge plein. Titre optionnel [#1].
% ============================================================
\tcbset{popbase/.style={breakable,enhanced,colback=white,boxrule=2.3pt,arc=9pt,
  shadow={3pt}{-3pt}{0pt}{popink},left=14pt,right=14pt,top=11pt,bottom=11pt,before skip=11pt,after skip=15pt,
  fontupper=\popsemi\fontsize{10.6}{14.5}\selectfont\color{popink},
  fontlower=\popsemi\fontsize{10.6}{14.5}\selectfont\color{popink}}}
% Coche dessinée (le glyphe ✓ n'existe pas dans les polices du document)
\newcommand{\popcheck}[1]{\tikz[baseline=-0.5ex]\draw[#1,line width=1.6pt,line cap=round,line join=round](-0.13,0.0)--(-0.04,-0.09)--(0.13,0.1);}
\newcommand{\popboxhead}[3]{\poppill{#1}{#2}{white}\ifx\relax#3\relax\else\hspace{6pt}{\popxbold\fontsize{10.6}{12}\selectfont\color{popink}#3}\fi\par\vspace{7pt}}

\newtcolorbox{aretenir}[1][]{popbase,colframe=popgreen,before upper={\popboxhead{À retenir}{popgreen}{#1}}}
\newtcolorbox{exemplebox}[1][]{popbase,colframe=poporange,before upper={\popboxhead{Exemple}{poporange}{#1}}}
\newtcolorbox{definition}[1][]{popbase,colframe=popblue,before upper={\popboxhead{Définition}{popblue}{#1}}}
\newtcolorbox{propriete}[1][]{popbase,colframe=poppurple,before upper={\popboxhead{Propriété}{poppurple}{#1}}}
\newtcolorbox{methode}[1][]{popbase,colframe=popgold,before upper={\popboxhead{Méthode}{popgold}{#1}}}
\newtcolorbox{attention}[1][]{popbase,colframe=poppink,before upper={\popboxhead{Attention}{poppink}{#1}}}
\newtcolorbox{experience}[1][]{popbase,colframe=popblue,colback=popblueL,before upper={\popboxhead{Expérience}{popblue}{#1}}}
% « Le savais-tu ? » : carte violette pleine (contexte, culture, Cameroun)
\newtcolorbox{savaistu}[1][]{popbase,colback=poppurple,colframe=popink,
  fontupper=\popsemi\fontsize{10.4}{14.2}\selectfont\color{white},
  before upper={\poppill{Le savais-tu\,?}{white}{poppurple}\ifx\relax#1\relax\else\hspace{6pt}{\popxbold\color{white}#1}\fi\par\vspace{7pt}}}
% Exercice résolu : énoncé en haut, \tcblower puis la solution sur fond crème
\newcounter{popexo}\newcounter{popexr}
\newtcolorbox{exoresolu}[1][]{popbase,colframe=poporange,colbacklower=popcream,
  segmentation style={popink,line width=1.2pt,dashed},
  before upper={\stepcounter{popexr}\popboxhead{Exercice résolu \thepopexr}{poporange}{#1}},
  before lower={\poppill{Solution}{popink}{popcream}\par\vspace{6pt}}}
% Exercice (non résolu en place) — la correction est dans la partie Corrigés
\newtcolorbox{exercice}[1][]{popbase,colframe=popink,
  before upper={\stepcounter{popexo}\popboxhead{Exercice \thepopexo}{popink}{#1}}}
% Corrigé
\newtcolorbox{corrige}[1][]{popbase,colframe=popgreen,colback=popgreenL,
  before upper={\popboxhead{Corrigé}{popgreen}{#1}}}
% Objectifs / prérequis / bilan
\newtcolorbox{objectifs}{popbase,colframe=popink,colback=poporangeL,
  before upper={\popboxhead{Objectifs de la leçon}{popink}{}}}
\newtcolorbox{prerequis}{popbase,colframe=popink,colback=popblueL,
  before upper={\popboxhead{Prérequis}{popblue}{}}}
\newtcolorbox{bilan}{popbase,colframe=popink,colback=white,boxrule=2.8pt,
  before upper={\popboxhead{Fiche bilan}{poporange}{L'essentiel à connaître pour l'examen}}}
% Figure encadrée avec légende
\newtcolorbox{popfigure}[1][]{enhanced,breakable=false,colback=white,colframe=popline,boxrule=1.4pt,arc=8pt,
  left=10pt,right=10pt,top=10pt,bottom=8pt,before skip=10pt,after skip=12pt,halign=center,
  lower separated=false,halign lower=center,
  fontlower=\popsemi\fontsize{9}{11.5}\selectfont\color{popmuted},
  #1}
\newcommand{\legende}[1]{\par\vspace{6pt}{\popsemi\fontsize{9}{11.5}\selectfont\color{popmuted}#1}}

% ============================================================
% Signature OnBuch+
% ============================================================
\newcommand{\poplogoinline}[1]{{\poplogo\fontsize{#1}{#1}\selectfont\color{popink}OnBuch\color{poporange}+}}
\newcommand{\signatureonbuch}{%
  \begin{tcolorbox}[enhanced,colback=popink,colframe=popink,boxrule=2.2pt,arc=10pt,
      drop shadow=poporange,left=14pt,right=14pt,top=10pt,bottom=10pt,before skip=0pt,after skip=0pt]
    \tikz[baseline=-0.7ex]\node[circle,fill=popgreen,draw=popcream,line width=1.3pt,minimum size=22pt,inner sep=0]{\popcheck{white}};%
    \hspace{9pt}\begin{minipage}[c]{0.62\linewidth}
      {\popxbold\fontsize{12}{13}\selectfont\color{popcream}Signé\ \textbullet\ {\poplogo OnBuch\color{poporange}+}}\par\vspace{2pt}
      {\raggedright\popsemi\fontsize{8}{10.5}\selectfont\color{popline}\MakeUppercase{Équipe pédagogique OnBuch+}\\\MakeUppercase{Conforme au programme officiel MINESEC}\par}
    \end{minipage}\hfill
    {\poplogo\fontsize{22}{22}\selectfont\color{popcream}OnBuch\color{poporange}+}
  \end{tcolorbox}%
}

% ============================================================
% Page de garde
% #1 titre · #2 matière · #3 niveau · #4 module · #5 n° leçon · #6 durée ·
% #7 description / objectifs (texte ou itemize)
% ============================================================
\newcommand{\pagegarde}[7]{%
  \thispagestyle{empty}
  \newgeometry{a4paper,margin=1.7cm,top=1.6cm,bottom=1.4cm}%
  \begin{tikzpicture}[remember picture,overlay]
    \fill[poporange!18] ([xshift=-3.2cm,yshift=-3.6cm]current page.north east) circle (4.6cm);
    \fill[poppurple!14] ([xshift=2.4cm,yshift=7.6cm]current page.south west) circle (3.8cm);
  \end{tikzpicture}%
  {\poplogo\fontsize{30}{30}\selectfont\color{popink}OnBuch\color{poporange}+}\hfill
  \raisebox{0.6ex}{\poppill{Cours signé \textbullet\ Premium}{popink}{popcream}}\par
  \vspace{1pt}\popsquiggle{poppurple}\par
  \vspace{8pt}
  \begin{tcolorbox}[enhanced,colback=poporange,colframe=popink,boxrule=2.8pt,arc=13pt,
      drop shadow=popink,left=18pt,right=18pt,top=12pt,bottom=13pt,after skip=10pt]
    \poppill{#2 \textbullet\ #3}{popink}{popcream}\hspace{5pt}\poppill{#4}{white}{popink}\par\vspace{8pt}
    {\popdisplay\fontsize{14}{14}\selectfont\color{popink}LEÇON #5}\par\vspace{4pt}
    {\raggedright\hyphenpenalty=10000\exhyphenpenalty=10000\popdisplay\fontsize{25}{27}\selectfont\color{popcream}\MakeUppercase{#1}\par}
    \vspace{8pt}
    {\popxbold\fontsize{9.5}{9.5}\selectfont\color{popink}\MakeUppercase{Durée conseillée : #6}}
  \end{tcolorbox}
  \begin{tcolorbox}[enhanced,colback=white,colframe=popink,boxrule=2.2pt,arc=10pt,
      drop shadow=popink,left=16pt,right=16pt,top=10pt,bottom=10pt,after skip=8pt]
    \poppill{Dans cette leçon}{poppurple}{white}\par\vspace{5pt}
    \popsemi\fontsize{9.3}{12.6}\selectfont\color{popink}#7
  \end{tcolorbox}
  \noindent\begin{minipage}[t]{0.487\linewidth}
    \begin{tcolorbox}[enhanced,colback=popgreen,colframe=popink,boxrule=2.2pt,arc=10pt,
        drop shadow=popink,left=11pt,right=11pt,top=10pt,bottom=10pt,height=3.1cm,valign=center]
      \tcbox[colback=white,colframe=popink,boxrule=1.3pt,arc=5pt,boxsep=0pt,nobeforeafter,
        left=3pt,right=3pt,top=3pt,bottom=3pt,valign=center]{\qrcode[height=1.9cm]{https://play.google.com/store/apps/details?id=com.luvvix.onbuch&hl=fr}}%
      \hspace{8pt}\begin{minipage}[c]{0.5\linewidth}
        {\popdisplay\fontsize{11}{12.5}\selectfont\color{white}\MakeUppercase{Télécharge OnBuch}}\par\vspace{4pt}
        {\raggedright\popsemi\fontsize{8.4}{11}\selectfont\color{white}Gratuit \textbullet\ Google Play\\Tuteur IA, annales, concours, résultats\par}
      \end{minipage}
    \end{tcolorbox}
  \end{minipage}\hfill
  \begin{minipage}[t]{0.487\linewidth}
    \begin{tcolorbox}[enhanced,colback=poppurple,colframe=popink,boxrule=2.2pt,arc=10pt,
        drop shadow=popink,left=11pt,right=11pt,top=10pt,bottom=10pt,height=3.1cm,valign=center]
      \tikz[baseline=-0.7ex]\node[circle,fill=white,draw=popink,line width=1.3pt,minimum size=1.6cm,inner sep=0]{\popdisplay\fontsize{24}{24}\selectfont\color{poppurple}+};%
      \hspace{8pt}\begin{minipage}[c]{0.6\linewidth}
        {\popdisplay\fontsize{11}{12.5}\selectfont\color{white}\MakeUppercase{Passe à OnBuch+}}\par\vspace{4pt}
        {\raggedright\popsemi\fontsize{8.4}{11}\selectfont\color{white}Tous les cours signés, Tuteur IA illimité, fiches PDF — onglet OnBuch+ de l'app\par}
      \end{minipage}
    \end{tcolorbox}
  \end{minipage}
  \vspace{8pt}
  \popcontenu
  \vfill
  \signatureonbuch
  \restoregeometry
  \newpage
}

% Rangée « Ce cours contient » (page de garde)
\newcommand{\poptile}[3]{% #1 couleur #2 gros texte #3 légende
  \begin{tcolorbox}[enhanced,colback=#1,colframe=popink,boxrule=2pt,arc=9pt,shadow={2.5pt}{-2.5pt}{0pt}{popink},
      left=6pt,right=6pt,top=9pt,bottom=9pt,halign=center,valign=center,height=1.75cm,width=\linewidth,nobeforeafter]
    {\popdisplay\fontsize{12.5}{13}\selectfont\color{white}\MakeUppercase{#2}}\par\vspace{3pt}
    {\centering\popsemi\fontsize{7.8}{9.5}\selectfont\color{white}#3\par}
  \end{tcolorbox}}
\newcommand{\popcontenu}{%
  \par\noindent{\popxboldls\fontsize{7.5}{7.5}\selectfont\color{popmuted}CE COURS SIGNÉ CONTIENT}\par\vspace{7pt}
  \noindent\begin{minipage}[t]{0.235\linewidth}\poptile{popblue}{Cours}{complet, illustré, pas à pas}\end{minipage}\hfill
  \begin{minipage}[t]{0.235\linewidth}\poptile{poporange}{Méthodes}{et exercices résolus}\end{minipage}\hfill
  \begin{minipage}[t]{0.235\linewidth}\poptile{poppink}{Exercices}{type examen + corrigés}\end{minipage}\hfill
  \begin{minipage}[t]{0.235\linewidth}\poptile{popgreen}{Fiche bilan}{l'essentiel à retenir}\end{minipage}\par}

% Sommaire
\newtcolorbox{sommaire}{enhanced,colback=white,colframe=popink,boxrule=2.2pt,arc=10pt,
  drop shadow=popink,left=18pt,right=18pt,top=13pt,bottom=13pt,before skip=6pt,after skip=20pt,
  before upper={\poppill{Sommaire}{poppurple}{white}\par\vspace{9pt}\popsemi\fontsize{10.6}{14.5}\selectfont\color{popink}}}

% Signature de fin de cours — poussée tout en bas de la dernière page
\newcommand{\findecours}{%
  \par\vfill\nopagebreak
  \begin{center}\popsquiggle{poporange}\end{center}\vspace{4pt}
  \signatureonbuch
}
% Glyphes absents de Fira Math : redéfinis à partir de glyphes disponibles
\AtBeginDocument{%
  \def\setminus{\mathbin{\backslash}}\let\smallsetminus\setminus
  \def\vdots{\mathord{\vbox{\baselineskip3.2pt\lineskiplimit0pt\kern4pt\hbox{.}\hbox{.}\hbox{.}}}}%
  \def\ddots{\mathinner{\mkern1mu\raise6.5pt\hbox{.}\mkern2mu\raise3.6pt\hbox{.}\mkern2mu\raise0.7pt\hbox{.}\mkern1mu}}%
  \def\triangle{\mathord{\scalebox{0.9}{$\symup{\Delta}$}}}%
  \def\nabla{\mathord{\raisebox{\depth}{\scalebox{1}[-1]{$\symup{\Delta}$}}}}%
  \def\checkmark{\popcheck{popdark}}%
  \def\star{\mathbin{\ast}}\def\bigstar{\mathord{\ast}}%
  \def\diamond{\mathbin{\scalebox{0.6}{$\Diamond$}}}%
  \def\bigcup{\mathop{\mathchoice{\raisebox{-0.3em}{\scalebox{1.7}{$\cup$}}}{\scalebox{1.25}{$\cup$}}{\cup}{\cup}}\displaylimits}%
  \def\bigcap{\mathop{\mathchoice{\raisebox{-0.3em}{\scalebox{1.7}{$\cap$}}}{\scalebox{1.25}{$\cap$}}{\cap}{\cap}}\displaylimits}%
  \def\lesseqqgtr{\mathrel{\lessgtr}}\def\bigoplus{\mathop{\oplus}\displaylimits}%
}
\providecommand{\ssi}{si et seulement si} % abréviation fréquente

%% FIGFIX-BEGIN v3 — correctifs globaux des figures (docs/onbuch-generation/tools/figfix)
\usepackage{adjustbox}\usepackage{etoolbox}
% toute figure TikZ/pgfplots plus large que la ligne est réduite à la largeur disponible
\BeforeBeginEnvironment{tikzpicture}{\begin{adjustbox}{max width=\linewidth}}
\AfterEndEnvironment{tikzpicture}{\end{adjustbox}}
% légendes pgfplots lisibles par-dessus les courbes
\pgfplotsset{every axis/.append style={legend style={fill=white,fill opacity=0.92,text opacity=1,draw=black!15,inner sep=3pt}}}
% étiquettes d'arêtes (syntaxe edge["..."]) avec fond blanc
\tikzset{every edge quotes/.append style={fill=white,inner sep=1.2pt}}
% bulles de cartes mentales : texte à la ligne dans la bulle, bulle un peu plus grande
\tikzset{every concept/.append style={text width=2.0cm,align=center,minimum size=2.3cm,inner sep=2pt}}
%% FIGFIX-END
\begin{document}

\pagegarde{Les composés aromatiques}{Chimie}{1ère C}{Module 1 : Chimie organique}{04}{5 h}{Cette leçon présente la structure singulière du benzène, molécule plane à électrons délocalisés, et ses principaux dérivés aromatiques. Elle étudie ensuite les réactions caractéristiques du noyau benzénique (addition et surtout substitution : halogénation, nitration, sulfonation, Friedel et Crafts) ainsi que les mesures de sécurité liées à leur manipulation.
\vspace{4pt}
\begin{multicols}{2}\small
\begin{itemize}
\item À la fin de cette leçon, je sais décrire la géométrie de la molécule de benzène (planéité, angles, longueurs de liaisons)
\item À la fin de cette leçon, je sais expliquer la délocalisation des électrons π et représenter symboliquement le noyau benzénique
\item À la fin de cette leçon, je sais interpréter les propriétés spécifiques du benzène (stabilité, préférence pour la substitution) par sa structure
\item À la fin de cette leçon, je sais citer et nommer quelques composés aromatiques usuels (toluène, phénol, naphtaline, TNT...)
\item À la fin de cette leçon, je sais nommer les trois dérivés disubstitués du benzène (ortho, méta, para)
\item À la fin de cette leçon, je sais écrire les équations-bilans des réactions d'addition (H2, Cl2) sur le benzène
\end{itemize}\end{multicols}}

\begin{sommaire}
\begin{multicols}{2}
\begin{enumerate}
\item Structure du benzène
\item Exemples de composés aromatiques et nomenclature
\item Combustion du benzène
\item Réactions d'addition sur le benzène
\item Réactions de substitution : la réactivité privilégiée du benzène
\item Sécurité et utilisations des composés aromatiques
\item Travaux pratiques
\item Méthodes et exercices résolus
\item Exercices
\item Corrigés des exercices
\item Fiche bilan
\end{enumerate}
\end{multicols}
\end{sommaire}

\begin{objectifs}
\begin{itemize}
\item À la fin de cette leçon, je sais décrire la géométrie de la molécule de benzène (planéité, angles, longueurs de liaisons)
\item À la fin de cette leçon, je sais expliquer la délocalisation des électrons π et représenter symboliquement le noyau benzénique
\item À la fin de cette leçon, je sais interpréter les propriétés spécifiques du benzène (stabilité, préférence pour la substitution) par sa structure
\item À la fin de cette leçon, je sais citer et nommer quelques composés aromatiques usuels (toluène, phénol, naphtaline, TNT...)
\item À la fin de cette leçon, je sais nommer les trois dérivés disubstitués du benzène (ortho, méta, para)
\item À la fin de cette leçon, je sais écrire les équations-bilans des réactions d'addition (H2, Cl2) sur le benzène
\item À la fin de cette leçon, je sais écrire les équations-bilans des réactions de substitution : halogénation, nitration, sulfonation, Friedel et Crafts
\item À la fin de cette leçon, je sais écrire l'équation-bilan de la combustion complète du benzène
\item À la fin de cette leçon, je sais citer les mesures de sécurité pour manipuler le benzène et les dérivés nitrés
\end{itemize}
\end{objectifs}

\begin{prerequis}
\begin{itemize}
\item Formules brutes, semi-développées et représentation de Lewis des molécules organiques
\item Liaisons covalentes simples et doubles, hybridation du carbone (notion vue en hydrocarbures)
\item Réactivité des alcènes : addition de H2, Cl2, Br2 et test à l'eau de brome
\item Écriture et équilibrage d'équations-bilans chimiques
\item Nomenclature des hydrocarbures (alcanes, alcènes, alcynes)
\end{itemize}
\end{prerequis}

\newpage

\coursec{Structure du benzène}

\courssub{Situation d'accroche : le mystère du liquide à l'odeur particulière}

À Douala, une usine fabrique des explosifs destinés aux carrières de pierre : le TNT. Ce produit célèbre est obtenu à partir d'un liquide incolore à l'odeur particulière, le \cle{benzène}. Au marché Mokolo de Yaoundé, on vend aussi des insecticides et des boules anti-mites (la naphtaline) dont l'odeur caractéristique, que tu connais sûrement, provient de molécules de la même famille. Les parfums, les colorants, de nombreux médicaments (comme le paracétamol) contiennent également ce même motif de base : le \cle{noyau benzénique}.

Mais voici le mystère. Sur le papier, le benzène de formule brute \ce{C6H6} semble extrêmement insaturé : il lui « manque » huit atomes d'hydrogène par rapport à l'alcane à six carbones, l'hexane \ce{C6H14}. On s'attendrait donc à ce qu'il se comporte comme un alcène, voire comme un composé encore plus réactif. Or l'expérience montre l'inverse : \textbf{le benzène ne décolore pas l'eau de brome à froid}, alors que le cyclohexène le fait instantanément !

\textbf{Problématique :} quelle est la véritable structure de la molécule de benzène, et comment cette structure explique-t-elle sa stabilité surprenante et sa réactivité particulière ?

\courssub{Historique : la formule de Kekulé (1865)}

Dès 1825, le physicien anglais Michael Faraday isole le benzène à partir du gaz d'éclairage. Les analyses montrent que sa formule brute est \ce{C6H6}. Pendant quarante ans, les chimistes s'arrachent les cheveux : comment six atomes de carbone et six atomes d'hydrogène peuvent-ils s'organiser de manière stable ?

En 1865, le chimiste allemand \cle{August Kekulé} propose une structure révolutionnaire : les six atomes de carbone forment un \textbf{cycle} (un hexagone), chaque carbone porte un atome d'hydrogène, et le cycle contient \textbf{trois doubles liaisons alternées} avec trois liaisons simples. Mais un problème subsiste : où placer les doubles liaisons ? Il existe en effet deux façons de les disposer, parfaitement symétriques. Kekulé suggère alors que la molécule « oscille » rapidement entre ces deux arrangements.

\begin{savaistu}[Le rêve du serpent]
Kekulé raconta lui-même qu'il aurait eu l'idée de la structure cyclique du benzène lors d'une rêverie devant sa cheminée : il vit des atomes danser, puis se transformer en un \textbf{serpent qui se mordait la queue}. Cette image lui aurait inspiré le cycle à six carbones. Qu'il s'agisse d'un vrai rêve ou d'une jolie anecdote, cette intuition a révolutionné la chimie organique : c'était la première fois qu'on imaginait une molécule en anneau. Aujourd'hui, la production mondiale de benzène dépasse 50 millions de tonnes par an !
\end{savaistu}

\courssub{Géométrie de la molécule : un hexagone plan et régulier}

Les techniques modernes (diffraction des rayons X, spectroscopie) ont permis de mesurer précisément la géométrie du benzène. Les résultats sont remarquables :

\begin{itemize}
\item la molécule est \textbf{plane} : les 6 atomes de carbone et les 6 atomes d'hydrogène sont tous dans le même plan ;
\item les 6 atomes de carbone forment un \textbf{hexagone régulier} : les six angles $\widehat{C-C-C}$ valent exactement \ang{120} ;
\item chaque atome de carbone est \textbf{hybridé $sp^2$} : il forme trois liaisons $\sigma$ (deux avec ses voisins carbone, une avec un hydrogène) disposées dans un plan à \ang{120}, exactement comme dans l'éthylène ;
\item il reste sur chaque carbone une \textbf{orbitale $p$ non hybridée}, perpendiculaire au plan du cycle, contenant un électron.
\end{itemize}

Mais la mesure la plus surprenante concerne les longueurs des liaisons carbone--carbone : \textbf{les six liaisons sont rigoureusement identiques}, avec une longueur de \SI{139}{pm}. Or une liaison simple \ce{C-C} mesure \SI{154}{pm} (dans l'éthane) et une liaison double \ce{C=C} mesure \SI{134}{pm} (dans l'éthylène). La liaison du benzène est donc \textbf{intermédiaire} entre une simple et une double liaison !

\begin{center}
\begin{tabularx}{\linewidth}{|l|X|X|X|}
\hline
\rowcolor{popblueL}
\textbf{Liaison} & \textbf{Molécule de référence} & \textbf{Longueur (pm)} & \textbf{Énergie (kJ.mol$^{-1}$)} \\
\hline
Simple \ce{C-C} & Éthane \ce{C2H6} & 154 & 348 \\
\hline
\ce{C-C} du benzène & Benzène \ce{C6H6} & \textbf{139} & intermédiaire \\
\hline
Double \ce{C=C} & Éthylène \ce{C2H4} & 134 & 612 \\
\hline
\end{tabularx}
\end{center}

Ce résultat est incompatible avec la formule de Kekulé, qui prévoit trois liaisons longues (\SI{154}{pm}) et trois liaisons courtes (\SI{134}{pm}) alternées. Il faut donc aller plus loin.

\courssub{Les électrons $\pi$ délocalisés : la clé du mystère}

Chacun des six carbones $sp^2$ possède une orbitale $p$ perpendiculaire au plan du cycle, occupée par un électron. Ces six orbitales $p$, toutes parallèles entre elles et très proches, se recouvrent latéralement \textbf{non pas deux à deux} (comme dans un alcène), mais \textbf{toutes ensemble, sur tout le pourtour du cycle}. Elles fusionnent en un \cle{nuage électronique délocalisé} ayant la forme d'un anneau (comme un « beignet ») au-dessus et au-dessous du plan de la molécule.

Les \textbf{six électrons $\pi$} n'appartiennent donc plus à une liaison particulière : ils appartiennent à l'ensemble du cycle. On dit qu'ils sont \cle{délocalisés}.

\begin{popfigure}
\begin{center}
\begin{tikzpicture}[scale=0.9]
% Cycle : sommets aux angles 90, 150, 210, 270, 330, 30 (pointe en haut, comme chemfig)
\foreach \a in {90, 150, 210, 270, 330, 30} {
  \draw[thick,popdark,fill=popcream] (\a:1.3) -- (\a+60:1.3);
  \fill[popdark] (\a:1.3) circle (0.09);
}
% Orbitales p (ellipses verticales sur chaque carbone)
\foreach \a in {90, 150, 210, 270, 330, 30} {
  \draw[popblue,thick,fill=popblue!25,opacity=0.75] (\a:1.3) ++(0,0.55) ellipse (0.16 and 0.5);
  \draw[popblue,thick,fill=popblue!25,opacity=0.75] (\a:1.3) ++(0,-0.55) ellipse (0.16 and 0.5);
}
% Nuage délocalisé (anneaux)
\draw[poporange,very thick,opacity=0.9] (0,0.85) ellipse (1.75 and 0.55);
\draw[poporange,very thick,opacity=0.9] (0,-0.85) ellipse (1.75 and 0.55);
\node[poporange,align=center] at (3.6,1.3) {\small nuage $\pi$\\ \small délocalisé};
\draw[poporange,-{Stealth}] (2.6,1.15) -- (1.6,0.95);
\node[popblue,align=center] at (-3.7,1.3) {\small orbitales $p$\\ \small perpendiculaires\\ \small au plan};
\draw[popblue,-{Stealth}] (-2.6,1.15) -- (-1.55,1.0);
\end{tikzpicture}
\end{center}
\legende{Les six orbitales $p$ (bleu) sont perpendiculaires au plan du cycle ; leur recouvrement latéral continu crée le nuage $\pi$ délocalisé (orange), au-dessus et au-dessous du plan.}
\end{popfigure}

\begin{definition}[Noyau benzénique]
Le \cle{noyau benzénique} est un cycle plan de 6 atomes de carbone hybridés $sp^2$, dont les 6 électrons $\pi$ sont \textbf{délocalisés} sur l'ensemble du cycle. Toutes les liaisons carbone--carbone y sont identiques, de longueur intermédiaire (\SI{139}{pm}) entre une liaison simple et une liaison double.
\end{definition}

\begin{definition}[Composé aromatique]
Un \cle{composé aromatique} est un composé organique comportant au moins un \textbf{noyau benzénique} à électrons $\pi$ délocalisés. Le nom vient de l'« arôme » (l'odeur) souvent agréable des premiers composés de cette famille découverts (vanille, amandes amères, girofle).
\end{definition}

\courssub{Représentation symbolique du noyau benzénique}

Puisque les électrons $\pi$ ne sont pas localisés entre deux carbones précis, on ne peut pas représenter le benzène par une seule formule avec des doubles liaisons fixes. Deux conventions coexistent :

\begin{itemize}
\item \textbf{les deux formules de Kekulé} : deux hexagones avec trois doubles liaisons alternées, séparés par une double flèche de \textbf{mésomérie} $\leftrightarrow$. La molécule réelle n'est ni l'une ni l'autre : c'est un \emph{hybride de résonance}, « moyenne » des deux ;
\item \textbf{la représentation au cercle} : un hexagone avec un \textbf{cercle intérieur}, qui symbolise élégamment le nuage des 6 électrons $\pi$ délocalisés sur tout le cycle. C'est la représentation moderne, la plus fidèle à la réalité.
\end{itemize}

\begin{popfigure}
\begin{center}
\chemfig{*6((-H)=(-H)-(-H)=(-H)-(-H)=(-H)-)}
\hspace{0,8cm}
$\xrightarrow{\text{mésomérie}}$
\hspace{0,8cm}
\chemfig{*6((-H)-(-H)=(-H)-(-H)=(-H)-(-H)=)}
\hspace{0,8cm}
$\equiv$
\hspace{0,8cm}
\begin{tikzpicture}[baseline=(current bounding box.center)]
\draw[thick,popdark] (90:1.1) \foreach \a in {150,210,270,330,30} { -- (\a:1.1) } -- cycle;
\draw[thick,poporange] (0,0) circle (0.62);
\foreach \a in {90,150,210,270,330,30} {
  \node[popdark,font=\small] at (\a:1.55) {H};
  \draw[popdark,thick] (\a:1.1) -- (\a:1.38);
}
\end{tikzpicture}
\end{center}
\legende{Les deux formules limites de Kekulé (à gauche et au centre), liées par la mésomérie, et la représentation symbolique moderne : hexagone avec cercle intérieur (à droite).}
\end{popfigure}

\begin{attention}[Erreur fréquente au Probatoire]
Ne dessine \textbf{jamais} le benzène avec trois doubles liaisons localisées fixes sans préciser la délocalisation ! Si tu écris une formule de Kekulé, tu dois soit dessiner les \textbf{deux} formules séparées par la double flèche $\leftrightarrow$, soit utiliser l'hexagone au cercle. Un hexagone seul avec trois doubles liaisons laisserait croire qu'il existe des liaisons simples et doubles distinctes, ce qui est \textbf{faux} : les six liaisons sont identiques. De même, la double flèche $\leftrightarrow$ ne signifie pas que la molécule « bascule » d'une forme à l'autre : la molécule réelle est unique, c'est un hybride des deux représentations.
\end{attention}

\courssub{Conséquence expérimentale : le benzène n'est pas un alcène}

La délocalisation des électrons $\pi$ a une conséquence spectaculaire et mesurable : le benzène \textbf{ne réagit pas} comme un composé à doubles liaisons ordinaire.

\begin{experience}[Action de l'eau de brome sur le benzène et sur le cyclohexène]
\textbf{Protocole :} dans deux tubes à essais, on verse \SI{2}{mL} d'eau de brome (solution orange de \ce{Br2}). On ajoute dans le premier tube quelques gouttes de cyclohexène \ce{C6H10}, dans le second quelques gouttes de benzène \ce{C6H6}. On agite et on observe à froid, à l'abri de la lumière vive.

\textbf{Observations :}
\begin{itemize}
\item tube 1 (cyclohexène) : la solution se \textbf{décolore immédiatement} ;
\item tube 2 (benzène) : la solution reste \textbf{orange} : aucune réaction à froid.
\end{itemize}

\textbf{Interprétation :} le cyclohexène possède une vraie double liaison \ce{C=C} localisée, riche en électrons $\pi$ accessibles : elle réalise une \textbf{addition} du dibrome :
\[ \ce{C6H10 + Br2 -> C6H10Br2} \]
Dans le benzène, les électrons $\pi$ sont délocalisés et « piégés » dans le nuage aromatique stable : ils ne sont pas disponibles pour additionner \ce{Br2}. Le benzène ne décolore donc pas l'eau de brome à froid.
\end{experience}

\begin{methode}[Justifier la stabilité particulière du benzène]
Pour démontrer que le benzène n'est pas un alcène « ordinaire », suis cette démarche :
\begin{enumerate}
\item Rappelle la formule brute \ce{C6H6} : forte insaturation apparente (8 H de moins que l'hexane).
\item Cite le test à l'eau de brome : un alcène (ex. cyclohexène) décolore \ce{Br2} à froid par addition ; le benzène \textbf{ne le décolore pas}.
\item Conclus sur la structure : les électrons $\pi$ ne sont pas localisés dans des doubles liaisons, mais \textbf{délocalisés} sur tout le cycle.
\item Appuie-toi sur les mesures : 6 liaisons \ce{C-C} identiques de \SI{139}{pm}, intermédiaires entre simple et double.
\item Termine par l'énergie : la délocalisation stabilise la molécule (voir ci-dessous).
\end{enumerate}
\end{methode}

\courssub{Stabilité particulière du noyau aromatique : l'énergie de délocalisation}

La délocalisation des électrons $\pi$ \textbf{abaisse l'énergie} de la molécule : le benzène réel est beaucoup plus stable que ne le serait un hypothétique « cyclohexatriène » à trois doubles liaisons localisées. On peut le mesurer grâce à l'hydrogénation (addition de dihydrogène).

L'hydrogénation du cyclohexène (une seule double liaison) dégage \SI{120}{kJ.mol^{-1}} :
\[ \ce{C6H10 + H2 ->[Ni] C6H12} \qquad \Delta H = \SI{-120}{kJ.mol^{-1}} \]

Si le benzène possédait trois doubles liaisons \emph{indépendantes}, son hydrogénation complète en cyclohexane devrait dégager $3 \times 120 = \SI{360}{kJ.mol^{-1}}$. Or l'expérience donne seulement \SI{208}{kJ.mol^{-1}} :
\[ \ce{C6H6 + 3 H2 ->[Ni,\ \text{chauffage}] C6H12} \qquad \Delta H = \SI{-208}{kJ.mol^{-1}} \]

Le benzène « réel » dégage donc \SI{152}{kJ.mol^{-1}} de \textbf{moins} que prévu : il était déjà plus stable de cette quantité d'énergie. Cette différence est appelée \cle{énergie de stabilisation par délocalisation} (ou énergie de résonance).

\begin{exoresolu}[Calcul de l'énergie de délocalisation]
\textbf{Énoncé :} L'hydrogénation d'une mole de cyclohexène dégage \SI{120}{kJ.mol^{-1}}. L'hydrogénation complète d'une mole de benzène en cyclohexane dégage \SI{208}{kJ.mol^{-1}}. Calculer l'énergie de stabilisation par délocalisation du benzène et commenter.
\tcblower
\textbf{Solution détaillée :}

\textbf{Étape 1 :} si le benzène contenait trois doubles liaisons localisées et indépendantes, l'énergie dégagée par son hydrogénation serait :
\[ \Delta H_{\text{prévu}} = 3 \times (-120) = \SI{-360}{kJ.mol^{-1}} \]

\textbf{Étape 2 :} la valeur expérimentale est :
\[ \Delta H_{\text{exp}} = \SI{-208}{kJ.mol^{-1}} \]

\textbf{Étape 3 :} l'écart entre les deux valeurs est l'énergie de délocalisation :
\[ E_{\text{déloc}} = 360 - 208 = \SI{152}{kJ.mol^{-1}} \]

\textbf{Conclusion :} le benzène est plus stable de \SI{152}{kJ.mol^{-1}} qu'un cycle à trois doubles liaisons localisées. Cette énergie considérable explique pourquoi le benzène préfère conserver son noyau aromatique (réactions de \textbf{substitution}) plutôt que de le détruire (réactions d'\textbf{addition}), contrairement aux alcènes.
\end{exoresolu}

\courssub{Complément Probatoire : comparaison benzène / cyclohexène}

Le tableau suivant résume la comparaison classique demandée au Probatoire entre le benzène et le cyclohexène, deux hydrocarbures cycliques à six carbones.

\begin{center}
\begin{tabularx}{\linewidth}{|X|X|X|}
\hline
\rowcolor{popblueL}
\textbf{Critère} & \textbf{Cyclohexène \ce{C6H10}} & \textbf{Benzène \ce{C6H6}} \\
\hline
Nature des liaisons & 1 double liaison \ce{C=C} localisée (\SI{134}{pm}) & 6 liaisons identiques (\SI{139}{pm}), électrons $\pi$ délocalisés \\
\hline
Eau de brome à froid & Décoloration immédiate (addition) & Aucune décoloration \\
\hline
Hydrogénation \ce{H2}/Ni & Facile, $\Delta H = \SI{-120}{kJ.mol^{-1}}$ & Difficile (catalyseur + chauffage), $\Delta H = \SI{-208}{kJ.mol^{-1}}$ \\
\hline
Réactivité privilégiée & \textbf{Addition} & \textbf{Substitution} (le noyau est conservé) \\
\hline
\end{tabularx}
\end{center}

\begin{aretenir}[L'essentiel de la structure du benzène]
\begin{itemize}
\item Formule brute : \ce{C6H6} ; molécule \textbf{plane}, hexagone régulier, angles de \ang{120}, carbones hybridés $sp^2$.
\item Les \textbf{6 électrons $\pi$ sont délocalisés} sur tout le cycle : nuage $\pi$ au-dessus et au-dessous du plan.
\item Les 6 liaisons \ce{C-C} sont \textbf{identiques} (\SI{139}{pm}), intermédiaires entre simple et double.
\item Représentation symbolique : \textbf{hexagone avec un cercle intérieur} (ou les deux formules de Kekulé liées par $\leftrightarrow$).
\item Conséquence : le benzène ne décolore \textbf{pas} l'eau de brome à froid et préfère la \textbf{substitution} à l'addition.
\item Stabilité exceptionnelle : énergie de délocalisation d'environ \SI{152}{kJ.mol^{-1}}.
\end{itemize}
\end{aretenir}

\begin{exercice}[Pour t'entraîner]
On réalise deux expériences avec un hydrocarbure inconnu $X$ de formule \ce{C6H6}. Expérience 1 : $X$ ne décolore pas l'eau de brome à froid. Expérience 2 : la diffraction des rayons X montre que toutes les liaisons carbone--carbone de $X$ ont la même longueur, \SI{139}{pm}.
\begin{enumerate}
\item Montrer que $X$ ne peut pas être le cyclohexa-1,3,5-triène à doubles liaisons localisées.
\item Proposer une représentation de $X$ cohérente avec ces deux expériences.
\item Prévoir le signe et l'ordre de grandeur de l'écart entre l'enthalpie d'hydrogénation expérimentale de $X$ et celle calculée pour trois doubles liaisons indépendantes.
\end{enumerate}
\end{exercice}

\begin{corrige}[Exercice n°1]
\textbf{1.} Un cyclohexatriène à doubles liaisons localisées posséderait de vraies liaisons \ce{C=C} : il décolorerait l'eau de brome à froid par addition, et présenterait deux longueurs de liaison distinctes (\SI{134}{pm} et \SI{154}{pm}). Or $X$ ne décolore pas \ce{Br2} et ses six liaisons sont identiques : les doubles liaisons ne sont donc pas localisées.

\textbf{2.} $X$ est le benzène : on le représente par un hexagone régulier avec un cercle intérieur symbolisant les 6 électrons $\pi$ délocalisés (ou par les deux formules de Kekulé séparées par $\leftrightarrow$).

\textbf{3.} L'enthalpie d'hydrogénation expérimentale est \textbf{moins exothermique} que la valeur calculée ($\SI{-208}{kJ.mol^{-1}}$ au lieu de $\SI{-360}{kJ.mol^{-1}}$) : l'écart, environ \SI{152}{kJ.mol^{-1}}, correspond à l'énergie de stabilisation par délocalisation du noyau aromatique.
\end{corrige}

% Fin de la section 1 : Structure du benzène

\coursec{Exemples de composés aromatiques et nomenclature}

Dans la section précédente, nous avons découvert la structure remarquable du benzène : un cycle plan de six atomes de carbone portant six électrons délocalisés, représenté symboliquement par un hexagone cerclé. Mais le benzène n'est qu'un point de départ. En remplaçant un ou plusieurs de ses atomes d'hydrogène par d'autres groupes d'atomes, on obtient une immense famille de composés : les \cle{dérivés du benzène}. Ce sont eux que l'on retrouve dans les solvants de peinture, les colorants, les médicaments comme l'aspirine, les insecticides et même les explosifs comme le TNT. Encore faut-il savoir les reconnaître et les nommer sans erreur : c'est précisément l'objectif de cette section, et c'est une compétence très souvent évaluée au Probatoire.

\courssub{Les dérivés monosubstitués du benzène}

\begin{definition}[Dérivé monosubstitué]
Un \cle{dérivé monosubstitué} du benzène est un composé aromatique obtenu en remplaçant \textbf{un seul} atome d'hydrogène du benzène \ce{C6H6} par un atome ou un groupe d'atomes appelé \cle{substituant}. Sa formule brute générale s'écrit $\ce{C6H5}$-X, où $\ce{C6H5}$- est le groupe \cle{phényle}.
\end{definition}

Puisque les six atomes d'hydrogène du benzène sont rigoureusement équivalents (conséquence de la délocalisation électronique), il n'existe qu'\textbf{un seul} dérivé monosubstitué pour chaque substituant : peu importe l'hydrogène remplacé, la molécule obtenue est la même. Voici les quatre exemples classiques à connaître par cœur :

\begin{center}
\begin{tabularx}{\linewidth}{|l|l|X|}
\hline
\rowcolor{popblueL}
\textbf{Substituant} & \textbf{Nom du composé} & \textbf{Utilisation principale} \\
\hline
\ce{-CH3} (méthyle) & \textbf{Toluène} (méthylbenzène) & Solvant des peintures et colles ; matière première du TNT \\
\hline
\ce{-OH} (hydroxyle) & \textbf{Phénol} (hydroxybenzène) & Désinfectant, matière première des résines et de l'aspirine \\
\hline
\ce{-Cl} (chlore) & \textbf{Chlorobenzène} & Solvant, intermédiaire de synthèse (autrefois pour le DDT) \\
\hline
\ce{-NO2} (nitro) & \textbf{Nitrobenzène} & Intermédiaire de synthèse des colorants et de l'aniline \\
\hline
\end{tabularx}
\end{center}

\begin{attention}[Le phénol n'est pas un alcool]
Dans le phénol \ce{C6H5-OH}, le groupe \ce{-OH} est porté directement par un carbone du noyau benzénique. On ne dit jamais « alcool benzylique » pour cette molécule (l'alcool benzylique est \ce{C6H5-CH2-OH}) : le phénol possède des propriétés propres, notamment un caractère acide marqué, qui le distinguent radicalement des alcools. Retiens : \textbf{un \ce{-OH} sur un noyau benzénique = un phénol}, pas un alcool.
\end{attention}

\courssub{Les dérivés disubstitués : ortho, méta, para}

Dès que \textbf{deux} substituants sont fixés sur le noyau, une nouveauté apparaît : les deux groupes peuvent occuper des positions relatives différentes. Il existe exactement \textbf{trois} possibilités, et trois seulement.

\begin{definition}[Positions ortho, méta et para]
Sur un noyau benzénique portant deux substituants, on distingue trois positions relatives :
\begin{itemize}
\item \cle{ortho} (noté \emph{o-}) : les deux substituants sont sur deux carbones \textbf{adjacents} ; positions \textbf{1,2} ;
\item \cle{méta} (noté \emph{m-}) : les deux substituants sont séparés par \textbf{un} carbone ; positions \textbf{1,3} ;
\item \cle{para} (noté \emph{p-}) : les deux substituants sont \textbf{opposés} sur le cycle ; positions \textbf{1,4}.
\end{itemize}
\end{definition}

L'exemple de référence est celui des \cle{xylènes} (diméthylbenzènes, \ce{C6H4(CH3)2}), mélange de solvants très utilisé dans les laboratoires et l'industrie des peintures :

\begin{itemize}
\item \textbf{ortho-xylène} : 1,2-diméthylbenzène ;
\item \textbf{méta-xylène} : 1,3-diméthylbenzène ;
\item \textbf{para-xylène} : 1,4-diméthylbenzène.
\end{itemize}

\begin{popfigure}
\begin{center}
\setchemfig{atom sep=1.7em}
\chemfig{*6((-CH_3)=(-CH_3)-=-=)}
\hspace{1.2cm}
\chemfig{*6((-CH_3)=-(-CH_3)=-=)}
\hspace{1.2cm}
\chemfig{*6((-CH_3)=-=(-CH_3)-=)}
\end{center}
\vspace{0.3cm}
\begin{center}
\begin{tabularx}{\linewidth}{XXX}
\centering \textbf{ortho-xylène} (1,2) & \centering \textbf{méta-xylène} (1,3) & \centering \arraybackslash \textbf{para-xylène} (1,4)
\end{tabularx}
\end{center}
\legende{Les trois isomères du xylène \ce{C6H4(CH3)2}. Représentation de Kekulé (liaisons doubles localisées) pour montrer la connectivité ; les deux groupes méthyle sont adjacents (ortho), séparés d'un carbone (méta) ou opposés (para).}
\end{popfigure}

\begin{methode}[Nommer un dérivé disubstitué du benzène]
Pour nommer correctement un dérivé disubstitué, suis toujours la même démarche :
\begin{enumerate}
\item \textbf{Numéroter} les six carbones du cycle de façon à attribuer aux deux substituants les \textbf{indices les plus petits possibles} (1,2 plutôt que 1,6 ; 1,3 plutôt que 1,5 ; 1,4 plutôt que 1,5 ou 1,6).
\item En cas d'égalité entre deux numérotations, donner l'indice le plus petit au substituant cité \textbf{en premier dans l'ordre alphabétique}.
\item \textbf{Citer les substituants par ordre alphabétique} (chloro avant méthyle, méthyle avant nitro), chacun précédé de son indice de position.
\item Terminer par le nom de la chaîne principale : « benzène », ou un nom d'usage (toluène, phénol, aniline, acide benzoïque) si le groupe prioritaire le justifie.
\end{enumerate}
Exemple : le dérivé \ce{Cl-C6H4-NO2} avec les deux groupes opposés se nomme \textbf{1-chloro-4-nitrobenzène} (ou \emph{para}-chloronitrobenzène), et non « 1-nitro-4-chlorobenzène » : chloro passe avant nitro dans l'alphabet.
\end{methode}

\begin{attention}[Piège classique du Probatoire : méta contre para]
L'erreur la plus fréquente dans les copies est de confondre \textbf{méta (1,3)} et \textbf{para (1,4)}. Astuce visuelle infaillible : dessine le cycle comme un hexagone pointe en haut et numérote les sommets dans le sens des aiguilles d'une montre à partir du premier substituant. En \emph{para}, le deuxième substituant est \textbf{à l'opposé exact} du premier (les deux groupes sont sur un même axe de symétrie, « face à face »). En \emph{méta}, il y a \textbf{un seul carbone} entre les deux substituants. Vérifie aussi la cohérence : 1,2 = ortho, 1,3 = méta, 1,4 = para. Toute autre numérotation (1,5 ; 1,6) signifie que tu n'as pas choisi les indices les plus petits.
\end{attention}

\begin{exoresolu}[Nommer des dérivés disubstitués]
Énoncé : on considère les trois dérivés du benzène suivants : (a) deux groupes \ce{-CH3} en positions 1 et 4 ; (b) un groupe \ce{-OH} et un groupe \ce{-NO2} sur deux carbones adjacents ; (c) deux atomes de chlore séparés par un carbone du cycle. Donner le nom systématique et le nom usuel (avec ortho, méta ou para) de chaque composé.
\tcblower
Solution détaillée :
\begin{enumerate}
\item[(a)] Deux méthyles en positions 1 et 4 : les groupes sont opposés. Nom systématique : \textbf{1,4-diméthylbenzène} ; nom usuel : \textbf{para-xylène} (\emph{p}-xylène).
\item[(b)] \ce{-OH} et \ce{-NO2} sur des carbones adjacents : positions 1,2. Le groupe hydroxyle étant prioritaire (nom d'usage « phénol »), on numérote le carbone portant \ce{-OH} en position 1. Nom : \textbf{2-nitrophénol}, ou \textbf{ortho-nitrophénol}.
\item[(c)] Deux chlores séparés par un carbone : positions 1,3. Nom : \textbf{1,3-dichlorobenzène}, ou \textbf{méta-dichlorobenzène} (\emph{m}-dichlorobenzène).
\end{enumerate}
Vérification de la règle des indices : pour (a), 1,4 est la plus petite numérotation possible (1,4 et non 1,5 ou 2,5) ; pour (c), 1,3 et non 1,5.
\end{exoresolu}

\courssub{Les dérivés polysubstitués : TNT et acide picrique}

Lorsque trois substituants ou plus sont fixés sur le noyau, les préfixes ortho, méta et para ne suffisent plus : on utilise obligatoirement la \textbf{numérotation}. Deux molécules célèbres illustrent ce cas.

\begin{exemplebox}[Le TNT : triple nitration du toluène]
Le \cle{TNT} est le \textbf{2,4,6-trinitrotoluène} : un noyau benzénique portant un groupe méthyle \ce{-CH3} (position 1) et trois groupes nitro \ce{-NO2} (positions 2, 4 et 6). Il s'obtient industriellement par \textbf{nitration successive du toluène} en trois étapes, à l'aide du mélange sulfonitrique (acide nitrique concentré + acide sulfurique concentré) :
\[ \ce{C6H5-CH3 + 3 HNO3 ->[H2SO4\ conc.] C6H2(NO2)3-CH3 + 3 H2O} \]
Chaque étape est une \textbf{substitution électrophile} : un hydrogène du noyau est remplacé par un groupe \ce{-NO2}. L'acide sulfurique joue le rôle de \textbf{catalyseur} (il génère l'ion nitronium \ce{NO2+}, l'espèce réactive). La réaction est fortement \textbf{exothermique} : la température doit être contrôlée avec soin à chaque étape, la troisième nitration exigeant des conditions plus énergiques que les deux premières.
\end{exemplebox}

L'\cle{acide picrique} (2,4,6-trinitrophénol) est le cousin du TNT : le groupe \ce{-CH3} est remplacé par un groupe \ce{-OH}. C'est un acide fort... et un explosif encore plus puissant que le TNT, longtemps utilisé dans les obus avant d'être abandonné car il forme avec les métaux des sels (picrates) dangereusement sensibles aux chocs.

\courssub{Les composés à plusieurs noyaux benzéniques}

Certains composés aromatiques possèdent plusieurs noyaux benzéniques \textbf{condensés}, c'est-à-dire partageant un côté commun :

\begin{itemize}
\item le \cle{naphtaline} \ce{C10H8} : deux noyaux condensés. C'est le principe actif des \textbf{boules anti-mites} que l'on glisse dans les armoires ; il se sublime lentement à température ambiante, d'où son odeur caractéristique ;
\item l'\cle{anthracène} \ce{C14H10} : trois noyaux condensés en ligne, présent dans le goudron de houille et utilisé dans la fabrication de colorants.
\end{itemize}

\begin{popfigure}
\begin{center}
\setchemfig{atom sep=1.6em}
\chemfig{*6((-CH_3)=-=-=-)}
\hspace{0.9cm}
\chemfig{*6((-OH)=-=-=-)}
\hspace{0.9cm}
\chemfig{*6((-NO_2)=-=-=-)}
\hspace{0.9cm}
\chemfig{*6(=-*6(-=-=-)=-=-)}
\end{center}
\vspace{0.3cm}
\begin{center}
\begin{tabularx}{\linewidth}{XXXX}
\centering \textbf{toluène} & \centering \textbf{phénol} & \centering \textbf{nitrobenzène} & \centering \arraybackslash \textbf{naphtaline}
\end{tabularx}
\end{center}
\vspace{0.5cm}
\begin{center}
\setchemfig{atom sep=1.6em}
\chemfig{*6((-CH_3)=(-NO_2)-=(-NO_2)-=(-NO_2)-)}
\end{center}
\vspace{0.2cm}
\begin{center}\textbf{TNT} : 2,4,6-trinitrotoluène\end{center}
\legende{Galerie de composés aromatiques usuels. Représentation de Kekulé pour la connectivité ; le cercle symbolique (non dessiné ici) représente les six électrons délocalisés de chaque noyau benzénique.}
\end{popfigure}

\begin{savaistu}[Des boules anti-mites au TNT : une même famille]
Le naphtaline qui protège tes vêtements des mites et le TNT des explosifs appartiennent à la \textbf{même grande famille} : les composés aromatiques dérivés du benzène. Leur comportement commun s'explique par la présence du noyau benzénique et de ses électrons délocalisés. Autre exemple frappant : l'\textbf{aspirine} (acide acétylsalicylique), le médicament le plus consommé au monde, possède un noyau benzénique portant un groupe \ce{-COOH} et un groupe ester en position ortho. Le \textbf{DDT}, insecticide qui a sauvé des millions de vies contre le paludisme avant d'être interdit pour sa persistance dans l'environnement, contient deux noyaux chlorés. Solvants, colorants, médicaments, insecticides, explosifs : une seule structure de base, le noyau \ce{C6}, et une infinité d'applications.
\end{savaistu}

\courssub{Activité de manipulation : représenter les molécules aromatiques}

\begin{experience}[Construction de modèles moléculaires]
\textbf{Matériel} : boîte de modèles moléculaires (boules noires pour le carbone, blanches pour l'hydrogène, rouges pour l'oxygène, bleues pour l'azote, vertes pour le chlore ; bâtonnets pour les liaisons).

\textbf{Protocole} :
\begin{enumerate}
\item Construire le benzène \ce{C6H6} : vérifier que le cycle est \textbf{plan} et que les angles sont voisins de $120^\circ$.
\item Remplacer un hydrogène par un groupe \ce{-CH3} : obtenir le toluène. Vérifier qu'il n'existe qu'\textbf{un seul} dérivé monosubstitué possible.
\item Ajouter un deuxième groupe \ce{-CH3} : construire successivement les \textbf{trois} xylènes (ortho, méta, para) et constater qu'il n'existe aucune quatrième possibilité.
\item Construire le TNT en fixant trois groupes \ce{-NO2} en positions 2, 4 et 6 du toluène.
\end{enumerate}

\textbf{Observations} : le cycle reste plan dans tous les cas ; les substituants s'écartent les uns des autres (encombrement stérique), surtout en position ortho où ils sont voisins.

\textbf{Interprétation} : la planéité du noyau et l'équivalence des six hydrogènes expliquent le nombre limité d'isomères : un seul dérivé monosubstitué, trois dérivés disubstitués.
\end{experience}

\begin{exercice}[Pour t'entraîner]
1. Écrire la formule semi-développée et donner le nom du dérivé du benzène portant un groupe \ce{-CH3} et un atome de chlore en position para.
2. Combien existe-t-il de trichlorobenzènes isomères ? Donner la numérotation de chacun.
3. Le paradichlorobenzène est utilisé comme anti-mite en remplacement du naphtaline. À quel isomère (1,2 ; 1,3 ou 1,4) correspond son nom ?
\end{exercice}

\begin{corrige}[Exercice]
1. \emph{para}-chlorotoluène, nom systématique \textbf{1-chloro-4-méthylbenzène} : un noyau \ce{C6H4} portant \ce{-Cl} et \ce{-CH3} en positions opposées.
2. Il existe \textbf{trois} trichlorobenzènes : 1,2,3-trichlorobenzène, 1,2,4-trichlorobenzène et 1,3,5-trichlorobenzène (toute autre numérotation se ramène à l'une de celles-ci par la règle des indices les plus petits).
3. « Para » signifie positions 1,4 : c'est le \textbf{1,4-dichlorobenzène}.
\end{corrige}

\begin{aretenir}[L'essentiel de la section]
\begin{itemize}
\item Un dérivé \textbf{monosubstitué} du benzène est unique : toluène (\ce{-CH3}), phénol (\ce{-OH}), chlorobenzène (\ce{-Cl}), nitrobenzène (\ce{-NO2}).
\item Un dérivé \textbf{disubstitué} existe sous \textbf{trois} isomères : \textbf{ortho} (1,2), \textbf{méta} (1,3), \textbf{para} (1,4).
\item Nomenclature : indices les plus petits possibles, substituants cités par \textbf{ordre alphabétique}.
\item Au-delà de deux substituants, on numérote : TNT = \textbf{2,4,6-trinitrotoluène}, obtenu par triple nitration du toluène ; acide picrique = 2,4,6-trinitrophénol.
\item Les noyaux peuvent se condenser : \textbf{naphtaline} (2 noyaux, anti-mites), \textbf{anthracène} (3 noyaux).
\item Les composés aromatiques sont partout : solvants, colorants, aspirine, DDT, TNT — d'où l'importance des règles de sécurité que nous étudierons à la fin de la leçon.
\end{itemize}
\end{aretenir}

\coursec{Combustion du benzène}

Dans la section précédente, tu as découvert que le benzène \ce{C6H6} est le chef de file d'une grande famille de composés aromatiques : toluène, xylènes, phénol, aniline\ldots{} Tous partagent le même noyau benzénique à électrons délocalisés. Mais à quoi servent concrètement ces molécules, et comment réagissent-elles ? Nous abordons maintenant leurs \cle{propriétés chimiques}, en commençant par la plus universelle des réactions des hydrocarbures : la \cle{combustion}.

Tu as peut-être déjà remarqué, lors d'une fête de quartier à Yaoundé ou à Douala, la flamme jaune et très fumeuse d'une lampe tempête mal réglée ou d'un feu de pneus : cette fumée noire est de la suie, du carbone imbrûlé. Le benzène, extrêmement riche en carbone, brûle exactement de cette manière spectaculaire. Comprendre pourquoi, c'est l'objectif de cette section.

\courssub{La combustion complète du benzène}

\textbf{Intuition d'abord.} Brûler un hydrocarbure, c'est le faire réagir avec le dioxygène de l'air. Si l'apport en dioxygène est suffisant, chaque atome de carbone de la molécule est oxydé jusqu'au dioxyde de carbone \ce{CO2}, et chaque atome d'hydrogène jusqu'à l'eau \ce{H2O}. C'est la \cle{combustion complète}.

\begin{definition}[Combustion complète d'un hydrocarbure]
La \cle{combustion complète} d'un hydrocarbure de formule brute $C_xH_y$ est sa réaction totale avec le dioxygène \ce{O2}. Elle produit exclusivement du \cle{dioxyde de carbone} \ce{CO2} et de la \cle{vapeur d'eau} \ce{H2O}, en dégageant une grande quantité d'énergie thermique (réaction fortement \emph{exothermique}). L'équation-bilan générale s'écrit :
\[ C_xH_y + \left(x + \frac{y}{4}\right) O_2 \longrightarrow x CO_2 + \frac{y}{2} H_2O \]
\end{definition}

Appliquons ce résultat au benzène, de formule brute \ce{C6H6} (donc $x = 6$ et $y = 6$). Le coefficient du dioxygène vaut $x + \dfrac{y}{4} = 6 + \dfrac{6}{4} = 6 + 1{,}5 = 7{,}5 = \dfrac{15}{2}$. On obtient d'abord :
\[ \ce{C6H6 + 15/2 O2 -> 6 CO2 + 3 H2O} \]
Un coefficient fractionnaire étant inesthétique en chimie, on multiplie toute l'équation par 2.

\begin{aretenir}[Équation-bilan de la combustion complète du benzène]
\[ \ce{2 C6H6 + 15 O2 -> 12 CO2 + 6 H2O} \]
Cette réaction est \emph{totale} (flèche simple), \emph{exothermique}, et ne nécessite ni catalyseur ni chauffage prolongé : une simple flamme ou une étincelle suffit à l'amorcer, ensuite elle s'auto-entretient.
\end{aretenir}

\begin{methode}[Équilibrer la combustion d'un hydrocarbure $C_xH_y$]
Pour équilibrer à coup sûr l'équation-bilan de la combustion complète de n'importe quel hydrocarbure, suis toujours le même ordre, dit « ordre C -- H -- O » :
\begin{enumerate}
\item \textbf{Équilibrer le carbone} : placer le coefficient $x$ devant \ce{CO2}. Tous les atomes de carbone de $C_xH_y$ se retrouvent dans le dioxyde de carbone.
\item \textbf{Équilibrer l'hydrogène} : placer le coefficient $\dfrac{y}{2}$ devant \ce{H2O} (chaque molécule d'eau contient 2 atomes d'hydrogène).
\item \textbf{Équilibrer l'oxygène en dernier} : compter les atomes d'oxygène déjà placés à droite, soit $2x + \dfrac{y}{2}$, puis écrire devant \ce{O2} le coefficient $\dfrac{2x + y/2}{2} = x + \dfrac{y}{4}$.
\item \textbf{Éliminer les fractions} : si le coefficient de \ce{O2} est fractionnaire (dénominateur 2 ou 4), multiplier \emph{tous} les coefficients de l'équation par ce dénominateur.
\item \textbf{Vérifier} : compter les atomes de chaque élément dans les deux membres.
\end{enumerate}
\end{methode}

\begin{exemplebox}[Équilibrage pas à pas de la combustion du benzène]
Appliquons la méthode à \ce{C6H6} ($x = 6$, $y = 6$).
\begin{itemize}
\item \textbf{Étape 1 (carbone)} : 6 atomes de C à gauche $\Rightarrow$ coefficient 6 devant \ce{CO2} :
\[ \ce{C6H6 + O2 -> 6 CO2 + H2O} \]
\item \textbf{Étape 2 (hydrogène)} : 6 atomes de H à gauche $\Rightarrow$ coefficient $\dfrac{6}{2} = 3$ devant \ce{H2O} :
\[ \ce{C6H6 + O2 -> 6 CO2 + 3 H2O} \]
\item \textbf{Étape 3 (oxygène)} : à droite, on compte $6 \times 2 + 3 \times 1 = 15$ atomes d'oxygène. Il faut donc $\dfrac{15}{2}$ molécules de \ce{O2} :
\[ \ce{C6H6 + 15/2 O2 -> 6 CO2 + 3 H2O} \]
\item \textbf{Étape 4 (suppression de la fraction)} : on multiplie tous les coefficients par 2 :
\[ \ce{2 C6H6 + 15 O2 -> 12 CO2 + 6 H2O} \]
\item \textbf{Étape 5 (vérification)} : à gauche, C : $2 \times 6 = 12$ ; H : $2 \times 6 = 12$ ; O : $15 \times 2 = 30$. À droite, C : 12 ; H : $6 \times 2 = 12$ ; O : $12 \times 2 + 6 = 30$. L'équation est équilibrée.
\end{itemize}
\end{exemplebox}

\begin{attention}[Erreur fréquente : le coefficient fractionnaire oublié]
Beaucoup d'élèves s'arrêtent à l'équation $\ce{C6H6 + 15/2 O2 -> 6 CO2 + 3 H2O}$ et la recopient telle quelle sur leur copie. C'est \textbf{faux} en tant qu'équation-bilan finale : une équation-bilan s'écrit avec des coefficients \emph{entiers} (et les plus petits possibles). Dès que tu obtiens $\dfrac{15}{2}$ devant \ce{O2}, pense immédiatement à \textbf{doubler tous les coefficients}. Autre piège : ne doubler que le coefficient de \ce{O2} et oublier les autres — cela détruit la conservation des atomes ! Vérifie toujours chaque élément à la fin.
\end{attention}

\courssub{Une flamme très fuligineuse et éclairante : pourquoi ?}

\textbf{Observation.} Lorsqu'on enflamme quelques gouttes de benzène dans une coupelle (expérience réalisée uniquement par le professeur, sous hotte), la flamme est \emph{jaune-orangée, très éclairante}, et dégage d'\emph{épaisses fumées noires} qui déposent de la suie sur toute surface froide placée au-dessus.

\textbf{Interprétation.} Tout s'explique par la \cle{teneur en carbone} exceptionnellement élevée du benzène. Comparons les pourcentages massiques en carbone de quelques combustibles :

\begin{center}
\begin{tabularx}{\linewidth}{|l|c|c|X|}
\hline
\rowcolor{popblueL} \textbf{Combustible} & \textbf{Formule} & \textbf{\% massique en C} & \textbf{Aspect de la flamme} \\
\hline
Méthane & \ce{CH4} & $\dfrac{12}{16} = \SI{75}{\%}$ & bleue, propre, peu éclairante \\
\hline
Hexane & \ce{C6H14} & $\dfrac{72}{86} \approx \SI{84}{\%}$ & jaune, légèrement fuligineuse \\
\hline
Benzène & \ce{C6H6} & $\dfrac{72}{78} \approx \SI{92}{\%}$ & jaune vif, très fuligineuse et très éclairante \\
\hline
\end{tabularx}
\end{center}

Avec environ \SI{92}{\%} de carbone en masse, le benzène est l'un des hydrocarbures les plus « carbonés ». Or, dans la flamme réelle, l'apport local en dioxygène est souvent insuffisant pour oxyder instantanément tout ce carbone : une partie des atomes de carbone s'agrège en \emph{particules solides microscopiques} (les \cle{suies}). Portées à très haute température dans la flamme, ces particules incandescentes émettent une lumière jaune intense — d'où le caractère \emph{éclairant} — puis s'échappent sous forme de fumée noire — d'où le caractère \emph{fuligineux}.

\begin{popfigure}
\begin{center}
\begin{tikzpicture}[scale=1.0]
% Coupelle
\draw[thick, popdark] (-2.2,0) .. controls (-2.2,-0.9) and (2.2,-0.9) .. (2.2,0);
\draw[thick, popdark] (-2.2,0) -- (2.2,0);
\draw[thick, popdark] (-1.6,-0.55) -- (-1.6,-1.0);
\draw[thick, popdark] (1.6,-0.55) -- (1.6,-1.0);
\draw[thick, popdark] (-2.0,-1.0) -- (2.0,-1.0);
\node[popdark, font=\small] at (0,-1.45) {coupelle contenant du benzène};
% Flamme (enveloppes successives)
\fill[poporange!25] (0,0.1) .. controls (-1.5,1.2) and (-1.1,2.6) .. (0,3.6)
   .. controls (1.1,2.6) and (1.5,1.2) .. (0,0.1);
\fill[poporange!60] (0,0.15) .. controls (-1.0,1.0) and (-0.75,2.0) .. (0,2.8)
   .. controls (0.75,2.0) and (1.0,1.0) .. (0,0.15);
\fill[popgold] (0,0.2) .. controls (-0.55,0.9) and (-0.4,1.5) .. (0,2.1)
   .. controls (0.4,1.5) and (0.55,0.9) .. (0,0.2);
% Fumées noires (suies)
\draw[very thick, popmuted, decorate, decoration={coil, aspect=0.4, segment length=6pt, amplitude=4pt}] (0,3.6) -- (0,5.4);
\draw[very thick, popmuted!70, decorate, decoration={coil, aspect=0.4, segment length=6pt, amplitude=4pt}] (-0.5,3.4) -- (-0.9,5.2);
\draw[very thick, popmuted!70, decorate, decoration={coil, aspect=0.4, segment length=6pt, amplitude=4pt}] (0.5,3.4) -- (0.9,5.2);
\fill[popmuted] (-0.9,5.3) circle (0.10);
\fill[popmuted] (0,5.5) circle (0.12);
\fill[popmuted] (0.9,5.3) circle (0.10);
\fill[popmuted!80] (-0.3,5.1) circle (0.07);
\fill[popmuted!80] (0.4,5.0) circle (0.07);
% Annotations
\node[align=left, font=\small, text=popdark] at (4.6,2.4) {flamme jaune vif\\ \emph{très éclairante}};
\draw[-{Stealth}, popdark] (3.2,2.4) -- (0.9,2.2);
\node[align=left, font=\small, text=popmuted] at (4.6,4.9) {fumées noires :\\ particules de carbone\\ (suies) incandescentes};
\draw[-{Stealth}, popmuted] (3.1,4.9) -- (0.4,4.7);
\node[align=left, font=\small, text=popblue] at (-4.6,1.2) {benzène : \SI{92}{\%}\\ de carbone en masse};
\draw[-{Stealth}, popblue] (-3.0,1.2) -- (-1.2,0.5);
\end{tikzpicture}
\end{center}
\legende{Combustion du benzène : la forte teneur en carbone (\SI{92}{\%} en masse) provoque la formation de particules de carbone incandescentes, responsables d'une flamme jaune très éclairante et de fumées noires fuligineuses.}
\end{popfigure}

\begin{savaistu}[Les lampes à acétylène et les feux de brousse]
Le même phénomène explique pourquoi l'acétylène \ce{C2H2} (\SI{92,3}{\%} de carbone, comme le benzène !) brûle avec une flamme si lumineuse qu'on l'utilisait autrefois dans les lampes des mineurs et des cyclistes. À l'inverse, le gaz butane des bouteilles de cuisine, plus pauvre en carbone, brûle avec une flamme bleue propre lorsque le brûleur est bien réglé. Une flamme de gazinière qui jaunit et noircit les marmites signale un mauvais réglage de l'arrivée d'air : la combustion devient incomplète.
\end{savaistu}

\courssub{La combustion incomplète : un danger mortel}

\textbf{Intuition.} Que se passe-t-il si le dioxygène vient à manquer ? L'oxydation du carbone s'arrête « en chemin ». Au lieu du dioxyde de carbone \ce{CO2}, on obtient du \cle{monoxyde de carbone} \ce{CO}, et même du carbone solide \ce{C} (les suies) si le manque de dioxygène est sévère.

\begin{definition}[Combustion incomplète]
La \cle{combustion incomplète} d'un hydrocarbure se produit lorsque la quantité de dioxygène disponible est insuffisante. Elle produit du \cle{monoxyde de carbone} \ce{CO} (gaz incolore, inodore et \emph{très toxique}), de l'eau, et éventuellement du \cle{carbone} solide (suies noires). Pour le benzène, on peut écrire par exemple :
\[ \ce{2 C6H6 + 9 O2 -> 12 CO + 6 H2O} \qquad\text{ou, en cas de grand déficit en \ce{O2} :}\qquad \ce{2 C6H6 + 3 O2 -> 12 C + 6 H2O} \]
\end{definition}

\begin{attention}[Le monoxyde de carbone, tueur invisible]
Le monoxyde de carbone \ce{CO} est \emph{incolore, inodore et sans saveur} : impossible à détecter sans appareil. Inhalé, il se fixe sur l'hémoglobine du sang à la place du dioxygène (il s'y lie environ 300 fois plus fortement que \ce{O2}) et provoque l'asphyxie des cellules : maux de tête, vertiges, puis perte de conscience et mort. Chaque année, des intoxications surviennent lorsqu'on fait brûler du charbon de bois ou un groupe électrogène dans une pièce mal ventilée. \textbf{Règle de sécurité absolue} : toute combustion (cuisine, brasero, groupe électrogène) doit se faire dans un espace \emph{bien aéré}. Au laboratoire, la combustion du benzène ne se réalise jamais en dehors d'une hotte ventilée.
\end{attention}

\begin{exoresolu}[Volume de dioxygène nécessaire à la combustion du benzène]
\textbf{Énoncé.} On brûle complètement une masse $m = \SI{7,8}{g}$ de benzène \ce{C6H6} dans le dioxygène.
\begin{enumerate}
\item Écrire l'équation-bilan de la réaction.
\item Calculer la quantité de matière de benzène brûlé.
\item En déduire le volume de dioxygène nécessaire, mesuré dans les conditions où le volume molaire vaut $V_m = \SI{24,0}{L.mol^{-1}}$.
\item Calculer la masse de dioxyde de carbone produite.
\end{enumerate}
Données : $M(\text{C}) = \SI{12}{g.mol^{-1}}$ ; $M(\text{H}) = \SI{1}{g.mol^{-1}}$ ; $M(\text{O}) = \SI{16}{g.mol^{-1}}$.
\tcblower
\textbf{Solution détaillée.}
\begin{enumerate}
\item \textbf{Équation-bilan} (méthode C -- H -- O, puis doublement des coefficients) :
\[ \ce{2 C6H6 + 15 O2 -> 12 CO2 + 6 H2O} \]
\item \textbf{Quantité de benzène.} La masse molaire du benzène vaut :
\[ M(\ce{C6H6}) = 6 \times 12 + 6 \times 1 = \SI{78}{g.mol^{-1}} \]
D'où :
\[ n(\ce{C6H6}) = \frac{m}{M} = \frac{7{,}8}{78} = \SI{0,10}{mol} \]
\item \textbf{Volume de dioxygène.} Dressons le tableau d'avancement (réaction totale) :
\begin{center}
\begin{tabular}{lcccc}
\toprule
 & \ce{2 C6H6} & \ce{+ 15 O2} & \ce{-> 12 CO2} & \ce{+ 6 H2O} \\
\midrule
État initial (mol) & 0{,}10 & $n_0$ & 0 & 0 \\
État final (mol) & $0{,}10 - 2x_{max}$ & $n_0 - 15x_{max}$ & $12x_{max}$ & $6x_{max}$ \\
\bottomrule
\end{tabular}
\end{center}
Le benzène est entièrement consommé : $0{,}10 - 2x_{max} = 0$, donc $x_{max} = \SI{0,050}{mol}$.
La quantité de dioxygène \emph{nécessaire} (conditions stœchiométriques) est :
\[ n(\ce{O2}) = 15\,x_{max} = 15 \times 0{,}050 = \SI{0,75}{mol} \]
D'où le volume :
\[ V(\ce{O2}) = n \times V_m = 0{,}75 \times 24{,}0 = \SI{18}{L} \]
\item \textbf{Masse de dioxyde de carbone.} $n(\ce{CO2}) = 12\,x_{max} = 12 \times 0{,}050 = \SI{0,60}{mol}$. Avec $M(\ce{CO2}) = 12 + 2 \times 16 = \SI{44}{g.mol^{-1}}$ :
\[ m(\ce{CO2}) = 0{,}60 \times 44 = \SI{26,4}{g} \]
\end{enumerate}
\textbf{Conclusion.} Brûler seulement \SI{7,8}{g} de benzène (environ une cuillère à soupe) consomme \SI{18}{L} de dioxygène et rejette \SI{26,4}{g} de \ce{CO2} : les combustions sont de grosses consommatrices d'oxygène, d'où l'importance vitale de la ventilation.
\end{exoresolu}

\courssub{Le benzène n'est plus un carburant : un intermédiaire industriel majeur}

Au vu de son pouvoir énergétique élevé, on pourrait imaginer le benzène comme carburant. Historiquement, il a d'ailleurs été mélangé à l'essence pour augmenter son indice d'octane. Ce n'est \textbf{plus le cas aujourd'hui}, pour deux raisons majeures :

\begin{itemize}
\item \textbf{Toxicité} : le benzène est \emph{cancérogène avéré} (il provoque des leucémies) et sa combustion incomplète dégage du monoxyde de carbone et des suies cancérogènes. Sa teneur dans les essences est désormais limitée par la réglementation à moins de \SI{1}{\%} en volume.
\item \textbf{Valeur ajoutée} : le benzène vaut bien plus cher comme \emph{matière première} de l'industrie chimique que comme simple combustible.
\end{itemize}

\begin{aretenir}[Le benzène, intermédiaire industriel majeur]
Le benzène ne sert quasiment plus de carburant. Il est aujourd'hui l'un des \cle{intermédiaires industriels} les plus importants de la pétrochimie mondiale : il sert de point de départ à la fabrication des \emph{plastiques} (polystyrène, nylon via le cyclohexane), des \emph{caoutchoucs synthétiques}, des \emph{détergents}, des \emph{colorants}, des \emph{explosifs} (comme le TNT, que nous préparerons par nitration dans la suite de la leçon), des \emph{insecticides} et de nombreux \emph{médicaments}. Les réactions qui permettent ces transformations — additions et surtout substitutions — feront l'objet des sections suivantes.
\end{aretenir}

\begin{savaistu}[Pétrochimie et aromatiques au Cameroun]
La SONARA (Société Nationale de Raffinage) de Limbé traite le pétrole brut pour produire essences, gasoil et kérosène. Le raffinage et le reformage catalytique des coupes pétrolières sont précisément les procédés industriels qui génèrent les hydrocarbures aromatiques (benzène, toluène, xylènes — le fameux « BTX »), ensuite acheminés vers les industries chimiques du monde entier pour fabriquer plastiques, fibres et médicaments. Comprendre la chimie du benzène, c'est comprendre l'une des portes d'entrée de l'industrie moderne.
\end{savaistu}

\begin{exercice}[Combustion du toluène]
Le toluène, de formule brute \ce{C7H8}, est le dérivé méthylé du benzène.
\begin{enumerate}
\item Calculer le pourcentage massique en carbone du toluène et prévoir qualitativement l'aspect de sa flamme.
\item Écrire et équilibrer l'équation-bilan de sa combustion complète dans le dioxygène (on éliminera tout coefficient fractionnaire).
\item Quelle masse de toluène faut-il brûler pour produire \SI{1,4}{mol} de dioxyde de carbone ?
\end{enumerate}
Données : $M(\text{C}) = \SI{12}{g.mol^{-1}}$ ; $M(\text{H}) = \SI{1}{g.mol^{-1}}$.
\end{exercice}

\begin{corrige}[Exercice : combustion du toluène]
\begin{enumerate}
\item $M(\ce{C7H8}) = 7 \times 12 + 8 \times 1 = \SI{92}{g.mol^{-1}}$. Pourcentage massique en carbone : $\dfrac{7 \times 12}{92} = \dfrac{84}{92} \approx \SI{91,3}{\%}$. Cette teneur est très proche de celle du benzène (\SI{92}{\%}) : la flamme du toluène est donc également \emph{jaune, très fuligineuse et éclairante}.
\item Méthode C -- H -- O avec $x = 7$, $y = 8$ : coefficient de \ce{CO2} : 7 ; coefficient de \ce{H2O} : $\dfrac{8}{2} = 4$ ; oxygène à droite : $7 \times 2 + 4 = 18$, d'où $\dfrac{18}{2} = 9$ devant \ce{O2}. Le coefficient est déjà entier :
\[ \ce{C7H8 + 9 O2 -> 7 CO2 + 4 H2O} \]
\item D'après l'équation, 1 mol de toluène produit 7 mol de \ce{CO2}. Il faut donc $n(\ce{C7H8}) = \dfrac{1{,}4}{7} = \SI{0,20}{mol}$, soit une masse $m = 0{,}20 \times 92 = \SI{18,4}{g}$ de toluène.
\end{enumerate}
\end{corrige}

\coursec{Réactions d'addition sur le benzène}

Dans la section précédente, nous avons vu que la combustion du benzène détruit totalement la molécule : le noyau benzénique est brisé et tout le carbone finit en dioxyde de carbone. Nous allons maintenant étudier une famille de réactions plus douces, mais qui ont un point commun avec la combustion : elles font \emph{disparaître le caractère aromatique} du benzène. Ce sont les \cle{réactions d'addition}. Tu vas découvrir pourquoi elles sont difficiles, comment les réaliser, et comment l'une d'elles a donné naissance à un insecticide célèbre... et redouté.

\courssub{Qu'est-ce qu'une réaction d'addition sur le benzène ?}

\courssub{Intuition : casser un édifice très stable}

Imagine le noyau benzénique comme un édifice dont les six électrons $\pi$ délocalisés forment le ciment. Cet édifice est exceptionnellement solide : c'est précisément cette stabilité qui explique que le benzène ne se comporte \emph{pas} comme un alcène ordinaire, alors que sa formule brute $\ce{C6H6}$ laisse croire à la présence de trois doubles liaisons.

Une réaction d'addition consiste à fixer des atomes supplémentaires sur la molécule. Mais pour cela, il faut \textbf{rompre la délocalisation} des électrons $\pi$ : autrement dit, démolir le ciment de l'édifice. Cela coûte cher en énergie. Voilà l'idée directrice de toute cette section :

\begin{definition}[Réaction d'addition sur le benzène]
Une \cle{réaction d'addition} sur le benzène est une réaction au cours de laquelle des atomes (hydrogène, chlore, etc.) se fixent sur les atomes de carbone du cycle, avec \textbf{disparition du caractère aromatique} : les électrons $\pi$ ne sont plus délocalisés et le cycle devient un cycloalcane (ou un dérivé de cycloalcane).
\end{definition}

\begin{aretenir}[L'essentiel]
L'addition sur le benzène est \textbf{difficile} car elle détruit la délocalisation des électrons $\pi$, source de la grande stabilité du noyau aromatique. Elle exige donc des \textbf{conditions énergétiques} : catalyseur, chauffage ou lumière ultraviolette.
\end{aretenir}

\courssub{Comparaison avec les alcènes : pourquoi le benzène résiste}

Tu as étudié les alcènes : tu sais qu'ils décolorent instantanément l'eau de brome et qu'ils s'hydrogènent facilement en présence d'un catalyseur, à température modérée. Le benzène, lui, \textbf{ne décolore pas l'eau de brome} et ne réagit avec le dihydrogène que dans des conditions bien plus sévères.

\begin{center}
\begin{tabularx}{\linewidth}{|l|X|X|}
\hline
\rowcolor{popblueL} \textbf{Critère} & \textbf{Alcène (ex. éthène)} & \textbf{Benzène} \\
\hline
Liaisons insaturées & 1 double liaison localisée & 6 électrons $\pi$ délocalisés \\
\hline
Eau de brome & Décoloration immédiate & Aucune réaction \\
\hline
Hydrogénation & Facile (Ni, température douce) & Difficile (Ni, Pt ou Raney, chauffage) \\
\hline
Produit de l'addition & Alcane & Cycloalcane (aromaticité perdue) \\
\hline
\end{tabularx}
\end{center}

\begin{attention}[Le piège classique]
Ne dis jamais que le benzène « possède trois doubles liaisons comme trois alcènes ». S'il en était ainsi, il réagirait aussi vite qu'un alcène. C'est précisément parce que les électrons $\pi$ sont \textbf{délocalisés} sur tout le cycle que le benzène est beaucoup moins réactif en addition. Cette stabilité exceptionnelle s'appelle la \cle{stabilité aromatique} (ou énergie de stabilisation aromatique, environ \SI{150}{kJ.mol^{-1}}).
\end{attention}

\courssub{Hydrogénation du benzène : obtention du cyclohexane}

\courssub{La réaction}

En présence de dihydrogène et d'un catalyseur métallique (nickel, platine ou nickel de Raney), et sous chauffage, le benzène fixe \textbf{trois molécules de dihydrogène} : le cycle se sature complètement et l'on obtient le \cle{cyclohexane}, de formule $\ce{C6H12}$.

\[ \ce{C6H6 + 3 H2 ->[Ni,\ \text{chauffage}] C6H12} \]

\begin{popfigure}
\begin{center}
\chemfig{*6(=-=-=-)} \hspace{1cm} $+ \; 3\,\mathrm{H_2}$ \hspace{1cm}
\raisebox{0.6em}{\tikz{\draw[-{Stealth},very thick,popink] (0,0) -- (3,0) node[midway,above,popdark]{\small Ni ou Pt, $\theta$ élevée};}}
\hspace{1cm} \chemfig{*6(------)}
\end{center}
\legende{Hydrogénation catalytique du benzène : le cycle aromatique (à gauche, représentation de Kekulé avec doubles liaisons alternées) est transformé en cyclohexane saturé (à droite, liaisons simples seulement). Le catalyseur (nickel, platine ou nickel de Raney) et le chauffage fournissent l'énergie nécessaire pour rompre la délocalisation.}
\end{popfigure}

\textbf{Commentaires sur les conditions opératoires :}
\begin{itemize}
\item le \textbf{catalyseur} (Ni, Pt ou nickel de Raney) adsorbe le dihydrogène à sa surface et fragilise la liaison $\ce{H-H}$, ce qui abaisse l'énergie d'activation ;
\item le \textbf{chauffage} apporte l'énergie indispensable pour vaincre la stabilité aromatique ;
\item l'addition est \textbf{totale} : on n'obtient jamais de cyclohexène intermédiaire isolé dans ces conditions, les trois « doubles liaisons » sont hydrogénées d'un coup.
\end{itemize}

\begin{attention}[Erreur fréquente : le nombre de moles de dihydrogène]
Beaucoup d'élèves écrivent $\ce{C6H6 + H2 -> C6H8}$. C'est \textbf{faux} ! Le cycle benzénique comporte l'équivalent de \textbf{trois} doubles liaisons : il faut donc \textbf{3 moles de $\ce{H2}$} pour saturer complètement le cycle. Retiens : $\ce{C6H6 + 3 H2 -> C6H12}$. Vérifie toujours la conservation des atomes : 6 C et $6 + 6 = 12$ H de chaque côté.
\end{attention}

\begin{exoresolu}[Fabrication industrielle du cyclohexane]
Une unité pétrochimique traite \SI{78}{g} de benzène par hydrogénation totale en présence de nickel de Raney. On suppose le rendement de la réaction égal à \SI{100}{\%}.

\textbf{1.} Écrire l'équation-bilan de la réaction.

\textbf{2.} Calculer la quantité de matière de benzène utilisée.

\textbf{3.} En déduire le volume de dihydrogène nécessaire, mesuré dans les conditions normales de température et de pression ($V_m = \SI{22,4}{L.mol^{-1}}$), puis la masse de cyclohexane obtenue.

\emph{Données :} $M(\ce{C}) = \SI{12}{g.mol^{-1}}$ ; $M(\ce{H}) = \SI{1}{g.mol^{-1}}$.

\tcblower

\textbf{1. Équation-bilan :}
\[ \ce{C6H6 + 3 H2 ->[Ni\ Raney,\ \text{chauffage}] C6H12} \]

\textbf{2. Quantité de matière de benzène :}

La masse molaire du benzène vaut :
\[ M(\ce{C6H6}) = 6 \times 12 + 6 \times 1 = \SI{78}{g.mol^{-1}} \]
\[ n(\ce{C6H6}) = \frac{m}{M} = \frac{78}{78} = \SI{1,0}{mol} \]

\textbf{3. Volume de dihydrogène et masse de cyclohexane :}

Dressons le tableau d'avancement (réaction supposée totale) :

\begin{center}
\begin{tabularx}{\linewidth}{|l|X|X|X|}
\hline
\rowcolor{popblueL} & $\ce{C6H6}$ & $\ce{3 H2}$ & $\ce{C6H12}$ \\
\hline
État initial (mol) & $1{,}0$ & $3{,}0$ (excès) & $0$ \\
\hline
État final (mol) & $0$ & $0$ & $1{,}0$ \\
\hline
\end{tabularx}
\end{center}

D'après les coefficients stœchiométriques, il faut $3$ moles de $\ce{H2}$ pour $1$ mole de benzène :
\[ n(\ce{H2}) = 3 \times 1{,}0 = \SI{3,0}{mol} \]
\[ V(\ce{H2}) = n \times V_m = 3{,}0 \times 22{,}4 = \SI{67,2}{L} \]

Il se forme $1{,}0$ mole de cyclohexane. Sa masse molaire est :
\[ M(\ce{C6H12}) = 6 \times 12 + 12 \times 1 = \SI{84}{g.mol^{-1}} \]
\[ m(\ce{C6H12}) = n \times M = 1{,}0 \times 84 = \SI{84}{g} \]

\textbf{Conclusion :} il faut \SI{67,2}{L} de dihydrogène (CNTP) et l'on obtient \SI{84}{g} de cyclohexane. Remarque la cohérence : $78 + 3 \times 2 = 84$ g, la masse se conserve bien.
\end{exoresolu}

\begin{savaistu}[Le cyclohexane, matière première du nylon]
Le cyclohexane produit par hydrogénation du benzène est un intermédiaire industriel majeur : oxydé, il conduit à l'acide adipique et au caprolactame, les briques de fabrication du \textbf{nylon}. Chaque année, des millions de tonnes de benzène subissent ainsi cette hydrogénation dans les raffineries du monde entier, y compris dans les unités de traitement pétrochimique comme celles associées à la SONARA à Limbe.
\end{savaistu}

\courssub{Chloration du benzène : obtention de l'hexachlorocyclohexane}

\courssub{La réaction sous lumière ultraviolette}

Sous l'action de la \textbf{lumière ultraviolette} (UV) et en l'absence de catalyseur, le benzène peut fixer \textbf{trois molécules de dichlore} : six atomes de chlore se fixent sur les six atomes de carbone du cycle. On obtient l'\cle{hexachlorocyclohexane}, de formule brute $\ce{C6H6Cl6}$.

\[ \ce{C6H6 + 3 Cl2 ->[UV] C6H6Cl6} \]

\begin{popfigure}
\begin{center}
\chemfig{*6(=-=-=-)} \hspace{1cm} $+ \; 3\,\mathrm{Cl_2}$ \hspace{1cm}
\raisebox{0.6em}{\tikz{\draw[-{Stealth},very thick,poporange] (0,0) -- (3,0) node[midway,above,popdark]{\small lumière UV};}}
\hspace{1cm} \chemfig{*6((-Cl)-(-Cl)-(-Cl)-(-Cl)-(-Cl)-(-Cl))}
\end{center}
\legende{Chloration du benzène sous lumière ultraviolette : trois molécules de dichlore s'additionnent sur le cycle. Chaque carbone porte désormais un atome d'hydrogène et un atome de chlore : le produit est l'hexachlorocyclohexane $\ce{C6H6Cl6}$, un cycloalcane chloré qui n'a plus rien d'aromatique.}
\end{popfigure}

\textbf{Commentaires :}
\begin{itemize}
\item le rôle des \textbf{UV} est de rompre la liaison $\ce{Cl-Cl}$ par homolyse (formation de radicaux $\ce{Cl.}$) et de fournir l'énergie nécessaire à l'attaque du cycle aromatique ;
\item le produit obtenu est un \textbf{cyclohexane hexasubstitué} : chacun des six carbones porte un H et un Cl, d'où la formule $\ce{C6H6Cl6}$ ;
\item il existe plusieurs \textbf{stéréoisomères} de l'hexachlorocyclohexane selon la position relative des atomes de chlore (axiaux/équatoriaux) ; l'isomère $\gamma$, appelé \cle{lindane}, est le plus actif comme insecticide.
\end{itemize}

\begin{attention}[Addition ou substitution ? Regarde les conditions !]
Le couple benzène + dichlore peut donner deux réactions \textbf{différentes} selon les conditions :
\begin{itemize}
\item \textbf{lumière UV, sans catalyseur} $\Rightarrow$ \textbf{addition} : $\ce{C6H6 + 3 Cl2 -> C6H6Cl6}$ ;
\item \textbf{présence d'un catalyseur} ($\ce{FeCl3}$, $\ce{AlCl3}$), à l'abri de la lumière $\Rightarrow$ \textbf{substitution} : un atome de chlore remplace un atome d'hydrogène, on obtient le chlorobenzène $\ce{C6H5Cl}$ (nous l'étudierons dans la prochaine section).
\end{itemize}
Retiens ce réflexe : \emph{UV = addition ; catalyseur = substitution}.
\end{attention}

\courssub{L'hexachlorocyclohexane : un insecticide célèbre}

\begin{savaistu}[Le HCH, insecticide miracle devenu polluant redouté]
L'hexachlorocyclohexane (HCH), et particulièrement son isomère $\gamma$ appelé \textbf{lindane}, a été massivement utilisé à partir des années 1940 comme \textbf{insecticide} : contre les criquets ravageurs de cultures, les moustiques vecteurs du paludisme, les poux et les termites. Bon marché et très efficace, il a été épandu en quantités énormes dans les zones tropicales, y compris en Afrique centrale.

Mais on a découvert ensuite que le HCH est \textbf{persistant} dans l'environnement (il se dégrade très lentement), qu'il \textbf{s'accumule} dans les graisses des êtres vivants tout au long de la chaîne alimentaire (bioaccumulation), et qu'il est \textbf{toxique} pour le système nerveux et suspecté d'être cancérigène. Il figure aujourd'hui parmi les \emph{polluants organiques persistants} (POP) interdits ou sévèrement réglementés par la convention de Stockholm (2001). C'est un exemple frappant de la double face de la chimie : une molécule peut rendre d'immenses services... à condition d'en maîtriser les risques.
\end{savaistu}

Cet exemple illustre parfaitement notre famille de situations : la \textbf{préparation d'un insecticide} repose ici sur une réaction d'addition du dichlore sur le benzène, réalisée sous irradiation ultraviolette.

\courssub{Interprétation énergétique : pourquoi l'addition coûte-t-elle si cher ?}

Rassemblons les faits observés :
\begin{itemize}
\item l'hydrogénation exige un \textbf{catalyseur} et un \textbf{chauffage} ;
\item la chloration exige la \textbf{lumière UV} ;
\item le benzène, contrairement aux alcènes, \textbf{ne réagit pas} avec le brome à froid.
\end{itemize}

Toutes ces observations s'expliquent par une seule cause : la \textbf{délocalisation des six électrons $\pi$} sur l'ensemble du cycle confère au benzène une stabilité supplémentaire considérable, appelée \emph{énergie de stabilisation aromatique} (environ \SI{150}{kJ.mol^{-1}}). Or, toute addition \textbf{supprime cette délocalisation} : les électrons $\pi$ deviennent localisés dans des liaisons simples, et le produit obtenu (cyclohexane ou hexachlorocyclohexane) n'est plus aromatique.

\begin{propriete}[Règle générale de réactivité du benzène]
L'addition sur le benzène \textbf{détruit le caractère aromatique} : elle est donc \textbf{difficile} et n'a lieu que dans des conditions énergétiques (catalyseur + chauffage, ou lumière UV). C'est pourquoi le benzène préfère, lorsque c'est possible, les réactions de \textbf{substitution}, qui \emph{conservent} le noyau aromatique — ce sera l'objet de la section suivante.
\end{propriete}

\begin{methode}[Écrire et équilibrer une équation d'addition sur le benzène]
\begin{enumerate}
\item Identifier le réactif qui s'additionne : $\ce{H2}$ (hydrogénation) ou $\ce{Cl2}$ (chloration).
\item Se rappeler que le cycle benzénique comporte l'équivalent de \textbf{trois} doubles liaisons : il faut donc \textbf{3 molécules} de réactif par mole de benzène.
\item Écrire le produit : cyclohexane $\ce{C6H12}$ avec $\ce{H2}$ ; hexachlorocyclohexane $\ce{C6H6Cl6}$ avec $\ce{Cl2}$.
\item Préciser les \textbf{conditions} au-dessus de la flèche : Ni/Pt/Raney + chauffage pour $\ce{H2}$ ; UV pour $\ce{Cl2}$.
\item Vérifier la conservation des atomes des deux côtés de la flèche.
\end{enumerate}
\end{methode}

\begin{exercice}[Pour t'entraîner]
On fait réagir \SI{39}{g} de benzène avec du dichlore en présence de lumière ultraviolette, en supposant la réaction totale.

\textbf{1.} Écrire l'équation-bilan de la réaction et nommer le produit obtenu.

\textbf{2.} Calculer la quantité de matière de dichlore nécessaire.

\textbf{3.} Calculer la masse d'hexachlorocyclohexane obtenue.

\emph{Données :} $M(\ce{C}) = \SI{12}{g.mol^{-1}}$ ; $M(\ce{H}) = \SI{1}{g.mol^{-1}}$ ; $M(\ce{Cl}) = \SI{35,5}{g.mol^{-1}}$.
\end{exercice}

\begin{corrige}[Exercice]
\textbf{1.} Équation-bilan : $\ce{C6H6 + 3 Cl2 ->[UV] C6H6Cl6}$. Le produit est l'\textbf{hexachlorocyclohexane} (HCH), un insecticide.

\textbf{2.} $M(\ce{C6H6}) = 78$ g.mol$^{-1}$, donc $n(\ce{C6H6}) = 39/78 = \SI{0,50}{mol}$. D'après la stœchiométrie : $n(\ce{Cl2}) = 3 \times 0{,}50 = \SI{1,5}{mol}$.

\textbf{3.} $M(\ce{C6H6Cl6}) = 6 \times 12 + 6 \times 1 + 6 \times 35{,}5 = 72 + 6 + 213 = \SI{291}{g.mol^{-1}}$. Il se forme $0{,}50$ mol de HCH : $m = 0{,}50 \times 291 = \SI{145,5}{g}$.
\end{corrige}

\begin{aretenir}[Bilan de la section]
\begin{itemize}
\item Une \textbf{addition} sur le benzène fixe des atomes sur le cycle et \textbf{détruit le caractère aromatique} : elle est difficile et exige des conditions énergétiques.
\item \textbf{Hydrogénation} : $\ce{C6H6 + 3 H2 ->[Ni/Pt,\ \text{chauffage}] C6H12}$ (cyclohexane).
\item \textbf{Chloration} : $\ce{C6H6 + 3 Cl2 ->[UV] C6H6Cl6}$ (hexachlorocyclohexane, insecticide).
\item La difficulté de l'addition s'explique par la \textbf{délocalisation des électrons $\pi$}, que l'addition supprime : c'est le prix à payer pour saturer le cycle.
\end{itemize}
\end{aretenir}

\coursec{Réactions de substitution : la réactivité privilégiée du benzène}

Dans la section précédente, nous avons vu que le benzène peut, dans des conditions énergiques (rayonnement UV, hydrogénation catalytique à chaud), subir des réactions d'\cle{addition} qui détruisent le nuage électronique délocalisé : le cyclohexane ou le 1,2,3,4,5,6-hexachlorocyclohexane obtenus ne sont plus aromatiques. Mais ces réactions sont \emph{l'exception}. Le comportement le plus naturel, le plus fréquent du benzène est tout autre : plutôt que de casser son précieux système $\pi$ délocalisé, il préfère \textbf{remplacer l'un de ses atomes d'hydrogène} par un autre atome ou groupe d'atomes, tout en gardant intact son noyau aromatique. C'est la \cle{réaction de substitution}, la réaction caractéristique des composés aromatiques. C'est elle qui permet de fabriquer le TNT, les insecticides, les sulfamides et bien des matières plastiques.

\courssub{Principe général de la substitution}

\begin{definition}[Réaction de substitution sur le benzène]
Une \cle{réaction de substitution} est une réaction au cours de laquelle \textbf{un atome d'hydrogène du noyau benzénique est remplacé par un autre atome ou groupe d'atomes} $X$, avec \textbf{conservation du caractère aromatique} du cycle. L'équation générale s'écrit :
\[ \mathrm{C_6H_6 + X_2 \xrightarrow{\text{catalyseur}} C_6H_5X + HX} \]
Le produit obtenu est un \cle{dérivé monosubstitué} du benzène, encore aromatique.
\end{definition}

L'idée intuitive est simple : le nuage $\pi$ délocalisé confère au benzène une stabilité exceptionnelle (environ \SI{150}{kJ.mol^{-1}} de « bonus » de stabilité par rapport à un cycle hypothétique à trois doubles liaisons localisées). Lors d'une addition, cette stabilité est perdue ; lors d'une substitution, elle est préservée. Le benzène « choisit » donc la voie qui conserve son nuage $\pi$ : \textbf{le benzène préfère la substitution à l'addition}.

\begin{popfigure}
\centering
\chemfig{**6(------)} \hspace{0,8cm} $+$ \hspace{0,3cm} $X_2$ \hspace{0,8cm} $\xrightarrow{\text{catalyseur}}$ \hspace{0,8cm} \chemfig{**6((-X)-----)} \hspace{0,8cm} $+$ \hspace{0,3cm} HX
\legende{Schéma général d'une substitution électrophile sur le benzène : un hydrogène du noyau est remplacé par le groupe $X$ ; le cycle aromatique est conservé.}
\end{popfigure}

\begin{methode}[Écrire correctement l'équation d'une substitution]
\begin{enumerate}
\item \textbf{Identifier le groupe entrant} ($\ce{Cl}$, $\ce{Br}$, $\ce{NO2}$, $\ce{SO3H}$, groupe alkyle $R$\ldots) apporté par le réactif ($\ce{Cl2}$, $\ce{HNO3}$, $\ce{H2SO4}$, $\ce{RCl}$\ldots).
\item \textbf{Remplacer un H du noyau} par ce groupe : le benzène \ce{C6H6} devient \ce{C6H5-G} où G est le groupe entrant.
\item \textbf{Écrire le sous-produit} : l'hydrogène expulsé se combine avec le reste du réactif pour donner \ce{HCl}, \ce{HBr} ou \ce{H2O}.
\item \textbf{Préciser les conditions} au-dessus de la flèche : catalyseur (\ce{FeCl3}, \ce{AlCl3}, \ce{H2SO4} concentré) et température éventuelle.
\item \textbf{Vérifier l'équilibre} de l'équation en comptant les atomes de chaque côté.
\end{enumerate}
\end{methode}

\courssub{L'halogénation : chloration et bromuration}

En présence d'un catalyseur appelé \cle{acide de Lewis} (trichlorure de fer \ce{FeCl3} ou trichlorure d'aluminium \ce{AlCl3} pour le chlore ; tribromure de fer \ce{FeBr3} pour le brome), le benzène réagit avec les dihalogènes \emph{à froid et à l'abri de la lumière} selon une substitution :

\[ \ce{C6H6 + Cl2 ->[FeCl3] C6H5Cl + HCl} \]
\[ \ce{C6H6 + Br2 ->[FeBr3] C6H5Br + HBr} \]

Le produit est le \textbf{chlorobenzène} (ou le \textbf{bromobenzène}), et l'on observe un dégagement de chlorure d'hydrogène (ou de bromure d'hydrogène), gaz qui fume à l'air humide. Le rôle du catalyseur est de polariser la molécule \ce{Cl2} pour générer une espèce électrophile $\ce{Cl+}$ capable d'attaquer le nuage $\pi$.

\begin{attention}[Addition ou substitution ? Ne confonds pas les conditions !]
C'est le piège classique du Probatoire. Le \textbf{même couple de réactifs} (benzène + dichlore) donne deux réactions différentes selon les conditions :
\begin{itemize}
\item \textbf{Sous lumière UV, sans catalyseur} : \emph{addition} de 3 molécules de \ce{Cl2} $\rightarrow$ hexachlorocyclohexane \ce{C6H6Cl6} (le cycle n'est plus aromatique) ;
\item \textbf{À l'ombre, en présence de \ce{FeCl3}} : \emph{substitution} $\rightarrow$ chlorobenzène \ce{C6H5Cl} + \ce{HCl} (le cycle reste aromatique).
\end{itemize}
Retiens : \textbf{catalyseur (\ce{FeCl3}) = substitution ; UV = addition}. Écrire une addition de \ce{Cl2} en présence de \ce{FeCl3} est une erreur grave.
\end{attention}

\courssub{La nitration : vers les explosifs}

La \cle{nitration} introduit le groupe nitro $\ce{-NO2}$ sur le noyau. On utilise l'acide nitrique concentré \ce{HNO3} en présence d'acide sulfurique concentré \ce{H2SO4} (qui joue le rôle de catalyseur en générant l'ion nitronium $\ce{NO2+}$, l'espèce électrophile active), en chauffant modérément vers \SIrange{50}{60}{\celsius} :

\[ \ce{C6H6 + HNO3 ->[\ce{H2SO4}\ \text{conc.}][\SIrange{50}{60}{\celsius}] C6H5NO2 + H2O} \]

Le produit est le \textbf{nitrobenzène}, liquide jaune pâle à odeur d'amande amère, intermédiaire de la fabrication de l'aniline (teintures).

\begin{popfigure}
\centering
\chemfig{**6(------)} \hspace{0,6cm} $+$ \hspace{0,2cm} \ce{HNO3} \hspace{0,6cm} $\xrightarrow[\SIrange{50}{60}{\celsius}]{\ce{H2SO4}\ \text{conc.}}$ \hspace{0,6cm} \chemfig{**6((-NO_2)-----)} \hspace{0,6cm} $+$ \hspace{0,2cm} \ce{H2O}
\legende{Nitration du benzène : le groupe nitro $\ce{-NO2}$ remplace un hydrogène ; l'acide sulfurique concentré est le catalyseur, l'eau est le sous-produit.}
\end{popfigure}

\begin{exoresolu}[Nitration du benzène : masse de nitrobenzène obtenue]
On réalise la nitration de $m = \SI{39}{g}$ de benzène avec un excès d'acide nitrique. Le rendement de la réaction est de \SI{100}{\%}. Calculer la masse de nitrobenzène obtenue. \par
\emph{Données :} $M(\ce{C6H6}) = \SI{78}{g.mol^{-1}}$ ; $M(\ce{C6H5NO2}) = \SI{123}{g.mol^{-1}}$.
\tcblower
\textbf{Solution détaillée.}\par
\textbf{Étape 1 : équation de la réaction.}
\[ \ce{C6H6 + HNO3 ->[\ce{H2SO4}] C6H5NO2 + H2O} \]
\textbf{Étape 2 : quantité de matière de benzène.}
\[ n(\ce{C6H6}) = \frac{m}{M} = \frac{39}{78} = \SI{0,50}{mol} \]
\textbf{Étape 3 : tableau d'avancement.} L'acide nitrique étant en excès, le benzène est le réactif limitant.
\begin{center}
\begin{tabular}{|l|c|c|c|c|}
\hline
\rowcolor{popblueL} Équation & \ce{C6H6} & \ce{HNO3} & \ce{C6H5NO2} & \ce{H2O} \\
\hline
État initial (mol) & 0,50 & excès & 0 & 0 \\
\hline
État final (mol) & $0,50 - x_{max}$ & excès & $x_{max}$ & $x_{max}$ \\
\hline
\end{tabular}
\end{center}
Le benzène disparaît totalement : $0,50 - x_{max} = 0$, donc $x_{max} = \SI{0,50}{mol}$.
D'après la stœchiométrie (coefficients 1 pour 1) :
\[ n(\ce{C6H5NO2}) = x_{max} = \SI{0,50}{mol} \]
\textbf{Étape 4 : masse de nitrobenzène.}
\[ m(\ce{C6H5NO2}) = n \times M = 0,50 \times 123 = \SI{61,5}{g} \]
\textbf{Conclusion :} la nitration de \SI{39}{g} de benzène fournit, avec un rendement de \SI{100}{\%}, $\boxed{\SI{61,5}{g}}$ de nitrobenzène.
\end{exoresolu}

\courssub{Nitrations successives du toluène : la fabrication du TNT}

Le groupe méthyle $\ce{-CH3}$ du toluène \emph{active} le noyau benzénique : la nitration du toluène est plus facile que celle du benzène et peut se poursuivre jusqu'à l'introduction de \textbf{trois groupes nitro}, aux positions 2, 4 et 6. On obtient le \textbf{2,4,6-trinitrotoluène}, célèbre sous le nom de \cle{TNT}, puissant explosif.

\begin{popfigure}
\centering
\chemfig{**6((-CH_3)-----)} \hspace{0,4cm} $\xrightarrow{\ce{HNO3}/\ce{H2SO4}}$ \hspace{0,4cm} \chemfig{**6((-CH_3)(-NO_2)----)} \hspace{0,4cm} $\xrightarrow{\ce{HNO3}/\ce{H2SO4}}$ \hspace{0,4cm} \chemfig{**6((-CH_3)(-NO_2)(-H)(-NO_2)---)} \hspace{0,4cm} $\xrightarrow{\ce{HNO3}/\ce{H2SO4}}$ \hspace{0,4cm} \chemfig{**6((-CH_3)(-NO_2)(-H)(-NO_2)(-H)(-NO_2))}
\legende{Nitrations successives du toluène : toluène $\rightarrow$ 2-nitrotoluène (ortho) $\rightarrow$ 2,4-dinitrotoluène $\rightarrow$ 2,4,6-trinitrotoluène (TNT). Chaque étape consomme \ce{HNO3} et libère \ce{H2O}.}
\end{popfigure}

L'équation globale de la trinitration s'écrit :
\[ \ce{C6H5CH3 + 3 HNO3 ->[\ce{H2SO4}\ \text{conc.}] C6H2(NO2)3CH3 + 3 H2O} \]

\begin{attention}[Sécurité avec les dérivés nitrés]
Les dérivés polynitrés du benzène et du toluène sont des \textbf{explosifs} ou des précurseurs d'explosifs : leur manipulation est strictement réservée aux professionnels, dans des conditions contrôlées (température, absence de choc et de flamme). De plus, les vapeurs de benzène sont \textbf{cancérigènes} : toute manipulation s'effectue sous hotte, avec gants et lunettes, loin de toute flamme.
\end{attention}

\courssub{La sulfonation}

En chauffant le benzène avec l'acide sulfurique concentré \ce{H2SO4} (ou mieux, avec l'anhydride sulfurique \ce{SO3} / \ce{H2SO4} fumant), on introduit le groupe sulfonique $\ce{-SO3H}$ sur le noyau. Cette réaction est \textbf{réversible} :

\[ \ce{C6H6 + H2SO4 <=>[\text{chauffage}] C6H5SO3H + H2O} \]

Le produit est l'\textbf{acide benzènesulfonique}. Cette réaction est à la base de la fabrication de nombreux détergents (alkylbenzènesulfonates) et de sulfamides, premiers antibiotiques de l'histoire.

\begin{attention}[Réversibilité de la sulfonation]
Contrairement à la nitration ou l'halogénation, la sulfonation est un \textbf{équilibre}. Le chauffage favorise le sens direct (sulfonation) ; la dilution aqueuse et le réchauffage favorisent le sens inverse (désulfonation). Cette propriété est exploitée en synthèse pour protéger une position du noyau (groupe $\ce{-SO3H}$ encombrant, directeur \emph{méta}) avant de le retirer ensuite.
\end{attention}

\courssub{La réaction de Friedel et Crafts : alkylation du benzène}

Découverte en 1877 par Charles Friedel et James Crafts, cette réaction permet de fixer un \textbf{groupe alkyle} $R-$ sur le noyau benzénique, en faisant réagir le benzène avec un \cle{chlorure d'alkyle} $\ce{R-Cl}$ en présence de chlorure d'aluminium anhydre \ce{AlCl3} comme catalyseur :

\[ \ce{C6H6 + R-Cl ->[AlCl3] C6H5-R + HCl} \]

\textbf{Exemple fondamental : la préparation du toluène.} Avec le chlorométhane \ce{CH3Cl} ($R = \ce{CH3}$), on obtient le toluène :

\[ \ce{C6H6 + CH3Cl ->[AlCl3] C6H5CH3 + HCl} \]

\begin{popfigure}
\centering
\chemfig{**6(------)} \hspace{0,6cm} $+$ \hspace{0,2cm} \ce{CH3Cl} \hspace{0,6cm} $\xrightarrow{\ce{AlCl3}}$ \hspace{0,6cm} \chemfig{**6((-CH_3)-----)} \hspace{0,6cm} $+$ \hspace{0,2cm} \ce{HCl}
\legende{Réaction de Friedel et Crafts : le benzène réagit avec le chlorométhane en présence de \ce{AlCl3} pour donner le toluène et du chlorure d'hydrogène.}
\end{popfigure}

\begin{savaistu}[Friedel et Crafts : une réaction centenaire toujours industrielle]
Mise au point en 1877, la réaction de Friedel et Crafts est encore utilisée massivement aujourd'hui : l'industrie y a recours pour fabriquer l'\textbf{éthylbenzène} \ce{C6H5-C2H5} (à partir du benzène et du chloroéthane ou de l'éthylène), précurseur du \textbf{styrène}, lui-même monomère du \textbf{polystyrène} — la matière plastique des gobelets, des emballages protecteurs et des isolants. Une découverte du XIX\textsuperscript{e} siècle qui remplit encore nos poubelles\ldots et nos laboratoires !
\end{savaistu}

\begin{attention}[Limites de l'alkylation de Friedel-Crafts]
Deux phénomènes limitent cette réaction en synthèse fine :
\begin{itemize}
\item \textbf{Polyalkylation :} le produit (ex. toluène) est \emph{plus réactif} que le benzène de départ (groupe alkyle activateur). Il réagit à son tour pour donner du xylène, puis du triméthylbenzène, etc. On utilise un \textbf{excès de benzène} pour limiter ce problème.
\item \textbf{Réarrangement du carbocation :} avec les chlorures d'alkyle ramifiés (ex. \ce{CH3-CHCl-CH3}), le carbocation formé peut se réarranger, donnant un mélange d'isomères.
\end{itemize}
Au Probatoire, on retient l'équation idéale avec un chlorométhane ou chloroéthane (pas de réarrangement possible) et un excès de benzène.
\end{attention}

\courssub{Nomenclature des dérivés disubstitués : ortho, méta, para}

Lorsque deux substituants sont présents sur le noyau benzénique, trois isomères de constitution sont possibles selon leurs positions relatives. On utilise les préfixes \textbf{ortho-} (1,2-), \textbf{méta-} (1,3-) et \textbf{para-} (1,4-), abrégés \textbf{o-}, \textbf{m-}, \textbf{p-}.

\begin{definition}[Les trois isomères disubstitués]
Pour un benzène disubstitué de formule \ce{C6H4X2} (deux substituants identiques $X$) ou \ce{C6H4XY} (deux substituants différents $X$ et $Y$) :
\begin{itemize}
\item \textbf{Isomère ortho (1,2-)} : les deux substituants sont sur deux carbones \emph{voisins}.
\item \textbf{Isomère méta (1,3-)} : les deux substituants sont séparés par \emph{un carbone}.
\item \textbf{Isomère para (1,4-)} : les deux substituants sont en \emph{position opposée} sur le cycle.
\end{itemize}
\end{definition}

\begin{popfigure}
\centering
\begin{tabular}{ccc}
\chemfig{**6((-X)(-Y)----)} & \chemfig{**6((-X)(-H)(-Y)---)} & \chemfig{**6((-X)(-H)(-H)(-Y)--)} \\
\textbf{ortho} (1,2) & \textbf{méta} (1,3) & \textbf{para} (1,4)
\end{tabular}
\legende{Les trois isomères de disubstitution du benzène. Pour les xylènes (diméthylbenzènes) : ortho-xylène, méta-xylène, para-xylène.}
\end{popfigure}

\begin{aretenir}[Noms usuels à connaître]
\begin{itemize}
\item \ce{C6H4(CH3)2} : \textbf{xylènes} (ortho-, méta-, para-xylène).
\item \ce{C6H4Cl2} : \textbf{dichlorobenzènes} (ortho-, méta-, para-dichlorobenzène). Le para-dichlorobenzène est utilisé comme anti-mites.
\item \ce{C6H4(NO2)2} : \textbf{dinitrobenzènes}.
\end{itemize}
\end{aretenir}

\courssub{Bilan : pourquoi le benzène préfère-t-il la substitution ?}

\begin{aretenir}[L'essentiel sur les réactions de substitution]
\begin{itemize}
\item Dans une \textbf{substitution}, un hydrogène du noyau benzénique est remplacé par un groupe ($\ce{Cl}$, $\ce{Br}$, $\ce{NO2}$, $\ce{SO3H}$, $R$) et le \textbf{caractère aromatique est conservé}.
\item \textbf{Halogénation} : $\ce{C6H6 + Cl2 ->[FeCl3] C6H5Cl + HCl}$ (ou \ce{Br2}/\ce{FeBr3}).
\item \textbf{Nitration} : $\ce{C6H6 + HNO3 ->[\ce{H2SO4}][\SIrange{50}{60}{\celsius}] C6H5NO2 + H2O}$.
\item \textbf{Sulfonation} : $\ce{C6H6 + H2SO4 <=>[\text{chauffage}] C6H5SO3H + H2O}$ (équilibre).
\item \textbf{Friedel et Crafts} : $\ce{C6H6 + RCl ->[AlCl3] C6H5R + HCl}$ (ex. : benzène + chlorométhane $\rightarrow$ toluène). Attention à la polyalkylation.
\item Le toluène subit trois nitrations successives pour donner le \textbf{TNT} (2,4,6-trinitrotoluène).
\item Les dérivés disubstitués existent sous trois formes : \textbf{ortho} (1,2), \textbf{méta} (1,3), \textbf{para} (1,4).
\item \textbf{Interprétation} : la substitution conserve le nuage $\pi$ délocalisé, source de la stabilité exceptionnelle du benzène ; c'est pourquoi le benzène substitue bien plus facilement qu'il n'additionne.
\end{itemize}
\end{aretenir}

\begin{attention}[N'oublie jamais le catalyseur sur la flèche !]
L'erreur la plus fréquente au Probatoire est d'écrire une équation de substitution \textbf{sans préciser le catalyseur} au-dessus de la flèche. Sans \ce{FeCl3} (ou \ce{AlCl3}), le dichlore ne réagit pratiquement pas avec le benzène à froid ; sans \ce{H2SO4} concentré, la nitration est très lente ; sans \ce{AlCl3}, la réaction de Friedel et Crafts n'a pas lieu. Une équation sans ses conditions opératoires est une équation incomplète — et des points perdus.
\end{attention}

\begin{exercice}[Entraînement]
\begin{enumerate}
\item Écrire l'équation-bilan de la bromuration du benzène en précisant les conditions. Nommer le produit obtenu.
\item Écrire l'équation-bilan de la réaction de Friedel et Crafts entre le benzène et le chloroéthane \ce{C2H5Cl}. Nommer le produit.
\item Quelle masse de toluène peut-on obtenir par réaction de \SI{15,6}{g} de benzène avec un excès de chlorométhane, le rendement étant de \SI{100}{\%} ? ($M(\ce{C6H6}) = \SI{78}{g.mol^{-1}}$ ; $M(\ce{C6H5CH3}) = \SI{92}{g.mol^{-1}}$.)
\item Donner les noms et les formules semi-développées des trois isomères du dichlorobenzène \ce{C6H4Cl2}.
\end{enumerate}
\end{exercice}

\begin{corrige}[Exercice 1]
\begin{enumerate}
\item $\ce{C6H6 + Br2 ->[FeBr3] C6H5Br + HBr}$. Le produit est le \textbf{bromobenzène}.
\item $\ce{C6H6 + C2H5Cl ->[AlCl3] C6H5C2H5 + HCl}$. Le produit est l'\textbf{éthylbenzène}.
\item $n(\ce{C6H6}) = \dfrac{15,6}{78} = \SI{0,20}{mol}$. Stœchiométrie 1 pour 1 : $n(\text{toluène}) = \SI{0,20}{mol}$, d'où $m = 0,20 \times 92 = \SI{18,4}{g}$ de toluène.
\item Les trois isomères du dichlorobenzène :
\begin{itemize}
\item \textbf{1,2-dichlorobenzène (ortho-dichlorobenzène)} : \ce{Cl-C6H4-Cl} (Cl en positions 1 et 2).
\item \textbf{1,3-dichlorobenzène (méta-dichlorobenzène)} : \ce{Cl-C6H4-Cl} (Cl en positions 1 et 3).
\item \textbf{1,4-dichlorobenzène (para-dichlorobenzène)} : \ce{Cl-C6H4-Cl} (Cl en positions 1 et 4).
\end{itemize}
\end{enumerate}
\end{corrige}

\coursec{Sécurité et utilisations des composés aromatiques}

Dans la section précédente, nous avons vu que le benzène préfère la \cle{substitution} à l'addition, et que cette réactivité permet de préparer des dérivés très utiles : le nitrobenzène, le toluène, des alkylbenzènes... Or certains de ces dérivés, comme le trinitrotoluène (TNT) ou les insecticides chlorés, sont à la fois des produits industriels précieux et des substances dangereuses. Cette dualité — grande utilité, grand danger — est le cœur de cette dernière section. Un bon chimiste n'est pas seulement celui qui sait écrire une équation-bilan : c'est aussi celui qui sait \cle{manipuler sans se blesser}, \cle{protéger les autres} et \cle{gérer les déchets}. C'est un savoir-être autant qu'un savoir-faire.

\courssub{Toxicité et dangers du benzène}

\textbf{Pourquoi le benzène est-il si dangereux ?}

Le benzène \ce{C6H6} cumule trois dangers distincts qu'il faut bien distinguer :

\begin{itemize}
\item \textbf{Inflammabilité.} Le benzène est un liquide très volatil (température d'ébullition : \SI{80,1}{\celsius}) dont les vapeurs s'enflamment facilement. Son point éclair est d'environ \SI{-11}{\celsius} : cela signifie que même par une journée fraîche à Douala ou Yaoundé, l'air au-dessus d'un flacon de benzène contient assez de vapeurs pour s'enflammer au contact d'une flamme ou d'une étincelle.
\item \textbf{Toxicité aiguë.} L'inhalation de vapeurs de benzène provoque vertiges, maux de tête, somnolence ; à forte dose, elle peut entraîner la perte de conscience.
\item \textbf{Cancérogénicité.} C'est le danger le plus grave et le plus insidieux. Le benzène est classé \cle{cancérogène avéré pour l'Homme} (groupe 1 du CIRC, le Centre international de recherche sur le cancer). Une exposition répétée, même à de faibles concentrations, endommage la moelle osseuse et provoque des \cle{leucémies} (cancers du sang). Ce danger n'apparaît qu'après des années : on ne sent rien le jour de l'exposition, mais le risque s'accumule.
\end{itemize}

\begin{attention}[Le benzène est interdit de manipulation scolaire]
Parce que le benzène est cancérogène, \textbf{toute manipulation de benzène par des élèves est strictement interdite}. Au laboratoire scolaire, on le remplace par le \cle{toluène} \ce{C6H5-CH3}, beaucoup moins dangereux (non cancérogène, bien que toxique et inflammable), ou par des \textbf{simulations} (vidéos, logiciels). Ne confonds jamais « moins dangereux » avec « sans danger » : le toluène exige aussi la hotte et les gants. Autre piège : l'odeur douce et presque agréable du benzène ne trahit pas son danger — elle ne doit jamais être reniflée volontairement.
\end{attention}

\begin{savaistu}[Autrefois, on ne savait pas...]
Au XIX\textsuperscript{e} siècle et au début du XX\textsuperscript{e}, le benzène était utilisé partout : pour \textbf{dégraisser} les pièces métalliques et les mains des ouvriers, comme solvant dans les colles et les encres, et même dans certains \textbf{médicaments} et lotions capillaires ! Les cordonniers et les ouvriers d'imprimerie, exposés quotidiennement, ont payé un lourd tribut : c'est en observant des cas répétés de leucémies chez ces travailleurs que la médecine a établi, au milieu du XX\textsuperscript{e} siècle, le lien entre benzène et cancer du sang. Aujourd'hui, sa concentration dans l'essence est strictement limitée par la réglementation (moins de 1 \% en volume dans de nombreux pays). Cette histoire nous rappelle une leçon essentielle : \emph{l'absence de preuve du danger n'est pas la preuve de l'absence de danger}.
\end{savaistu}

\courssub{Les mesures de sécurité au laboratoire}

\textbf{L'intuition d'abord.} Se protéger d'un produit chimique, c'est couper les trois voies par lesquelles il peut nous atteindre : les \textbf{poumons} (inhalation), la \textbf{peau} (contact) et les \textbf{yeux} (projections). Chaque équipement de protection correspond à l'une de ces voies : la hotte protège les poumons, les gants et la blouse protègent la peau, les lunettes protègent les yeux.

\begin{methode}[Lire une fiche de sécurité avant toute manipulation]
Avant de manipuler un composé aromatique (toluène, xylène, phénol...), suis cette démarche en cinq étapes :
\begin{enumerate}
\item \textbf{Identifier les pictogrammes} sur l'étiquette du flacon : chaque losange rouge annonce un type de danger (inflammable, toxique, cancérogène...).
\item \textbf{Lire les mentions de danger} (phrases H) : par exemple « H225 : liquide et vapeurs très inflammables », « H350 : peut provoquer le cancer ».
\item \textbf{Choisir l'équipement de protection} en conséquence : blouse, lunettes de sécurité, gants adaptés (nitrile pour les solvants organiques).
\item \textbf{Préparer le poste de travail} : manipulation sous \cle{hotte ventilée} (les vapeurs sont aspirées et évacuées), aucune flamme ni appareil produisant des étincelles à proximité, verrerie en bon état.
\item \textbf{Prévoir l'après} : localiser le bidon de récupération des déchets organiques \emph{avant} de commencer, et savoir où se trouvent la douche de sécurité et le rince-œil.
\end{enumerate}
\end{methode}

\begin{popfigure}
\begin{center}
\begin{tikzpicture}[scale=1.0]
% Pictogramme 1 : inflammable (flamme)
\begin{scope}[shift={(0,0)}]
\draw[fill=white,line width=1.5pt,red!70!black,rotate=45] (-1.05,-1.05) rectangle (1.05,1.05);
\draw[fill=poporange!80!black,draw=popdark] (0,-0.55) .. controls (-0.55,-0.1) and (-0.35,0.25) .. (0,0.75) .. controls (0.35,0.25) and (0.55,-0.1) .. (0,-0.55);
\draw[fill=popgold!90!black,draw=none] (0,-0.45) .. controls (-0.3,-0.15) and (-0.18,0.1) .. (0,0.4) .. controls (0.18,0.1) and (0.3,-0.15) .. (0,-0.45);
\node[align=center,font=\small\bfseries] at (0,-2.0) {Inflammable};
\node[align=center,font=\footnotesize,text width=3.2cm] at (0,-2.75) {Liquide et vapeurs\\très inflammables};
\end{scope}
% Pictogramme 2 : toxique (tête de mort)
\begin{scope}[shift={(4.6,0)}]
\draw[fill=white,line width=1.5pt,red!70!black,rotate=45] (-1.05,-1.05) rectangle (1.05,1.05);
\draw[fill=popdark,draw=popdark] (0,0.15) circle (0.42);
\draw[fill=popdark,draw=popdark] (-0.22,-0.32) rectangle (0.22,-0.15);
\draw[fill=white,draw=none] (-0.16,0.22) circle (0.10);
\draw[fill=white,draw=none] (0.16,0.22) circle (0.10);
\draw[fill=white,draw=none] (0,0.0) circle (0.06);
\draw[line width=2pt,popdark] (-0.5,-0.55) -- (0.5,-0.85);
\draw[line width=2pt,popdark] (-0.5,-0.85) -- (0.5,-0.55);
\node[align=center,font=\small\bfseries] at (0,-2.0) {Toxique};
\node[align=center,font=\footnotesize,text width=3.2cm] at (0,-2.75) {Nocif ou mortel\\par inhalation, contact\\ou ingestion};
\end{scope}
% Pictogramme 3 : danger pour la santé (silhouette + étoile)
\begin{scope}[shift={(9.2,0)}]
\draw[fill=white,line width=1.5pt,red!70!black,rotate=45] (-1.05,-1.05) rectangle (1.05,1.05);
\draw[fill=popdark,draw=none] (0,0.42) circle (0.20);
\draw[fill=popdark,draw=none] (-0.38,-0.55) -- (-0.38,0.05) .. controls (-0.38,0.28) and (0.38,0.28) .. (0.38,0.05) -- (0.38,-0.55) -- cycle;
\draw[fill=white,draw=none] (0,0.02) -- (0.06,-0.10) -- (0.20,-0.08) -- (0.09,-0.19) -- (0.12,-0.33) -- (0,-0.25) -- (-0.12,-0.33) -- (-0.09,-0.19) -- (-0.20,-0.08) -- (-0.06,-0.10) -- cycle;
\node[align=center,font=\small\bfseries] at (0,-2.0) {Danger pour la santé};
\node[align=center,font=\footnotesize,text width=3.2cm] at (0,-2.75) {Cancérogène, mutagène,\\toxique pour la reproduction\\(cas du benzène)};
\end{scope}
\end{tikzpicture}
\end{center}
\legende{Les trois pictogrammes de danger essentiels pour les composés aromatiques. Le benzène porte les trois à la fois ; le toluène porte les deux premiers.}
\end{popfigure}

\begin{aretenir}[Les règles d'or pour manipuler un solvant aromatique]
\begin{itemize}
\item Sous \cle{hotte ventilée}, jamais sur la paillasse ouverte.
\item \cle{Gants}, \cle{lunettes} et \cle{blouse} obligatoires.
\item Aucune \cle{flamme} ni source d'étincelle à proximité (les vapeurs sont plus lourdes que l'air et s'accumulent au niveau du sol).
\item Préférer le \cle{toluène} au benzène dès que c'est possible.
\item Flacon refermé immédiatement après usage ; jamais de pipetage à la bouche.
\end{itemize}
\end{aretenir}

\courssub{Le danger particulier des dérivés nitrés}

\textbf{Pourquoi les groupes nitro rendent-ils une molécule explosive ?}

Rappelle-toi la nitration du benzène étudiée dans la section précédente : l'action du mélange sulfonitrique (acide nitrique + acide sulfurique concentré) fixe un groupe nitro \ce{-NO2} sur le noyau. Lorsque \emph{plusieurs} groupes nitro sont fixés sur le même noyau, la molécule devient un réservoir d'oxygène concentré : elle contient, dans sa propre structure, de quoi oxyder violemment ses atomes de carbone et d'hydrogène. Sa décomposition est alors brutale, très exothermique et produit un énorme volume de gaz : c'est une \cle{explosion}.

Le 2,4,6-trinitrotoluène (\cle{TNT}) s'obtient par nitration poussée du toluène :

\[ \ce{C6H5CH3 + 3 HNO3 ->[H2SO4\ conc.] C6H2(NO2)3CH3 + 3 H2O} \]

Il s'agit de \textbf{trois substitutions successives} : les trois groupes nitro se placent en positions 2, 4 et 6 par rapport au groupe méthyle \ce{-CH3}. Sa décomposition explosive s'écrit :

\[ \ce{2 C7H5N3O6 -> 3 N2 + 5 H2O + 7 CO + 7 C} \]

Vérifie l'équilibrage : 14 atomes de carbone, 10 d'hydrogène, 6 d'azote et 12 d'oxygène de chaque côté. Remarque surtout que la réaction ne consomme \textbf{aucun dioxygène extérieur} : tout l'oxydant est dans la molécule. C'est précisément ce qui distingue un explosif d'un combustible ordinaire.

L'\cle{acide picrique} (2,4,6-trinitrophénol), obtenu par nitration du phénol, est un autre explosif historique, encore plus dangereux car il forme avec les métaux des sels (picrates) extrêmement sensibles aux chocs.

\begin{exemplebox}[Le TNT : un explosif « sûr » entre des mains expertes]
Le TNT présente une propriété paradoxale et précieuse : il est \textbf{peu sensible aux chocs} et ne détone pas facilement. On peut le fondre (il fond vers \SI{80}{\celsius}) et le couler dans des obus sans qu'il explose. Il faut un \cle{détonateur} (une petite charge d'un explosif très sensible, comme l'azoture de plomb) pour déclencher sa détonation. C'est pourquoi les artificiers peuvent le transporter et le manipuler. \textbf{Mais attention} : cette stabilité ne concerne que l'explosion. Les vapeurs et les poussières de TNT restent \textbf{toxiques par contact cutané} : elles traversent la peau, provoquent des dermatites, des jaunisses et des atteintes du foie. Les ouvrières des usines de munitions de la Première Guerre mondiale, dont la peau jaunissait, étaient surnommées « canaries ». Stabilité explosive ne signifie donc pas innocuité chimique.
\end{exemplebox}

\begin{exoresolu}[Bilan de matière d'une nitration industrielle]
Une usine traite \SI{460}{g} de toluène \ce{C7H8} par nitration complète en TNT \ce{C7H5N3O6}. On donne $M(\ce{C7H8}) = \SI{92}{g.mol^{-1}}$ et $M(\ce{C7H5N3O6}) = \SI{227}{g.mol^{-1}}$. Quelle masse maximale de TNT peut-on obtenir ? Quelle masse d'acide nitrique \ce{HNO3} ($M = \SI{63}{g.mol^{-1}}$) est nécessaire ?
\tcblower
\textbf{Étape 1 : équation-bilan.}
\[ \ce{C7H8 + 3 HNO3 ->[H2SO4] C7H5N3O6 + 3 H2O} \]
\textbf{Étape 2 : quantité de toluène.}
\[ n(\ce{C7H8}) = \frac{m}{M} = \frac{460}{92} = \SI{5,0}{mol} \]
\textbf{Étape 3 : tableau d'avancement.}
\begin{center}
\begin{tabular}{|l|c|c|c|c|}
\hline
\rowcolor{popblueL} Équation & \ce{C7H8} & \ce{3 HNO3} & \ce{C7H5N3O6} & \ce{3 H2O} \\
\hline
État initial (mol) & 5,0 & $n$ & 0 & 0 \\
\hline
État final (mol) & 0 & 0 & 5,0 & 15,0 \\
\hline
\end{tabular}
\end{center}
D'après la stœchiométrie, $n(\ce{C7H5N3O6}) = n(\ce{C7H8}) = \SI{5,0}{mol}$ et $n(\ce{HNO3}) = 3 \times 5,0 = \SI{15}{mol}$.
\textbf{Étape 4 : masses.}
\[ m(\text{TNT}) = 5,0 \times 227 = \SI{1135}{g} \approx \SI{1,14}{kg} \]
\[ m(\ce{HNO3}) = 15 \times 63 = \SI{945}{g} \]
\textbf{Conclusion.} On peut obtenir au maximum environ \SI{1,14}{kg} de TNT, en consommant \SI{945}{g} d'acide nitrique. Ce calcul suppose un rendement de 100 \% ; en pratique, les nitrations successives ont des rendements inférieurs.
\end{exoresolu}

\begin{attention}[Pièges à éviter avec les composés nitrés]
\begin{itemize}
\item Ne jamais confondre \textbf{combustion} (réaction avec le \ce{O2} de l'air) et \textbf{détonation} (décomposition interne, sans oxygène extérieur).
\item Un dérivé nitré ne doit \textbf{jamais} être chauffé à flamme nue ni broyé au mortier.
\item Les gants sont indispensables : les dérivés nitrés passent \textbf{à travers la peau}.
\item Au lycée, on ne prépare \textbf{jamais} de dérivés polynitrés : la nitration étudiée en TP se limite à des démonstrations encadrées ou à des simulations.
\end{itemize}
\end{attention}

\courssub{Les utilisations des composés aromatiques}

Malgré leurs dangers, les composés aromatiques sont indispensables à l'industrie moderne. La solution n'est pas de les bannir, mais de les utiliser de manière contrôlée, dans des installations fermées et sécurisées. Voici les grandes familles d'applications :

\begin{center}
\begin{tabularx}{\linewidth}{|l|X|}
\hline
\rowcolor{popblueL} \textbf{Domaine} & \textbf{Exemples et rôle des composés aromatiques} \\
\hline
Solvants & Toluène et xylènes : dissolution des peintures, vernis, colles, encres d'imprimerie. \\
\hline
Matières plastiques & Le styrène \ce{C6H5-CH=CH2}, dérivé du benzène, se polymérise en \cle{polystyrène} (emballages, isolants). \\
\hline
Médicaments & De nombreux principes actifs contiennent un noyau benzénique : paracétamol, aspirine (acide acétylsalicylique), sulfamides. \\
\hline
Colorants & Les colorants azoïques (teintures textiles) et l'indigo sont des molécules aromatiques. \\
\hline
Insecticides & Le DDT et le lindane, dérivés chlorés du benzène, ont été massivement utilisés contre les moustiques (paludisme) avant d'être réglementés pour leur persistance dans l'environnement. \\
\hline
Explosifs & TNT, acide picrique : génie civil (mines, carrières, tunnels) et défense. \\
\hline
Carburants & Le toluène et les xylènes augmentent l'indice d'octane des essences. \\
\hline
\end{tabularx}
\end{center}

\begin{savaistu}[Les aromatiques au quotidien au Cameroun]
Le polystyrène expansé, issu du benzène, sert à fabriquer les glacières et les emballages de poisson fumé que l'on voit dans les marchés de Douala et de Yaoundé. Les colorants aromatiques teintent les pagnes wax. Et c'est la chimie aromatique qui fournit bon nombre de médicaments vendus en pharmacie, comme le paracétamol. La chimie du benzène, loin d'être abstraite, est présente dans ta vie quotidienne — d'où l'importance de comprendre aussi ses dangers.
\end{savaistu}

\courssub{Gestion des déchets au laboratoire}

La protection ne s'arrête pas à la fin de la manipulation : elle concerne aussi ce que l'on fait des restes. Un solvant aromatique versé à l'évier contamine le réseau d'eau, pollue les nappes phréatiques et peut dégager des vapeurs inflammables dans les canalisations.

\begin{aretenir}[Règles de gestion des déchets organiques]
\begin{itemize}
\item \textbf{Jamais} de benzène, toluène ou dérivé nitré versé à l'évier.
\item Les déchets de solvants organiques sont recueillis dans un \cle{bidon dédié}, étiqueté (« déchets organiques halogénés » ou « non halogénés »), fermé et stocké loin de toute flamme.
\item Les bidons sont ensuite traités par une filière spécialisée (incinération contrôlée ou régénération des solvants).
\item Toute verrerie ayant contenu un dérivé nitré est rincée avec précaution, et l'eau de rinçage rejoint le bidon de déchets, pas l'évier.
\end{itemize}
\end{aretenir}

\begin{exercice}[Pour vérifier que tu as compris]
\begin{enumerate}
\item Cite les trois dangers principaux du benzène et précise lequel justifie son interdiction en travaux pratiques scolaires.
\item Pourquoi remplace-t-on le benzène par le toluène au laboratoire ? Le toluène est-il sans danger ?
\item Écris l'équation-bilan de la décomposition explosive du TNT et explique pourquoi elle ne nécessite pas de dioxygène.
\item Un flacon porte les pictogrammes « flamme » et « danger pour la santé ». Quelles précautions de manipulation imposent-ils ?
\end{enumerate}
\end{exercice}

\begin{corrige}[Exercice n°1]
\begin{enumerate}
\item Inflammabilité, toxicité aiguë (inhalation) et cancérogénicité (leucémies). C'est la cancérogénicité, danger à long terme même à faible dose, qui justifie l'interdiction scolaire.
\item Le toluène n'est pas cancérogène et est moins toxique ; il donne des réactions de substitution analogues. Mais il reste inflammable et toxique : hotte, gants et lunettes restent obligatoires.
\item \ce{2 C7H5N3O6 -> 3 N2 + 5 H2O + 7 CO + 7 C}. La molécule contient ses propres atomes d'oxygène (six par molécule) : l'oxydant est interne, aucun \ce{O2} extérieur n'est requis.
\item « Flamme » impose l'absence de flamme et d'étincelle ; « danger pour la santé » impose hotte ventilée, gants, lunettes, blouse et limitation stricte de la durée d'exposition.
\end{enumerate}
\end{corrige}

En conclusion de cette leçon, retiens l'idée force : la structure électronique exceptionnelle du benzène — ses six électrons $\pi$ délocalisés — explique à la fois sa stabilité, sa préférence pour la substitution, et la richesse de ses dérivés. Ces dérivés sont des trésors industriels (plastiques, médicaments, colorants) mais aussi des substances à respecter (cancérogénicité, toxicité, explosivité). Le chimiste responsable est celui qui maîtrise les deux visages de cette chimie.

\newpage

\coursec{Travaux pratiques}

\courssub{Objectifs du TP}

À la fin de cette séance de travaux pratiques, tu seras capable de :

\begin{itemize}
\item comparer expérimentalement la réactivité d'un alcène (le cyclohexène) et d'un composé aromatique (le toluène) vis-à-vis du dibrome ;
\item montrer que le noyau benzénique, malgré ses trois « doubles liaisons » apparentes, \textbf{ne décolore pas} l'eau de brome dans les conditions ordinaires, contrairement aux alcènes ;
\item interpréter cette différence de comportement par la \cle{délocalisation} des six électrons $\pi$ du noyau benzénique, qui lui confère une stabilité particulière (caractère aromatique) ;
\item distinguer les conditions de l'\cle{addition} (alcènes, à l'obscurité, sans catalyseur) de celles de la \cle{substitution} électrophile sur le noyau (benzène et dérivés, en présence d'un catalyseur) ou de la substitution radicalaire sur la chaîne latérale du toluène (sous lumière vive) ;
\item appliquer les règles de sécurité liées à la manipulation des composés halogénés et aromatiques.
\end{itemize}

\begin{aretenir}[Problématique]
Le benzène \ce{C6H6} possède le même rapport hydrogène/carbone qu'un composé très insaturé. Pourtant, il ne se comporte pas comme un alcène. Comment le vérifier expérimentalement, et que nous apprend ce test sur la structure réelle du noyau benzénique ?
\end{aretenir}

\courssub{Matériel et produits}

\begin{center}
\rowcolors{1}{popblueL}{white}
\begin{tabularx}{\linewidth}{|l|X|}
\hline
\rowcolor{popblue!40} \textbf{Matériel} & \textbf{Produits} \\
\hline
2 tubes à essais propres et secs avec bouchons & Dibrome \ce{Br2} en solution dans le cyclohexane (environ \SI{1}{\%} en volume), ou eau de brome \\
\hline
2 pipettes Pasteur ou compte-gouttes & Cyclohexène \ce{C6H10} (à défaut : huile végétale riche en insaturations, par exemple huile de palme rouge décolorée ou huile d'arachide) \\
\hline
1 éprouvette graduée de \SI{10}{mL} & Toluène \ce{C6H5-CH3} (remplaçant du benzène, interdit en milieu scolaire car cancérigène) \\
\hline
1 support à tubes & Cyclohexane (solvant, pour tube témoin) \\
\hline
Gants en nitrile, lunettes de protection, blouse & Lampe puissante ou lumière solaire directe \\
\hline
Bidon de récupération des déchets halogénés & Papier pH (facultatif, pour détecter \ce{HBr}) \\
\hline
\end{tabularx}
\end{center}

\begin{savaistu}[Et si un produit manque ?]
Dans beaucoup de lycées camerounais, le cyclohexène n'est pas disponible. On peut le remplacer avantageusement par quelques gouttes d'\textbf{huile végétale} (huile de palme, huile d'arachide du marché) : les triglycérides insaturés qu'elle contiennent décolorent l'eau de brome, ce qui illustre parfaitement le test d'insaturation. Si le dibrome pur manque, l'eau de brome préparée par le laborantin convient très bien. Le toluène peut être remplacé par du xylène (diluant à peinture technique), avec les mêmes précautions.
\end{savaistu}

\courssub{Sécurité}

\begin{attention}[Pictogrammes et gestes qui sauvent]
\textbf{Produits dangereux manipulés :}
\begin{itemize}
\item \textbf{Dibrome} : pictogrammes \emph{corrosif} (GHS05) et \emph{dangereux pour la santé} (GHS08). Liquide roux très volatil, vapeurs âcres et toxiques, provoque de graves brûlures de la peau et des yeux. Manipulation \textbf{obligatoire sous hotte} ou en plein air, avec gants et lunettes.
\item \textbf{Toluène} : pictogrammes \emph{inflammable} (GHS02), \emph{dangereux pour la santé} (GHS08). Vapeurs nocives ; travailler loin de toute flamme.
\item \textbf{Cyclohexène et cyclohexane} : pictogramme \emph{inflammable} (GHS02).
\end{itemize}
\textbf{Gestes à adopter :}
\begin{enumerate}
\item Porter blouse, lunettes et gants pendant toute la séance ; cheveux longs attachés.
\item Prélever de \textbf{très petits volumes} (quelques gouttes, jamais plus de \SI{2}{mL}).
\item Ne jamais sentir directement un tube : éventer avec la main si nécessaire.
\item En cas de contact cutané avec le brome : rincer abondamment à l'eau, puis avec une solution de thiosulfate de sodium si disponible, et prévenir l'enseignant.
\item \textbf{Ne rien jeter à l'évier} : verser tous les résidus dans le bidon de récupération des déchets halogénés.
\item Le \textbf{benzène est formellement interdit} en travaux pratiques scolaires (cancérigène avéré, CMR 1A) : on le remplace par le toluène.
\end{enumerate}
\end{attention}

\courssub{Schéma du dispositif}

\begin{popfigure}
\begin{center}
\begin{tikzpicture}[scale=0.9, every node/.style={font=\small}]
% Tube 1 : cyclohexène
\draw[thick,popdark] (0,0) -- (0,3.2) (1.4,0) -- (1.4,3.2);
\draw[thick,popdark] (0,0) arc (180:360:0.7);
\fill[popgoldL] (0.02,0.05) -- (0.02,1.1) -- (1.38,1.1) -- (1.38,0.05) arc (180:360:0.68);
\draw[popgold,thick] (0.02,1.1) -- (1.38,1.1);
\node[popdark] at (0.7,-0.7) {Tube 1 : cyclohexène};
\node[popdark] at (0.7,-1.15) {+ solution de \ce{Br2}};
\node[poporange,font=\small\bfseries] at (0.7,1.6) {Décoloration immédiate};
% Tube 2 : toluène
\begin{scope}[xshift=4.5cm]
\draw[thick,popdark] (0,0) -- (0,3.2) (1.4,0) -- (1.4,3.2);
\draw[thick,popdark] (0,0) arc (180:360:0.7);
\fill[poporange!70] (0.02,0.05) -- (0.02,1.1) -- (1.38,1.1) -- (1.38,0.05) arc (180:360:0.68);
\draw[poporange,thick] (0.02,1.1) -- (1.38,1.1);
\node[popdark] at (0.7,-0.7) {Tube 2 : toluène};
\node[popdark] at (0.7,-1.15) {+ solution de \ce{Br2} (obscurité)};
\node[poporange,font=\small\bfseries] at (0.7,1.6) {Coloration persistante};
\end{scope}
% Pipette ajout Br2
\draw[thick,popblue] (2.6,3.6) -- (2.75,2.4) -- (2.9,3.6);
\draw[thick,popblue] (2.6,3.6) arc (180:360:0.15);
\node[popblue] at (2.75,4.1) {Gouttes de \ce{Br2}};
% Soleil pour tube 2 (étape lumière)
\begin{scope}[xshift=7.6cm,yshift=2.6cm]
\draw[fill=popgold,draw=popgold] (0,0) circle (0.28);
\foreach \a in {0,45,...,315} \draw[popgold,thick] (\a:0.36) -- (\a:0.6);
\node[popgold!60!black] at (0,-1.0) {Lumière vive};
\end{scope}
\draw[-{Stealth[length=3mm]},popgold,thick] (7.3,2.3) -- (6.2,1.6);
% Tube témoin (optionnel, esquissé petit à droite)
\begin{scope}[xshift=9.5cm, scale=0.6]
\draw[thick,popdark] (0,0) -- (0,3.2) (1.4,0) -- (1.4,3.2);
\draw[thick,popdark] (0,0) arc (180:360:0.7);
\fill[poporange!70] (0.02,0.05) -- (0.02,1.1) -- (1.38,1.1) -- (1.38,0.05) arc (180:360:0.68);
\draw[poporange,thick] (0.02,1.1) -- (1.38,1.1);
\node[popdark] at (0.7,-0.7) {Témoin};
\node[popdark] at (0.7,-1.1) {cyclohexane};
\end{scope}
\end{tikzpicture}
\end{center}
\legende{Dispositif du test comparatif : le tube 1 (alcène) se décolore immédiatement ; le tube 2 (aromatique) conserve la teinte orangée du brome à l'obscurité. Le tube témoin (cyclohexane) ne réagit pas non plus à l'obscurité.}
\end{popfigure}

\courssub{Protocole}

\begin{methode}[Déroulement de l'expérience]
\begin{enumerate}
\item Sous la hotte (ou en plein air), introduire dans le \textbf{tube 1} environ \SI{1}{mL} de cyclohexène (soit 15 à 20 gouttes), et dans le \textbf{tube 2} environ \SI{1}{mL} de toluène. Préparer un \textbf{tube témoin} avec \SI{1}{mL} de cyclohexane.
\item Ajouter dans chaque tube, goutte à goutte et en comptant, la solution de dibrome (environ 5 gouttes, soit \SI{0,25}{mL}) jusqu'à l'apparition d'une teinte orangée nette.
\item Boucher les tubes et agiter doucement pendant 30 secondes. Observer immédiatement et noter l'évolution des teintes.
\item Conserver les trois tubes \textbf{à l'obscurité} (placard ou boîte en carton) pendant 5 minutes, puis observer à nouveau.
\item Exposer ensuite le \textbf{tube 2} (toluène) et le \textbf{tube témoin} à la lumière vive (soleil direct ou lampe puissante) pendant 10 à 15 minutes. Observer et tester éventuellement les vapeurs du tube avec un papier pH humidifié.
\item Verser le contenu des trois tubes dans le bidon de déchets halogénés, rincer la verrerie, se laver les mains.
\end{enumerate}
\end{methode}

\courssub{Observations et résultats attendus}

\begin{center}
\rowcolors{1}{popgreenL}{white}
\begin{tabularx}{\linewidth}{|l|X|X|X|}
\hline
\rowcolor{popgreen!40} \textbf{Tube} & \textbf{Aussitôt après agitation} & \textbf{Après 5 min à l'obscurité} & \textbf{Après 10--15 min à la lumière} \\
\hline
1 : cyclohexène + \ce{Br2} & Décoloration immédiate (2 à 3 gouttes suffisent) & Solution incolore & Solution incolore, inchangée \\
\hline
2 : toluène + \ce{Br2} & Teinte orangée persistante & Teinte orangée persistante & Décoloration lente ; vapeurs âcres ; papier pH humidifié vire au rouge (dégagement de \ce{HBr}) \\
\hline
Témoin : cyclohexane + \ce{Br2} & Teinte orangée persistante & Teinte orangée persistante & Décoloration très lente ; papier pH vire au rouge (formation de \ce{HBr}) \\
\hline
\end{tabularx}
\end{center}

La décoloration quasi instantanée du tube 1 traduit une réaction \textbf{rapide} entre le dibrome et le cyclohexène. La persistance de la teinte dans le tube 2, à l'obscurité, montre qu'\textbf{aucune réaction notable} n'a lieu entre le dibrome et le noyau aromatique dans ces conditions. Le tube témoin confirme que le cyclohexane (alcane cyclique saturé) ne réagit pas non plus à l'obscurité.

\courssub{Exploitation}

\begin{exoresolu}[Questions guidées et réponses]
\textbf{Énoncé.}
\textbf{1.} Écrire l'équation-bilan de la réaction qui a lieu dans le tube 1. À quel type de réaction appartient-elle ? Préciser le nom du produit formé.
\textbf{2.} Pourquoi la décoloration du tube 1 est-elle un test caractéristique des insaturations ? Donner une application industrielle (savonnerie).
\textbf{3.} Le toluène possède, comme le benzène, trois doubles liaisons « apparentes » sur son cycle. Pourquoi ne décolore-t-il pas le brome à l'obscurité ?
\textbf{4.} Que se passe-t-il dans le tube 2 exposé à la lumière vive ? Quel indice expérimental prouve qu'il s'agit d'une substitution et non d'une addition ? Écrire l'équation-bilan et nommer le produit organique.
\textbf{5.} Quelle condition supplémentaire faudrait-il réunir pour que le toluène réagisse avec le brome \emph{sur le noyau} benzénique ? Écrire l'équation-bilan et citer les isomères formés.
\tcblower
\textbf{Solution détaillée.}
\textbf{1.} Dans le tube 1, le dibrome s'\textbf{additionne} sur la double liaison \ce{C=C} du cyclohexène :
\[ \ce{C6H10 + Br2 -> C6H10Br2} \]
Il se forme du \textbf{1,2-dibromocyclohexane}, incolore. C'est une \cle{réaction d'addition} électrophile, caractéristique des alcènes.

\textbf{2.} Un composé insaturé (alcène, alcyne, corps gras insaturé) consomme le dibrome par addition : la teinte orangée disparaît. Ce test à l'eau de brome (ou solution de dibrome) permet donc de \textbf{mettre en évidence une insaturation}. C'est d'ailleurs le principe de la mesure de l'\emph{indice d'iode} des huiles en savonnerie (on utilise \ce{I2} / \ce{ICl} mais le principe est le même : addition sur les doubles liaisons des acides gras).

\textbf{3.} Les trois doubles liaisons du cycle du toluène ne sont \textbf{pas de vraies doubles liaisons localisées} : les six électrons $\pi$ sont \cle{délocalisés} sur l'ensemble du cycle, formant un nuage électronique au-dessus et au-dessous du plan du noyau. Cette délocalisation confère au cycle une stabilité exceptionnelle (énergie de résonance $\approx$ \SI{150}{kJ.mol^{-1}}). Une addition de \ce{Br2} détruirait ce système délocalisé aromatique : elle est donc thermodynamiquement défavorable et ne se produit pas dans les conditions ordinaires. Le noyau aromatique préfère conserver son intégrité.

\textbf{4.} À la lumière vive, le tube 2 se décolore lentement et le papier pH humidifié rougit : il se dégage du bromure d'hydrogène \ce{HBr} (gaz acide). La formation de \ce{HBr} prouve qu'il s'agit d'une \cle{substitution} (un atome d'H est remplacé par Br, avec dégagement de \ce{HBr}) et non d'une addition (qui ne dégagerait rien). Sous lumière, le brome attaque la \textbf{chaîne latérale} méthyle du toluène par mécanisme \textbf{radicalaire} :
\[ \ce{C6H5-CH3 + Br2 ->[$h\nu$] C6H5-CH2Br + HBr} \]
Le produit formé est le \textbf{bromure de benzyle} (ou \emph{$\alpha$-bromotoluène}).

\textbf{5.} Pour que la substitution ait lieu \textbf{sur le noyau} benzénique (substitution électrophile aromatique), il faut un \textbf{catalyseur} de Lewis, le tribromure de fer \ce{FeBr3} (généré in situ par ajout de limaille de fer \ce{Fe} dans le mélange \ce{Br2}/toluène) :
\[ \ce{C6H5-CH3 + Br2 ->[FeBr3] C6H4Br-CH3 + HBr} \]
On obtient alors un mélange des deux isomères régiochimiques majoritaires : le \textbf{1-bromo-2-méthylbenzène} (\emph{ortho}-bromotoluène) et le \textbf{1-bromo-4-méthylbenzène} (\emph{para}-bromotoluène). Le groupe méthyle est \emph{ortho/para}-directeur. Sans catalyseur et à l'obscurité, le noyau reste inerte : c'est toute la différence entre le caractère aromatique et le caractère éthylénique.
\end{exoresolu}

\courssub{Conclusion}

Ce TP met en évidence de manière simple et spectaculaire le paradoxe du benzène : \textbf{riche en électrons $\pi$ mais pauvre en réactions d'addition}. Là où le cyclohexène décolore instantanément le brome par addition, le toluène — et à plus forte raison le benzène — conserve la teinte orangée à l'obscurité et sans catalyseur. Cette inertie apparente s'explique par la \cle{délocalisation} des six électrons $\pi$ sur tout le cycle, qui confère au noyau benzénique une stabilité particulière, appelée \cle{caractère aromatique}.

La réactivité normale du benzène et de ses dérivés n'est donc pas l'addition (qui détruirait le système aromatique) mais la \cle{substitution électrophile aromatique}, qui le préserve : halogénation en présence de \ce{FeBr3} (ou \ce{FeCl3}/\ce{Cl2}), nitration (mélange acide \ce{HNO3}/\ce{H2SO4}), sulfonation (\ce{H2SO4} concentré), réactions de Friedel et Crafts (alkylation, acylation). L'addition reste possible, mais dans des conditions énergiques (hydrogénation catalytique \ce{H2}/\ce{Ni} vers le cyclohexane, photochloration \ce{Cl2}/\ce{h\nu} vers l'hexachlorocyclohexane, lindane).

Retiens la phrase-clé de cette leçon : \emph{le benzène n'est pas un alcène à trois doubles liaisons, c'est un système aromatique délocalisé}. Cette compréhension est la clé pour prévoir et maîtriser les synthèses industrielles comme celle du TNT (trinitrotoluène) ou d'insecticides organophosphorés.

\newpage

\coursec{Méthodes et exercices résolus}

\coursec{Les composés aromatiques}

\courssub{Méthode 1 : Représenter le benzène et interpréter ses propriétés par sa structure}

\begin{methode}[Expliquer la stabilité du benzène]
Pour justifier le comportement particulier du benzène (\ce{C6H6}), suis toujours cette démarche :
\begin{enumerate}
\item \textbf{Rappeler la géométrie} : molécule \cle{plane}, hexagone régulier, 6 atomes de carbone hybridés $sp^2$, angles de \SI{120}{\degree}, 6 liaisons C--C \emph{identiques} de longueur intermédiaire (\SI{0,139}{nm}) entre une simple liaison (\SI{0,154}{nm}) et une double liaison (\SI{0,134}{nm}).
\item \textbf{Évoquer la délocalisation} : les 6 électrons $\pi$ sont \cle{délocalisés} sur l'ensemble du cycle (on parle de « sextet aromatique »), d'où la représentation symbolique par un hexagone contenant un cercle.
\item \textbf{Conclure sur la réactivité} : la délocalisation confère au noyau benzénique une grande stabilité. Le benzène \emph{ne décolore pas} l'eau de brome ni le permanganate de potassium à froid (contrairement aux alcènes) : il réagit difficilement par \cle{addition} et préfère la \cle{substitution}, qui conserve le noyau aromatique.
\end{enumerate}
\textbf{Réflexe} : dès qu'on te demande « pourquoi le benzène ne réagit-il pas comme un alcène ? », la réponse tient en un mot : \emph{délocalisation des électrons $\pi$}.
\end{methode}

\begin{exoresolu}[Le benzène est-il un alcène cyclique ?]
On réalise deux expériences à température ambiante : on verse quelques gouttes de dibrome dans du cyclohexène (expérience 1) puis dans du benzène (expérience 2). Dans le premier cas, la coloration rouge-orangé du dibrome disparaît immédiatement ; dans le second, elle persiste.
\begin{enumerate}
\item Écrire l'équation-bilan de la réaction observée dans l'expérience 1.
\item Expliquer pourquoi le benzène ne réagit pas de la même manière, alors que sa formule \ce{C6H6} suggère une forte insaturation.
\end{enumerate}
\tcblower
\textbf{1. Réaction avec le cyclohexène.} Le cyclohexène possède une véritable double liaison C=C qui réagit par \cle{addition} électrophile avec le dibrome :
\[ \ce{C6H10 + Br2 -> C6H10Br2} \]
Le produit obtenu est le 1,2-dibromocyclohexane, incolore : la décoloration observée est le test caractéristique des insaturations \emph{localisées}.

\textbf{2. Comportement du benzène.} Si le benzène possédait trois doubles liaisons classiques (formule de Kekulé), il devrait décolorer le dibrome trois fois plus vite. Or ce n'est pas le cas. L'explication repose sur la structure réelle de la molécule :
\begin{itemize}
\item les 6 électrons $\pi$ sont \textbf{délocalisés} sur tout le cycle (sextet aromatique) ;
\item les 6 liaisons C--C sont toutes identiques, de longueur \SI{0,139}{nm}, intermédiaire entre simple et double liaison ;
\item cette délocalisation stabilise fortement la molécule : rompre le système aromatique par une addition coûte beaucoup d'énergie.
\end{itemize}
\textbf{Conclusion} : le benzène n'est pas un alcène cyclique ; c'est un composé \cle{aromatique} dont les électrons délocalisés lui confèrent une stabilité exceptionnelle. Il ne réagit pas par addition dans les conditions ordinaires, mais par \emph{substitution}, qui préserve le noyau aromatique.
\end{exoresolu}

\courssub{Méthode 2 : Nommer les dérivés du benzène, en particulier les dérivés disubstitués}

\begin{methode}[Nomenclature des dérivés benzéniques]
\begin{enumerate}
\item \textbf{Monosubstitués} : on nomme le substituant suivi de « benzène » : \ce{C6H5-CH3} = méthylbenzène (nom courant : \cle{toluène}), \ce{C6H5-OH} = hydroxybenzène (\cle{phénol}), \ce{C6H5-NO2} = nitrobenzène.
\item \textbf{Disubstitués} : deux méthodes équivalentes.
\begin{itemize}
\item \emph{Numérotation} : on numérote le cycle pour obtenir les plus petits indices : 1,2- ; 1,3- ; 1,4-diméthylbenzène.
\item \emph{Préfixes traditionnels} : \cle{ortho} ($o$) = positions 1,2 ; \cle{méta} ($m$) = positions 1,3 ; \cle{para} ($p$) = positions 1,4.
\end{itemize}
\item \textbf{Cas particuliers à connaître} : les trois diméthylbenzènes s'appellent les \cle{xylènes} (ortho-, méta-, paraxylène). Quand un groupe caractéristique est prioritaire (OH, COOH, CHO), il porte l'indice 1 : 2-nitrophénol, 4-méthylphénol...
\end{enumerate}
\textbf{Réflexe} : pour un dérivé disubstitué, il n'existe que \textbf{trois} isomères de position : $o$, $m$, $p$. Apprends à les dessiner et à les nommer dans les deux systèmes.
\end{methode}

\begin{aretenir}[Composés aromatiques usuels à connaître]
\begin{center}
\begin{tabularx}{\linewidth}{|l|l|l|X|}
\hline
\rowcolor{popblueL} \textbf{Nom systématique} & \textbf{Nom usuel} & \textbf{Formule brute} & \textbf{Utilisation / Contexte} \\
\hline
Méthylbenzène & Toluène & \ce{C7H8} & Solvant, synthèse TNT, carburants \\
\hline
Hydroxybenzène & Phénol & \ce{C6H6O} & Désinfectant, résines, bakélite \\
\hline
Nitrobenzène & Nitrobenzène & \ce{C6H5NO2} & Synthèse aniline, parfumerie \\
\hline
Aminobenzène & Aniline & \ce{C6H7N} & Teintures, polyuréthanes \\
\hline
Acide benzoïque & Acide benzoïque & \ce{C7H6O2} & Conservateur (E210), médicaments \\
\hline
1,2/1,3/1,4-Diméthylbenzène & Xylènes ($o,m,p$) & \ce{C8H10} & Solvants industriels, essence \\
\hline
Naphtalène & Naphtalène & \ce{C10H8} & Antimite, synthèse colorants \\
\hline
Acide acétylsalicylique & Aspirine & \ce{C9H8O4} & Antalgique, antiagrégant plaquettaire \\
\hline
Paracétamol & Paracétamol & \ce{C8H9NO2} & Antalgique, antipyrétique \\
\hline
\end{tabularx}
\end{center}
\textbf{Note} : le phénol (extrait de l'écorce de saule, d'où « acide salicylique ») et l'aspirine illustrent le lien entre pharmacopée traditionnelle et chimie moderne.
\end{aretenir}

\begin{exoresolu}[Les trois xylènes]
La nitration du méthylbenzène (toluène) peut conduire, selon les conditions, à trois dérivés dinitrés isomères de formule brute \ce{C6H3(CH3)(NO2)2}... En fait, on s'intéresse ici aux trois diméthylbenzènes \ce{C6H4(CH3)2}, appelés xylènes, utilisés comme solvants industriels.
\begin{enumerate}
\item Écrire la formule semi-développée (représentation du cycle) de chacun des trois isomères.
\item Donner à chacun son nom systématique et son nom usuel (avec les préfixes ortho, méta, para).
\end{enumerate}
\tcblower
\textbf{1. et 2.} La formule brute commune est \ce{C8H10}. Les deux groupes méthyle peuvent occuper trois positions relatives sur le noyau :

\begin{center}
\begin{tabularx}{\linewidth}{|c|X|X|}
\hline
\rowcolor{popblueL} \textbf{Positions} & \textbf{Nom systématique} & \textbf{Nom usuel} \\
\hline
1,2 & 1,2-diméthylbenzène & \emph{ortho}-xylène ($o$-xylène) \\
\hline
1,3 & 1,3-diméthylbenzène & \emph{méta}-xylène ($m$-xylène) \\
\hline
1,4 & 1,4-diméthylbenzène & \emph{para}-xylène ($p$-xylène) \\
\hline
\end{tabularx}
\end{center}

Représentations (les sommets non marqués portent des atomes d'hydrogène) :

\begin{popfigure}
\centering
\begin{tikzpicture}
\node at (0,0) {\chemfig{**6((-CH_3)-(-CH_3)----)}};
\node at (0,-2.6) {\small 1,2-diméthylbenzène (\emph{ortho})};
\node at (5.5,0) {\chemfig{**6((-CH_3)--(-CH_3)---)}};
\node at (5.5,-2.6) {\small 1,3-diméthylbenzène (\emph{méta})};
\node at (11,0) {\chemfig{**6((-CH_3)---(-CH_3)--)}};
\node at (11,-2.6) {\small 1,4-diméthylbenzène (\emph{para})};
\end{tikzpicture}
\legende{Les trois isomères disubstitués du benzène : les xylènes.}
\end{popfigure}

\textbf{Conclusion} : il existe exactement trois dérivés diméthylés du benzène, correspondant aux positions ortho (1,2), méta (1,3) et para (1,4). Ce sont des isomères de position, de même formule brute \ce{C8H10}, mais de propriétés physiques différentes (points de fusion et d'ébullition distincts).
\end{exoresolu}

\courssub{Méthode 3 : Écrire et exploiter l'équation de combustion du benzène}

\begin{methode}[Combustion complète d'un composé aromatique]
\begin{enumerate}
\item Écrire les réactifs : hydrocarbure $+$ dioxygène \ce{O2} ; produits : \ce{CO2} $+$ \ce{H2O} (combustion \emph{complète}).
\item Équilibrer dans l'ordre : C, puis H, puis O. Pour le benzène :
\[ \ce{2 C6H6 + 15 O2 -> 12 CO2 + 6 H2O} \]
\item Pour les calculs stœchiométriques, dresser un \cle{tableau d'avancement} et utiliser $n = \dfrac{m}{M}$ ou $n = \dfrac{V}{V_m}$ pour les gaz.
\item \textbf{Signal particulier} : la combustion du benzène est souvent \emph{incomplète} (flammes très fuligineuses, dépôt de noir de carbone) car le rapport C/H est élevé. Le mentionner si l'énoncé évoque une fumée noire.
\end{enumerate}
\end{methode}

\begin{exoresolu}[Combustion du benzène]
On brûle complètement $m = \SI{7,8}{g}$ de benzène liquide dans le dioxygène de l'air.
\begin{enumerate}
\item Écrire l'équation-bilan de la combustion complète.
\item Calculer le volume de dioxygène nécessaire, mesuré dans les conditions normales ($V_m = \SI{22,4}{L.mol^{-1}}$).
\item Calculer la masse d'eau formée.
\end{enumerate}
Données : $M(\text{C}) = \SI{12}{g.mol^{-1}}$, $M(\text{H}) = \SI{1}{g.mol^{-1}}$, $M(\text{O}) = \SI{16}{g.mol^{-1}}$.
\tcblower
\textbf{1. Équation-bilan.}
\[ \ce{2 C6H6 + 15 O2 -> 12 CO2 + 6 H2O} \]
Vérification : 12 C, 12 H et 30 O de chaque côté.

\textbf{2. Volume de dioxygène.} Masse molaire du benzène :
\[ M(\ce{C6H6}) = 6 \times 12 + 6 \times 1 = \SI{78}{g.mol^{-1}} \]
Quantité de benzène brûlé :
\[ n(\ce{C6H6}) = \frac{m}{M} = \frac{7,8}{78} = \SI{0,10}{mol} \]
Tableau d'avancement (en mol) :
\begin{center}
\begin{tabularx}{\linewidth}{|l|X|X|X|X|}
\hline
\rowcolor{popblueL} & \ce{2C6H6} & \ce{15O2} & \ce{12CO2} & \ce{6H2O} \\
\hline
État initial & 0,10 & $n$ & 0 & 0 \\
\hline
État final & $0{,}10 - 2x_{max}$ & $n - 15x_{max}$ & $12x_{max}$ & $6x_{max}$ \\
\hline
\end{tabularx}
\end{center}
Le benzène est entièrement consommé : $2x_{max} = 0{,}10$ donc $x_{max} = \SI{0,050}{mol}$.
D'où :
\[ n(\ce{O2}) = 15\,x_{max} = 15 \times 0{,}050 = \SI{0,75}{mol} \]
\[ V(\ce{O2}) = n \times V_m = 0{,}75 \times 22{,}4 = \SI{16,8}{L} \]

\textbf{3. Masse d'eau formée.}
\[ n(\ce{H2O}) = 6\,x_{max} = 6 \times 0{,}050 = \SI{0,30}{mol} \]
\[ m(\ce{H2O}) = n \times M = 0{,}30 \times 18 = \SI{5,4}{g} \]

\textbf{Conclusion} : la combustion complète de \SI{7,8}{g} de benzène consomme \SI{16,8}{L} de dioxygène (CNTP) et produit \SI{5,4}{g} d'eau (et \SI{26,4}{g} de \ce{CO2}).
\end{exoresolu}

\courssub{Méthode 4 : Écrire les réactions d'addition sur le benzène}

\begin{methode}[Additions de \ce{H2} et \ce{Cl2} sur le noyau benzénique]
\begin{enumerate}
\item \textbf{Identifier les conditions} : l'addition sur le benzène n'est possible que dans des conditions énergiques qui rompent l'aromaticité :
\begin{itemize}
\item \ce{H2} : catalyseur (nickel, platine), chauffage, pression ;
\item \ce{Cl2} : lumière ultraviolette (UV).
\end{itemize}
\item \textbf{Compter les molécules ajoutées} : le noyau comporte l'équivalent de \emph{trois} doubles liaisons, donc \textbf{3 molécules} s'additionnent par noyau :
\[ \ce{C6H6 + 3 H2 ->[Ni] C6H12} \quad \text{(cyclohexane)} \]
\[ \ce{C6H6 + 3 Cl2 ->[UV] C6H6Cl6} \quad \text{(hexachlorocyclohexane)} \]
\item \textbf{Nommer les produits} : cyclohexane (hydrogénation) ; 1,2,3,4,5,6-hexachlorocyclohexane (chloration), dont l'isomère $\gamma$ est le \cle{lindane}, insecticide.
\end{enumerate}
\textbf{Réflexe} : addition = \emph{disparition} du caractère aromatique ; le produit n'est plus un composé benzénique mais un cyclohexane substitué.
\end{methode}

\begin{exoresolu}[Hydrogénation du benzène en cyclohexane]
Une unité pétrochimique (comme la SONARA à Limbé) transforme le benzène en cyclohexane, intermédiaire de la fabrication du nylon. On traite $V = \SI{50}{mL}$ de benzène liquide (densité $d = 0{,}88$) par le dihydrogène en présence de nickel.
\begin{enumerate}
\item Écrire l'équation-bilan de la réaction.
\item Calculer la quantité de matière de benzène traitée.
\item En déduire le volume de dihydrogène nécessaire (CNTP, $V_m = \SI{22,4}{L.mol^{-1}}$) et la masse de cyclohexane obtenue si la réaction est totale.
\end{enumerate}
Données : $M(\ce{C6H6}) = \SI{78}{g.mol^{-1}}$ ; $M(\ce{C6H12}) = \SI{84}{g.mol^{-1}}$ ; $\rho_{eau} = \SI{1,0}{g.mL^{-1}}$.
\tcblower
\textbf{1. Équation-bilan.}
\[ \ce{C6H6 + 3 H2 ->[Ni, \text{chaleur}] C6H12} \]

\textbf{2. Quantité de benzène.} Masse de benzène :
\[ m = d \times \rho_{eau} \times V = 0{,}88 \times 1{,}0 \times 50 = \SI{44}{g} \]
\[ n(\ce{C6H6}) = \frac{m}{M} = \frac{44}{78} = \SI{0,56}{mol} \]

\textbf{3. Bilans stœchiométriques.} D'après l'équation, 1 mol de benzène consomme 3 mol de \ce{H2} et donne 1 mol de cyclohexane :
\[ n(\ce{H2}) = 3 \times 0{,}56 = \SI{1,69}{mol} \]
\[ V(\ce{H2}) = n \times V_m = 1{,}69 \times 22{,}4 \approx \SI{37,9}{L} \]
\[ n(\ce{C6H12}) = \SI{0,56}{mol} \quad \Rightarrow \quad m(\ce{C6H12}) = 0{,}56 \times 84 \approx \SI{47,4}{g} \]

\textbf{Conclusion} : l'hydrogénation totale de \SI{50}{mL} de benzène nécessite environ \SI{38}{L} de dihydrogène (CNTP) et fournit environ \SI{47}{g} de cyclohexane. Remarque : la masse du produit est supérieure à celle du benzène de départ car on a \emph{ajouté} 6 atomes d'hydrogène par molécule.
\end{exoresolu}

\courssub{Méthode 5 : Écrire les réactions de substitution sur le benzène}

\begin{methode}[Les quatre substitutions à connaître]
La substitution est la réaction privilégiée du benzène : un atome d'hydrogène du noyau est \cle{remplacé} par un autre atome ou groupe, et le caractère aromatique est \emph{conservé}. Pour chaque réaction, retiens : réactif, catalyseur/conditions, produit, sous-produit.
\begin{enumerate}
\item \textbf{Halogénation} : \ce{Cl2} ou \ce{Br2} en présence d'un catalyseur (\ce{FeCl3}, \ce{AlCl3} ou \ce{Fe}) :
\[ \ce{C6H6 + Cl2 ->[FeCl3] C6H5Cl + HCl} \]
\item \textbf{Nitration} : acide nitrique concentré \ce{HNO3} en présence d'acide sulfurique concentré \ce{H2SO4}, à chaud :
\[ \ce{C6H6 + HNO3 ->[H2SO4] C6H5NO2 + H2O} \]
Produit : nitrobenzène (étape vers le TNT par nitrations successives du toluène).
\item \textbf{Sulfonation} : acide sulfurique concentré à chaud (réaction réversible) :
\[ \ce{C6H6 + H2SO4 -> C6H5SO3H + H2O} \]
Produit : acide benzènesulfonique.
\item \textbf{Friedel et Crafts (alkylation)} : chlorure d'alkyle R--Cl avec \ce{AlCl3} :
\[ \ce{C6H6 + R-Cl ->[AlCl3] C6H5-R + HCl} \]
Exemple : \ce{C6H6 + CH3Cl ->[AlCl3] C6H5CH3 + HCl} (toluène).
\end{enumerate}
\textbf{Réflexe de vérification} : dans une substitution, le produit organique contient toujours \textbf{6 carbones et 5 hydrogènes} sur le noyau (\ce{C6H5-X}) et il se forme toujours un sous-produit (\ce{HCl} ou \ce{H2O}).
\end{methode}

\begin{exoresolu}[Préparation du TNT]
Le trinitrotoluène (TNT), explosif bien connu, s'obtient par nitration du toluène \ce{C6H5CH3}. La nitration complète s'écrit en remplaçant trois atomes d'hydrogène du noyau par trois groupes nitro \ce{NO2}.
\begin{enumerate}
\item Écrire l'équation-bilan de la trinitration du toluène par l'acide nitrique.
\item On traite $m = \SI{46}{g}$ de toluène. Calculer la masse de TNT obtenue si la réaction est totale.
\end{enumerate}
Données : $M(\ce{C7H8}) = \SI{92}{g.mol^{-1}}$ ; $M(\ce{C7H5N3O6}) = \SI{227}{g.mol^{-1}}$.
\tcblower
\textbf{1. Équation-bilan.} Chaque nitration remplace un H du noyau par \ce{NO2} et libère une molécule d'eau. Pour trois nitrations :
\[ \ce{C6H5CH3 + 3 HNO3 ->[H2SO4] C6H2(NO2)3CH3 + 3 H2O} \]
soit, avec les formules brutes :
\[ \ce{C7H8 + 3 HNO3 -> C7H5N3O6 + 3 H2O} \]
Vérification : C : 7 = 7 ; H : $8 + 3 = 11$ et $5 + 6 = 11$ ; N : 3 = 3 ; O : 9 = 9. L'équation est équilibrée.

\textbf{2. Masse de TNT.} Quantité de toluène :
\[ n(\ce{C7H8}) = \frac{m}{M} = \frac{46}{92} = \SI{0,50}{mol} \]
D'après l'équation, 1 mol de toluène donne 1 mol de TNT :
\[ n(\text{TNT}) = \SI{0,50}{mol} \]
\[ m(\text{TNT}) = n \times M = 0{,}50 \times 227 = \SI{113,5}{g} \approx \SI{114}{g} \]

\textbf{Conclusion} : la trinitration de \SI{46}{g} de toluène fournit théoriquement environ \SI{114}{g} de TNT. La masse augmente fortement car on a greffé trois groupes \ce{NO2} (chacun de masse \SI{46}{g.mol^{-1}}) en ne retirant que trois atomes d'hydrogène.
\end{exoresolu}

\courssub{Méthode 6 : Distinguer addition et substitution, et appliquer les règles de sécurité}

\begin{methode}[Reconnaître le type de réaction et manipuler en sécurité]
\textbf{A. Diagnostic rapide du type de réaction :}
\begin{enumerate}
\item Regarder le \textbf{nombre de produits organiques} : un seul produit organique $\Rightarrow$ \cle{addition} ; un produit organique $+$ un sous-produit (\ce{HCl}, \ce{H2O}) $\Rightarrow$ \cle{substitution}.
\item Regarder la \textbf{formule du produit} : si le noyau a gagné des atomes sans en perdre (\ce{C6H6 -> C6H12}), c'est une addition ; si un H a été remplacé (\ce{C6H6 -> C6H5NO2}), c'est une substitution.
\item Regarder les \textbf{conditions} : Ni/chaleur ou UV $\Rightarrow$ addition ; \ce{FeCl3}, \ce{AlCl3}, \ce{H2SO4} concentré $\Rightarrow$ substitution.
\end{enumerate}
\textbf{B. Sécurité (benzène et dérivés nitrés) :}
\begin{enumerate}
\item Le benzène est \textbf{inflammable}, \textbf{volatil} et \textbf{cancérigène} : manipuler sous hotte, loin de toute flamme, avec gants et lunettes ; ne jamais en respirer les vapeurs.
\item Les dérivés nitrés (nitrobenzène, TNT) sont \textbf{toxiques} et \textbf{explosifs} : éviter chocs, frottements et chauffage brutal ; stocker en petites quantités, au frais.
\item Les mélanges nitrants (\ce{HNO3} + \ce{H2SO4} concentrés) sont très corrosifs : verser l'acide dans l'eau, jamais l'inverse, et refroidir le milieu réactionnel.
\end{enumerate}
\end{methode}

\begin{exoresolu}[Addition ou substitution ?]
Pour chacune des réactions suivantes, indiquer s'il s'agit d'une addition ou d'une substitution, en justifiant, et nommer le produit organique obtenu.
\begin{enumerate}
\item $\ce{C6H6 + 3 Cl2 ->[UV] C6H6Cl6}$
\item $\ce{C6H6 + Br2 ->[FeBr3] C6H5Br + HBr}$
\item $\ce{C6H6 + C2H5Cl ->[AlCl3] C6H5C2H5 + HCl}$
\item $\ce{C6H6 + 3 H2 ->[Ni] C6H12}$
\end{enumerate}
\tcblower
\textbf{1.} $\ce{C6H6 + 3 Cl2 ->[UV] C6H6Cl6}$ : \textbf{addition}. Un seul produit, la molécule a gagné 6 atomes de chlore sans rien perdre ; le caractère aromatique a disparu. Produit : le 1,2,3,4,5,6-hexachlorocyclohexane (dont l'isomère $\gamma$ est le lindane, insecticide).

\textbf{2.} $\ce{C6H6 + Br2 ->[FeBr3] C6H5Br + HBr}$ : \textbf{substitution} (bromuration). Un atome d'hydrogène du noyau est remplacé par un atome de brome ; il se forme le sous-produit \ce{HBr} ; le noyau aromatique est conservé. Produit : le bromobenzène.

\textbf{3.} $\ce{C6H6 + C2H5Cl ->[AlCl3] C6H5C2H5 + HCl}$ : \textbf{substitution} de Friedel et Crafts (alkylation). Un hydrogène du noyau est remplacé par le groupe éthyle \ce{C2H5}, avec formation de \ce{HCl}. Produit : l'éthylbenzène.

\textbf{4.} $\ce{C6H6 + 3 H2 ->[Ni] C6H12}$ : \textbf{addition} (hydrogénation). Trois molécules de dihydrogène s'ajoutent sur le noyau ; le produit n'est plus aromatique. Produit : le cyclohexane.

\textbf{Conclusion} : les réactions 1 et 4 sont des additions (conditions énergiques, perte de l'aromaticité) ; les réactions 2 et 3 sont des substitutions (catalyseurs de Lewis, conservation du noyau benzénique, formation d'un sous-produit).
\end{exoresolu}

\begin{attention}[Les 5 erreurs qui coûtent des points au Probatoire]
\begin{enumerate}
\item \textbf{Dessiner le benzène avec trois doubles liaisons alternées fixes} et l'assimiler à un « cyclohexatriène ». C'est faux : les 6 électrons $\pi$ sont délocalisés ; on représente le noyau par un hexagone avec un cercle, et toutes les liaisons C--C sont identiques.
\item \textbf{Confondre addition et substitution.} Addition : un seul produit, le noyau disparaît (\ce{C6H6 + 3H2 -> C6H12}). Substitution : un H remplacé, un sous-produit (\ce{HCl} ou \ce{H2O}) apparaît, le noyau est conservé. Le benzène réagit \emph{préférentiellement} par substitution.
\item \textbf{Oublier les conditions opératoires} dans les équations : \ce{FeCl3} (halogénation), \ce{H2SO4} concentré (nitration), \ce{AlCl3} (Friedel et Crafts), Ni/chaleur (hydrogénation), UV (chloration par addition). Une équation sans catalyseur est une équation incomplète.
\item \textbf{Mal équilibrer la combustion} : $\ce{2 C6H6 + 15 O2 -> 12 CO2 + 6 H2O}$. Équilibre dans l'ordre C, H, O, et n'hésite pas à mettre le coefficient 2 devant \ce{C6H6} pour éviter les fractions.
\item \textbf{Négliger la sécurité et les noms usuels} : le benzène est cancérigène (hotte, gants, pas de flamme) ; les dérivés nitrés sont explosifs. Et apprends les noms : toluène, phénol, xylènes, et les positions \emph{ortho} (1,2), \emph{méta} (1,3), \emph{para} (1,4) — question classique et facile à réussir !
\end{enumerate}
\end{attention}

\newpage

\coursec{Exercices}

\courssub{Je vérifie mes connaissances}

\begin{exercice}[QCM : structure du benzène]
Pour chaque affirmation, choisir la (ou les) bonne(s) réponse(s).
\begin{enumerate}
\item La molécule de benzène \ce{C6H6} est :
\begin{enumerate}
\item plane ; \item tétraédrique ; \item formée de 6 atomes de carbone hybridés $sp^2$.
\end{enumerate}
\item Dans le noyau benzénique, les électrons $\pi$ sont :
\begin{enumerate}
\item localisés sur trois doubles liaisons fixes ;
\item délocalisés sur l'ensemble du cycle ;
\item absents.
\end{enumerate}
\item La longueur des six liaisons carbone--carbone du benzène est :
\begin{enumerate}
\item celle d'une liaison simple \ce{C-C} ;
\item celle d'une liaison double \ce{C=C} ;
\item intermédiaire entre une liaison simple et une liaison double.
\end{enumerate}
\item La réaction caractéristique du benzène est :
\begin{enumerate}
\item l'addition ; \item la substitution ; \item l'élimination.
\end{enumerate}
\end{enumerate}
\end{exercice}

\begin{exercice}[Vrai ou faux, à justifier]
Répondre par VRAI ou FAUX, puis justifier brièvement.
\begin{enumerate}
\item Le benzène décolore rapidement l'eau de brome à température ambiante.
\item Le benzène peut réaliser une réaction d'addition avec le dihydrogène en présence de nickel.
\item La nitration du benzène nécessite un mélange d'acide nitrique et d'acide sulfurique concentrés.
\item Le 1,2-diméthylbenzène est appelé \emph{méta}-xylène.
\item Les dérivés nitrés du benzène sont des composés sans danger que l'on peut manipuler sans précaution.
\end{enumerate}
\end{exercice}

\begin{exercice}[Texte à trous]
Recopier et compléter le texte suivant avec les mots ou groupes de mots : \emph{délocalisés ; substitution ; plane ; catalyseur ; noyau benzénique ; ortho ; nitration}.

La molécule de benzène est \ldots\ldots\ldots : ses douze atomes sont dans un même plan. Ses six électrons $\pi$ sont \ldots\ldots\ldots sur tout le cycle, ce qui confère au \ldots\ldots\ldots une grande stabilité. C'est pourquoi le benzène préfère les réactions de \ldots\ldots\ldots, qui conservent le cycle aromatique. La réaction du benzène avec l'acide nitrique est une \ldots\ldots\ldots ; elle nécessite un \ldots\ldots\ldots : l'acide sulfurique concentré. Dans le 1,2-dichlorobenzène, les deux atomes de chlore sont en position \ldots\ldots\ldots .
\end{exercice}

\begin{exercice}[Définitions]
Définir en une ou deux phrases chacun des termes suivants :
\begin{enumerate}
\item composé aromatique ;
\item réaction de substitution électrophile ;
\item réaction de Friedel et Crafts ;
\item positions \emph{ortho}, \emph{méta} et \emph{para} sur un noyau benzénique disubstitué.
\end{enumerate}
\end{exercice}

\courssub{Je m'entraîne}

\begin{exercice}[Les xylènes]
\begin{enumerate}
\item Écrire la formule brute des trois dérivés diméthylés du benzène (xylènes).
\item Représenter chacun d'eux en formule semi-développée (on représentera le noyau par un hexagone portant les substituants).
\item Nommer chaque isomère en nomenclature systématique et préciser, pour chacun, s'il s'agit du dérivé \emph{ortho}, \emph{méta} ou \emph{para}.
\end{enumerate}
\end{exercice}

\begin{exercice}[Équations-bilans de substitution]
Compléter et équilibrer les équations-bilans des réactions suivantes, réalisées sur le benzène. On précisera pour chacune les conditions expérimentales (catalyseur, température) et le type de réaction.
\begin{enumerate}
\item \ce{C6H6 + Cl2 ->[FeCl3] \ldots + \ldots}
\item \ce{C6H6 + HNO3 ->[H2SO4] \ldots + \ldots}
\item \ce{C6H6 + SO3 -> \ldots} (sulfonation, à l'aide d'oléum)
\item \ce{C6H6 + CH3-CH2-Cl ->[AlCl3] \ldots + \ldots} (Friedel et Crafts)
\end{enumerate}
Nommer le produit organique obtenu dans chaque cas.
\end{exercice}

\begin{exercice}[Additions sur le benzène]
\begin{enumerate}
\item Écrire l'équation-bilan de l'addition du dihydrogène sur le benzène en présence de nickel. Nommer le produit obtenu.
\item Écrire l'équation-bilan de l'addition du dichlore sur le benzène sous l'action de la lumière. Donner la formule brute du produit obtenu (hexachlorocyclohexane, ancien insecticide).
\item Pourquoi dit-on que ces réactions d'addition font \og perdre \fg{} le caractère aromatique du benzène ?
\end{enumerate}
\end{exercice}

\begin{exercice}[Benzène et cyclohexène face au brome]
On dispose de deux flacons non étiquetés contenant l'un du benzène, l'autre du cyclohexène \ce{C6H10}.
\begin{enumerate}
\item Décrire un test simple, réalisable au laboratoire du lycée, permettant de distinguer les deux liquides. Préciser les observations attendues dans chaque cas.
\item Expliquer, en s'appuyant sur la structure électronique du noyau benzénique, pourquoi le benzène décolore difficilement l'eau de brome alors que le cyclohexène la décolore instantanément.
\item Écrire l'équation-bilan de la réaction du cyclohexène avec le dibrome.
\end{enumerate}
\end{exercice}

\courssub{Je me prépare au Probatoire}

\begin{exercice}[Combustion du benzène et du toluène]
\emph{Données :} $M(\ce{C}) = \SI{12}{g.mol^{-1}}$ ; $M(\ce{H}) = \SI{1}{g.mol^{-1}}$ ; $M(\ce{O}) = \SI{16}{g.mol^{-1}}$ ; volume molaire dans les conditions de l'expérience : $V_m = \SI{24}{L.mol^{-1}}$.

Le benzène \ce{C6H6} et le toluène (méthylbenzène) \ce{C7H8} sont des solvants aromatiques utilisés dans l'industrie.
\begin{enumerate}
\item Écrire et équilibrer l'équation-bilan de la combustion complète du benzène dans le dioxygène.
\item Écrire et équilibrer l'équation-bilan de la combustion complète du toluène dans le dioxygène.
\item On réalise la combustion complète de $m = \SI{15,6}{g}$ de benzène.
\begin{enumerate}
\item Calculer la quantité de matière de benzène brûlé.
\item À l'aide d'un tableau d'avancement, déterminer la quantité de dioxygène nécessaire.
\item En déduire le volume de dioxygène correspondant, mesuré dans les conditions de l'expérience.
\item Calculer la masse de dioxyde de carbone formée.
\end{enumerate}
\item La flamme du benzène est très fuligineuse (elle dégage beaucoup de suies). Proposer une explication en comparant la teneur en carbone du benzène à celle d'un alcane comme l'hexane \ce{C6H14}.
\end{enumerate}
\end{exercice}

\begin{exercice}[Préparation du nitrobenzène]
\emph{Données :} $M(\ce{C6H6}) = \SI{78}{g.mol^{-1}}$ ; $M(\ce{C6H5NO2}) = \SI{123}{g.mol^{-1}}$.

Au laboratoire, on prépare le nitrobenzène, intermédiaire de synthèse utilisé notamment dans la fabrication de colorants, par action d'un mélange d'acide nitrique et d'acide sulfurique concentrés sur le benzène :
\[ \ce{C6H6 + HNO3 ->[H2SO4] C6H5NO2 + H2O} \]
\begin{enumerate}
\item De quel type de réaction s'agit-il ? Justifier.
\item Quel est le rôle de l'acide sulfurique dans cette réaction ?
\item On fait réagir $m = \SI{78}{g}$ de benzène, l'acide nitrique étant en excès.
\begin{enumerate}
\item Calculer la quantité de matière de benzène mise en jeu.
\item En déduire la masse théorique de nitrobenzène attendue.
\end{enumerate}
\item À l'issue de l'expérience, on recueille en réalité $m' = \SI{98,4}{g}$ de nitrobenzène. Calculer le rendement de la préparation.
\item Citer deux précautions à prendre pour manipuler le benzène et les dérivés nitrés au laboratoire.
\end{enumerate}
\end{exercice}

\begin{exercice}[Du benzène au TNT : le trinitrotoluène]
\emph{Données :} $M(\ce{C7H8}) = \SI{92}{g.mol^{-1}}$ ; $M(\ce{TNT}) = M(\ce{C7H5N3O6}) = \SI{227}{g.mol^{-1}}$.

Le 2,4,6-trinitrotoluène (TNT) est un dérivé nitré du toluène \ce{C6H5-CH3}, obtenu industriellement par nitration successive du toluène. On admet que la nitration complète d'une mole de toluène consomme trois moles d'acide nitrique et donne une mole de TNT.
\begin{enumerate}
\item Rappeler la formule brute du toluène, puis celle du TNT.
\item Écrire, en utilisant les formules brutes, l'équation-bilan globale de la trinitration du toluène (il se forme également de l'eau).
\item Représenter la formule semi-développée du 2,4,6-trinitrotoluène (on représentera le noyau par un hexagone portant le groupe méthyle en position 1).
\item Sur le noyau du TNT, préciser les positions relatives (\emph{ortho}, \emph{méta} ou \emph{para}) des trois groupes nitro par rapport au groupe méthyle.
\item Une usine traite $m = \SI{46}{kg}$ de toluène.
\begin{enumerate}
\item Calculer la quantité de matière de toluène correspondante.
\item En déduire la masse théorique de TNT obtenue.
\end{enumerate}
\item Le TNT est un explosif puissant. Rédiger un court paragraphe (5 à 8 lignes) expliquant pourquoi la manipulation des dérivés nitrés aromatiques exige des mesures de sécurité strictes, en citant au moins trois règles concrètes à respecter au laboratoire.
\end{enumerate}
\end{exercice}

\begin{aretenir}[Résumé : Réactivité du benzène]
Le benzène (\ce{C6H6}) est le prototype des \textbf{composés aromatiques} : cycle plan, 6 atomes C $sp^2$, 6 électrons $\pi$ délocalisés (règle de Hückel $4n+2, n=1$). Cette délocalisation confère une \textbf{stabilité exceptionnelle} (énergie de résonance $\approx \SI{150}{kJ.mol^{-1}}$).
\begin{itemize}
\item \textbf{Combustion :} Complète $\rightarrow \ce{CO2} + \ce{H2O}$ ; flamme très fuligineuse (excès de C).
\item \textbf{Additions (conditions drastiques) :} $\ce{+ 3 H2 ->[Ni, T^\circ, P] C6H12}$ (cyclohexane) ; $\ce{+ 3 Cl2 ->[h\nu] C6H6Cl6}$ (hexachlorocyclohexane). \textbf{Perte de l'aromaticité}.
\item \textbf{Substitutions électrophiles (SEAr, réaction caractéristique) :} Conservent le noyau aromatique. Nécessitent un \textbf{électrophile fort} généré par un \textbf{catalyseur} (acide de Lewis \ce{FeCl3, AlCl3} ou acide de Brønsted \ce{H2SO4}).
\begin{itemize}
\item Halogénation : \ce{Cl2/FeCl3} $\rightarrow$ chlorobenzène.
\item Nitration : \ce{HNO3/H2SO4} $\rightarrow$ nitrobenzène.
\item Sulfonation : \ce{SO3/oleum} $\rightarrow$ acide benzènesulfonique (réversible).
\item Friedel-Crafts : Alkylation (\ce{R-Cl/AlCl3}) ou Acylation (\ce{RCOCl/AlCl3}).
\end{itemize}
\end{itemize}
\textbf{Nomenclature disubstitués :} positions 1,2 = \emph{ortho} ; 1,3 = \emph{méta} ; 1,4 = \emph{para}.
\textbf{Sécurité :} Benzène = CMR (cancérogène, mutagène, reprotoxique). Dérivés nitrés = toxiques + explosifs. Hotte, EPI, pas de flamme, quantités minimes.
\end{aretenir}

\begin{attention}[Pièges classiques au Probatoire]
\begin{itemize}
\item \textbf{Confondre addition et substitution :} Le benzène \textbf{ne décolore pas} l'eau de brome ni \ce{KMnO4} à froid (pas d'addition facile). C'est le test de distinction avec les alcènes.
\item \textbf{Équilibrage nitration/sulfonation :} Ne pas oublier l'eau produite ($\ce{HNO3 -> NO2+ + H2O}$ ; $\ce{SO3 + H2O -> H2SO4}$).
\item \textbf{Catalyseur Friedel-Crafts :} \ce{AlCl3} doit être \textbf{anhydre} (hydrolyse violente par l'eau). Le solvant doit être anhydre.
\item \textbf{Orientation :} Le groupe \ce{-CH3} (toluène) est \emph{ortho/para}-directeur \textbf{activant}. Le groupe \ce{-NO2} (nitrobenzène) est \emph{méta}-directeur \textbf{désactivant} (nitration du nitrobenzène très lente, conditions drastiques).
\item \textbf{Calculs :} Utiliser un \textbf{tableau d'avancement}. Vérifier l'unité de la masse (kg vs g). Rendement = $\frac{m_{exp}}{m_{th}} \times 100$.
\end{itemize}
\end{attention}

\begin{savaistu}[Le benzène au Cameroun et dans le monde]
\begin{itemize}
\item \textbf{Histoire :} Découvert par Faraday (1825) dans un résidu d'éclairage au gaz. Structure élucidée par Kekulé (1865) : rêve du serpent se mordant la queue (cycle). Théorie de la résonance (Pauling, 1930s).
\item \textbf{Industrie :} Le benzène est une \textbf{plate-forme chimique majeure} (pétrochimie : reformage catalytique, vapocraquage). Précurseur de : styrène (polystyrène), phénol (résines, bakélite), cyclohexane (nylon-6,6), colorants, médicaments (aspirine), détergents (LAS).
\item \textbf{Au Cameroun :} La \textbf{SONARA} (Société Nationale de Raffinage, Limbé) produit des coupes aromatiques (BTX : Benzène, Toluène, Xylènes) utilisées localement ou exportées. Le toluène et les xylènes sont des solvants courants dans les peintures (usines de Douala/Bassa) et l'artisanat. L'essence à la pompe contient des aromatiques (indice d'octane).
\item \textbf{Santé/Environnement :} Le benzène est classé \textbf{cancérogène certain (Groupe 1 CIRC)} : leucémies (LAM). Normes strictes : VLEP (Valeur Limite d'Exposition Professionnelle) \SI{1}{ppm} (\SI{3,25}{mg.m^{-3}}) en France/UE. Interdit dans les produits grand public (colle, correcteur). Surveillance de l'air ambiant (stations de mesure à Douala, Yaoundé).
\item \textbf{Pharmacopée :} Le noyau aromatique est le squelette de nombreux principes actifs : paracétamol, aspirine (acide acétylsalicylique), ibuprofène, anesthésiques locaux (lidocaïne), antihistaminiques.
\end{itemize}
\end{savaistu}

\begin{popfigure}
\begin{tikzpicture}[scale=0.9, every node/.style={font=\small}]
% Benzene resonance hybrid
\def\R{1.8}
\def\r{0.3}
\foreach \angle/\label in {90/1, 30/2, -30/3, -90/4, -150/5, 150/6} {
  \coordinate (C\label) at (\angle:\R);
  \fill[popblue] (C\label) circle (2.5pt);
}
\draw[popblue, thick] (C1) -- (C2) -- (C3) -- (C4) -- (C5) -- (C6) -- cycle;
\draw[popblue!50, thick, dashed] (C1) -- (C4) (C2) -- (C5) (C3) -- (C6);
% Circle for delocalization
\draw[poporange, thick] (0,0) circle (\R-0.25);
\node[poporange] at (0, -2.5) {Nuage $\pi$ délocalisé (6 e$^-$)};
% Orbitales pz
\foreach \angle in {90, 30, -30, -90, -150, 150} {
  \draw[popgreen, thick, ->] (\angle:\R) -- (\angle:\R+0.6) node[midway, right, popgreen] {$p_z$};
}
\node[popdark] at (0, 3.2) {Structure de résonance du benzène};
\end{tikzpicture}
\legende{Représentation du benzène : hybridation $sp^2$ (liaisons $\sigma$ en trait plein), orbitales $p_z$ perpendiculaires (flèches vertes) et nuage $\pi$ délocalisé (cercle orange). Les 6 liaisons C--C sont équivalentes (longueur 140 pm).}
\end{popfigure}

\begin{popfigure}
\begin{tikzpicture}[scale=0.8, every node/.style={font=\small}]
% Xylenes structures with chemfig-like drawing using TikZ
\def\drawXylene#1#2#3#4#5{ % x, y, pos1, pos2, label
  \begin{scope}[shift={(#1,#2)}]
    \foreach \a/\n in {90/1, 30/2, -30/3, -90/4, -150/5, 150/6} {
      \coordinate (C\n) at (\a:1.2);
      \fill[popblue] (C\n) circle (2pt);
    }
    \draw[popblue, thick] (C1)--(C2)--(C3)--(C4)--(C5)--(C6)--cycle;
    \draw[poporange, thick] (0,0) circle (0.95);
    % Substituants
    \foreach \pos in {#3,#4} {
      \coordinate (S) at (\pos*60-30:1.2);
      \draw[popdark, thick] (S) -- ++(\pos*60-30:0.6) node[popdark, anchor=center] {$\ce{CH3}$};
    }
    \node[popdark, align=center] at (0, -2.0) {#5};
  \end{scope}
}
\drawXylene{0}{0}{1}{2}{ortho-xylène\\1,2-diméthylbenzène}
\drawXylene{5}{0}{1}{3}{méta-xylène\\1,3-diméthylbenzène}
\drawXylene{10}{0}{1}{4}{para-xylène\\1,4-diméthylbenzène}
\end{tikzpicture}
\legende{Les trois isomères des xylènes (diméthylbenzènes). Le groupe méthyle est un substituant \emph{ortho/para}-directeur activant.}
\end{popfigure}

\begin{methode}[Calcul de rendement et tableau d'avancement]
\textbf{Étapes types pour un exercice de synthèse (Probatoire) :}
\begin{enumerate}
\item \textbf{Équation-bilan} équilibrée (formules brutes ou semi-développées). Identifier le réactif limitant (souvent le substrat aromatique).
\item \textbf{Données} : masses $m$, masses molaires $M$ (calculées si non données), volume molaire $V_m$ si gaz.
\item \textbf{Quantités de matière initiales} : $n = m/M$ (ou $n = V/V_m$).
\item \textbf{Tableau d'avancement} :
\begin{center}
\begin{tabular}{|c|c|c|c|c|}\hline
 & Réactif 1 & Réactif 2 & Produit 1 & Produit 2 \\\hline
État initial & $n_1$ & $n_2$ & 0 & 0 \\\hline
État final & $n_1 - a x$ & $n_2 - b x$ & $c x$ & $d x$ \\\hline
Avancement max $x_{max}$ & \multicolumn{4}{c|}{Déterminé par le réactif limitant} \\\hline
\end{tabular}
\end{center}
\item \textbf{Quantités finales} : $n_{final} = n_{initial} \pm \nu_i x_{max}$.
\item \textbf{Masse/Volume théorique} : $m_{th} = n_{final} \times M$ ou $V_{th} = n_{final} \times V_m$.
\item \textbf{Rendement} : $\eta = \frac{m_{exp}}{m_{th}} \times 100$ (ou $\frac{n_{exp}}{n_{th}}$). Toujours $\le 100\%$.
\end{enumerate}
\cle{Mots-clés : réactif limitant, avancement, rendement, masse théorique.}
\end{methode}

\begin{experience}[Test de distinction : Benzène vs Cyclohexène (Eau de brome)]
\textbf{Objectif :} Différencier un hydrocarbure aromatique (benzène) d'un alcène cyclique (cyclohexène) par leur réactivité vis-à-vis du brome.
\textbf{Matériel :} 2 tubes à essai, compte-gouttes, eau de brome (solution de \ce{Br2} dans l'eau, \SI{\sim 0,1}{mol.L^{-1}}, colorée jaune-orangé), benzène pur, cyclohexène pur, hotte aspirante (benzène CMR, cyclohexène inflammable/irritant).
\textbf{Protocole :}
\begin{enumerate}
\item Introduire \SI{2}{mL} d'eau de brome dans chaque tube.
\item Ajouter \SI{5}{gouttes} de benzène dans le tube 1, \SI{5}{gouttes} de cyclohexène dans le tube 2.
\item Boucher, agiter vigoureusement pendant \SI{30}{s}. Observer.
\end{enumerate}
\textbf{Observations :}
\begin{itemize}
\item \textbf{Tube 1 (Benzène) :} Deux phases liquides non miscibles. La phase aqueuse (inférieure) \textbf{reste colorée} (jaune-orangé). Pas de réaction visible à froid sans catalyseur.
\item \textbf{Tube 2 (Cyclohexène) :} \textbf{Décoloration immédiate et totale} de l'eau de brome. La solution devient incolore. Formation de 1,2-dibromocyclohexane (incolore).
\end{itemize}
\textbf{Interprétation :} Le cyclohexène possède une double liaison C=C \textbf{localisée}, riche en électrons $\pi$, qui réagit instantanément par \textbf{addition électrophile} sur \ce{Br2}. Le benzène a ses électrons $\pi$ \textbf{délocalisés} sur 6 atomes ; l'attaque sur le noyau briserait l'aromaticité (énergie d'activation trop élevée sans catalyseur \ce{FeBr3} ou lumière UV).
\textbf{Sécurité :} Manipulation sous hotte. Gants, lunettes. Benzène : cancérogène (éviter inhalation/contact). Cyclohexène : inflammable, irritant. Déchets : bac halogénés/organiques.
\end{experience}

\newpage

\coursec{Corrigés des exercices}

\begin{corrige}[Exercice 1 — QCM : structure du benzène]
\begin{enumerate}
\item \textbf{a} et \textbf{c}. La molécule de benzène est \cle{plane} : les 6 atomes de carbone et les 6 atomes d'hydrogène sont dans un même plan. Chaque carbone est hybridé $sp^2$ (géométrie triangulaire plane, angles de \SI{120}{\degree}). La réponse b est fausse : aucun atome du benzène n'est tétraédrique (hybridation $sp^3$).
\item \textbf{b}. Les six électrons $\pi$ sont \cle{délocalisés} sur l'ensemble du cycle : c'est le phénomène de \emph{mésomérie} (ou conjugaison) qui confère au noyau benzénique sa stabilité exceptionnelle. On ne peut pas parler de trois doubles liaisons fixes.
\item \textbf{c}. Les six liaisons \ce{C-C} du benzène sont toutes identiques, de longueur \SI{139}{pm}, valeur \cle{intermédiaire} entre une liaison simple (\SI{154}{pm}) et une liaison double (\SI{134}{pm}). C'est une preuve expérimentale de la délocalisation électronique.
\item \textbf{b}. La réaction caractéristique du benzène est la \cle{substitution} électrophile : elle permet de conserver le noyau aromatique et sa stabilité. L'addition est possible mais exige des conditions drastiques (catalyseur, lumière) et détruit l'aromaticité.
\end{enumerate}
\end{corrige}

\begin{corrige}[Exercice 2 — Vrai ou faux, à justifier]
\begin{enumerate}
\item \textbf{FAUX.} Contrairement aux alcènes, le benzène ne décolore pas l'eau de brome à température ambiante et sans catalyseur. La délocalisation des électrons $\pi$ stabilise le noyau et empêche l'addition du dibrome. (Avec un catalyseur comme \ce{FeBr3}, on observe une \emph{substitution}, pas une addition.)
\item \textbf{VRAI.} En présence de nickel (catalyseur), sous pression et à chaud, le benzène additionne trois molécules de dihydrogène pour donner le cyclohexane : \ce{C6H6 + 3 H2 ->[Ni] C6H12}.
\item \textbf{VRAI.} La nitration du benzène s'effectue avec le \emph{mélange sulfonitrique} : acide nitrique + acide sulfurique concentrés. L'acide sulfurique, catalyseur, permet la formation de l'entité électrophile active (ion nitronium \ce{NO2+}).
\item \textbf{FAUX.} Le 1,2-diméthylbenzène est l'\emph{ortho}-xylène (substituants sur deux carbones voisins). Le \emph{méta}-xylène est le 1,3-diméthylbenzène.
\item \textbf{FAUX.} Les dérivés nitrés aromatiques sont souvent \textbf{toxiques} (le nitrobenzène passe à travers la peau) et parfois \textbf{explosifs} (TNT, acide picrique). Leur manipulation exige des précautions strictes : gants, blouse, hotte, petites quantités.
\end{enumerate}
\end{corrige}

\begin{corrige}[Exercice 3 — Texte à trous]
La molécule de benzène est \textbf{plane} : ses douze atomes sont dans un même plan. Ses six électrons $\pi$ sont \textbf{délocalisés} sur tout le cycle, ce qui confère au \textbf{noyau benzénique} une grande stabilité. C'est pourquoi le benzène préfère les réactions de \textbf{substitution}, qui conservent le cycle aromatique. La réaction du benzène avec l'acide nitrique est une \textbf{nitration} ; elle nécessite un \textbf{catalyseur} : l'acide sulfurique concentré. Dans le 1,2-dichlorobenzène, les deux atomes de chlore sont en position \textbf{ortho}.
\end{corrige}

\begin{corrige}[Exercice 4 — Définitions]
\begin{enumerate}
\item \textbf{Composé aromatique} : composé organique dont la molécule contient au moins un \cle{noyau benzénique} (cycle à six atomes de carbone portant six électrons $\pi$ délocalisés), ce qui lui confère une stabilité et une réactivité particulières.
\item \textbf{Substitution électrophile} : réaction au cours de laquelle un atome d'hydrogène du noyau benzénique est remplacé par un groupe électrophile (accepteur d'électrons, ex. \ce{NO2+}, \ce{Cl+}, \ce{R+}), le cycle aromatique étant conservé.
\item \textbf{Réaction de Friedel et Crafts} : substitution électrophile au cours de laquelle un groupe alkyle (ex. \ce{-CH3}, \ce{-C2H5}) est fixé sur le noyau benzénique par action d'un halogénure d'alkyle \ce{R-Cl} en présence de chlorure d'aluminium \ce{AlCl3} comme catalyseur.
\item \textbf{Positions ortho, méta, para} : pour un noyau benzénique disubstitué, on dit que les deux substituants sont en position \emph{ortho} (1,2) s'ils sont portés par deux carbones voisins ; \emph{méta} (1,3) s'ils sont séparés par un carbone ; \emph{para} (1,4) s'ils sont en regard l'un de l'autre (positions opposées du cycle).
\end{enumerate}
\end{corrige}

\begin{corrige}[Exercice 5 — Les xylènes]
\begin{enumerate}
\item Les trois xylènes ont tous la même formule brute : $\ce{C6H4(CH3)2}$, c'est-à-dire \textbf{\ce{C8H10}}. Ce sont des isomères de position.
\item Formules semi-développées (le noyau est représenté par un hexagone, le cercle symbolise les électrons délocalisés) :
\begin{popfigure}
\begin{tikzpicture}[scale=0.9]
% ortho
\begin{scope}
\draw[thick,popblue] (30:1) -- (90:1) -- (150:1) -- (210:1) -- (270:1) -- (330:1) -- cycle;
\draw[popblue] (0,0) circle (0.62);
\draw[thick,popdark] (90:1) -- (90:1.55) node[above]{\footnotesize \ce{CH3}};
\draw[thick,popdark] (30:1) -- (30:1.55) node[above right=-2pt]{\footnotesize \ce{CH3}};
\node[below] at (0,-1.4) {\footnotesize 1,2-diméthylbenzène};
\node[below,popink] at (0,-1.85) {\footnotesize \emph{ortho}-xylène};
\end{scope}
% meta
\begin{scope}[xshift=4.6cm]
\draw[thick,popblue] (30:1) -- (90:1) -- (150:1) -- (210:1) -- (270:1) -- (330:1) -- cycle;
\draw[popblue] (0,0) circle (0.62);
\draw[thick,popdark] (90:1) -- (90:1.55) node[above]{\footnotesize \ce{CH3}};
\draw[thick,popdark] (330:1) -- (330:1.55) node[below right=-2pt]{\footnotesize \ce{CH3}};
\node[below] at (0,-1.4) {\footnotesize 1,3-diméthylbenzène};
\node[below,popink] at (0,-1.85) {\footnotesize \emph{méta}-xylène};
\end{scope}
% para
\begin{scope}[xshift=9.2cm]
\draw[thick,popblue] (30:1) -- (90:1) -- (150:1) -- (210:1) -- (270:1) -- (330:1) -- cycle;
\draw[popblue] (0,0) circle (0.62);
\draw[thick,popdark] (90:1) -- (90:1.55) node[above]{\footnotesize \ce{CH3}};
\draw[thick,popdark] (270:1) -- (270:1.55) node[below]{\footnotesize \ce{CH3}};
\node[below] at (0,-1.85) {\footnotesize 1,4-diméthylbenzène};
\node[below,popink] at (0,-2.3) {\footnotesize \emph{para}-xylène};
\end{scope}
\end{tikzpicture}
\legende{Les trois isomères du xylène (diméthylbenzène) : ortho (1,2), méta (1,3) et para (1,4).}
\end{popfigure}

\item Nomenclature :
\begin{center}
\begin{tabularx}{\linewidth}{|l|X|l|}
\hline
\rowcolor{popblueL} \textbf{Isomère} & \textbf{Nom systématique} & \textbf{Nom courant} \\
\hline
Groupes \ce{CH3} sur carbones voisins & 1,2-diméthylbenzène & \emph{ortho}-xylène \\
\hline
Groupes \ce{CH3} séparés par un carbone & 1,3-diméthylbenzène & \emph{méta}-xylène \\
\hline
Groupes \ce{CH3} en regard & 1,4-diméthylbenzène & \emph{para}-xylène \\
\hline
\end{tabularx}
\end{center}
\end{enumerate}
\textbf{Point de vigilance :} on numérote le cycle de façon à obtenir les \emph{plus petits indices possibles} pour les substituants.
\end{corrige}

\begin{corrige}[Exercice 6 — Équations-bilans de substitution]
\begin{enumerate}
\item \textbf{Chloration (halogénation)} — catalyseur : chlorure de fer(III) \ce{FeCl3}, à température ambiante :
\[ \ce{C6H6 + Cl2 ->[FeCl3] C6H5Cl + HCl} \]
Produit organique : le \textbf{chlorobenzène}.
\item \textbf{Nitration} — mélange sulfonitrique, \ce{H2SO4} concentré comme catalyseur, chauffage modéré (\SI{50}{\celsius} à \SI{60}{\celsius}) :
\[ \ce{C6H6 + HNO3 ->[H2SO4] C6H5NO2 + H2O} \]
Produit organique : le \textbf{nitrobenzène}.
\item \textbf{Sulfonation} — oléum (acide sulfurique fumant, source de \ce{SO3}), à chaud :
\[ \ce{C6H6 + SO3 -> C6H5-SO3H} \]
Produit organique : l'\textbf{acide benzènesulfonique}.
\item \textbf{Alkylation de Friedel et Crafts} — catalyseur : chlorure d'aluminium \ce{AlCl3} anhydre :
\[ \ce{C6H6 + CH3-CH2-Cl ->[AlCl3] C6H5-CH2-CH3 + HCl} \]
Produit organique : l'\textbf{éthylbenzène}.
\end{enumerate}
\textbf{Point de vigilance :} dans toutes ces réactions, un atome d'hydrogène du noyau est \emph{remplacé} par le groupe entrant : le noyau benzénique (6 électrons $\pi$ délocalisés) est conservé. Ne jamais écrire de produit d'addition ici.
\end{corrige}

\begin{corrige}[Exercice 7 — Additions sur le benzène]
\begin{enumerate}
\item \textbf{Hydrogénation catalytique} (nickel, sous pression, à chaud) :
\[ \ce{C6H6 + 3 H2 ->[Ni] C6H12} \]
Le produit obtenu est le \textbf{cyclohexane}. Trois molécules de dihydrogène s'additionnent sur le cycle.
\item \textbf{Chloration par addition} (sous l'action de la lumière ultraviolette, sans catalyseur de substitution) :
\[ \ce{C6H6 + 3 Cl2 ->[lumi\grave{e}re] C6H6Cl6} \]
Le produit a pour formule brute \textbf{\ce{C6H6Cl6}} : c'est l'hexachlorocyclohexane (dont l'isomère $\gamma$, le lindane, fut utilisé comme insecticide).
\item Dans le benzène, les six électrons $\pi$ sont délocalisés sur tout le cycle : c'est cette délocalisation qui constitue le \cle{caractère aromatique} et qui assure la stabilité du noyau. Lors d'une addition, les doubles liaisons « disparaissent » : les six carbones deviennent hybridés $sp^3$, la molécule n'est plus plane et les électrons $\pi$ ne sont plus délocalisés. Le cycle obtenu (cyclohexane ou hexachlorocyclohexane) est un \emph{cycloalcane saturé} : le caractère aromatique est perdu. C'est pourquoi ces additions, défavorables, n'ont lieu que dans des conditions énergétiques (catalyseur, lumière, pression).
\end{enumerate}
\end{corrige}

\begin{corrige}[Exercice 8 — Benzène et cyclohexène face au brome]
\begin{enumerate}
\item \textbf{Test à l'eau de brome.} Dans deux tubes à essais, verser quelques millilitres de chaque liquide, puis ajouter goutte à goutte de l'eau de brome (solution orangée de \ce{Br2}) et agiter, à température ambiante et à l'abri de la lumière vive.
\begin{itemize}
\item \textbf{Cyclohexène} : la coloration orangée du brome \emph{disparaît immédiatement} (décoloration rapide).
\item \textbf{Benzène} : la coloration orangée \emph{persiste} (pas de réaction dans ces conditions).
\end{itemize}
Le flacon dont le contenu décolore l'eau de brome contient le cyclohexène ; l'autre contient le benzène.
\item Le cyclohexène possède une \emph{vraie} double liaison \ce{C=C} localisée, riche en électrons et très réactive : le dibrome s'y \emph{additionne} facilement, d'où la décoloration instantanée. Dans le benzène, les six électrons $\pi$ sont \cle{délocalisés} sur l'ensemble du cycle : cette délocalisation stabilise fortement la molécule. Une addition de \ce{Br2} détruirait cette stabilité ; elle est donc défavorisée et ne se produit pas dans les conditions du test. Le benzène ne réagit avec le dibrome que par \emph{substitution}, en présence d'un catalyseur (\ce{FeBr3}).
\item Équation-bilan de l'addition du dibrome sur le cyclohexène :
\[ \ce{C6H10 + Br2 -> C6H10Br2} \]
Le produit est le 1,2-dibromocyclohexane (le dibrome s'additionne sur les deux carbones de la double liaison).
\end{enumerate}
\textbf{Point de vigilance :} la persistance de la teinte orangée avec le benzène ne signifie pas « absence totale de réactivité » : cela signifie que le benzène ne fait pas d'\emph{addition} dans ces conditions.
\end{corrige}

\begin{corrige}[Exercice 9 — Combustion du benzène et du toluène]
\begin{enumerate}
\item Combustion complète du benzène :
\[ \ce{2 C6H6 + 15 O2 -> 12 CO2 + 6 H2O} \]
(Vérification : 12 C, 12 H et 30 O de chaque côté.)
\item Combustion complète du toluène :
\[ \ce{C7H8 + 9 O2 -> 7 CO2 + 4 H2O} \]
(Vérification : 7 C, 8 H et 18 O de chaque côté.)
\item Combustion de $m = \SI{15,6}{g}$ de benzène.
\begin{enumerate}
\item Masse molaire du benzène : $M(\ce{C6H6}) = 6 \times 12 + 6 \times 1 = \SI{78}{g.mol^{-1}}$.
\[ n(\ce{C6H6}) = \frac{m}{M} = \frac{15,6}{78} = \SI{0,20}{mol} \]
\item Tableau d'avancement (équation simplifiée par 2 : \ce{C6H6 + 7,5 O2 -> 6 CO2 + 3 H2O}) :
\begin{center}
\begin{tabularx}{\linewidth}{|l|X|X|X|X|}
\hline
\rowcolor{popblueL} \textbf{État} & \ce{C6H6} & \ce{O2} & \ce{CO2} & \ce{H2O} \\
\hline
Initial (mol) & $0,20$ & $n$ & $0$ & $0$ \\
\hline
En cours (mol) & $0,20 - x$ & $n - 7,5x$ & $6x$ & $3x$ \\
\hline
Final (mol) & $0$ & $0$ & $1,20$ & $0,60$ \\
\hline
\end{tabularx}
\end{center}
Combustion complète : $x_{max} = \SI{0,20}{mol}$, d'où :
\[ n(\ce{O2}) = 7,5 \times 0,20 = \SI{1,50}{mol} \]
\item Volume de dioxygène :
\[ V(\ce{O2}) = n \times V_m = 1,50 \times 24 = \SI{36}{L} \]
\item Quantité puis masse de dioxyde de carbone :
\[ n(\ce{CO2}) = 6 \times 0,20 = \SI{1,20}{mol} \]
$M(\ce{CO2}) = 12 + 2 \times 16 = \SI{44}{g.mol^{-1}}$, donc :
\[ m(\ce{CO2}) = 1,20 \times 44 = \SI{52,8}{g} \]
\end{enumerate}
\item Le benzène \ce{C6H6} contient 6 atomes de carbone pour seulement 6 atomes d'hydrogène : sa teneur massique en carbone vaut $\frac{72}{78} \approx 92\,\%$, contre $\frac{72}{86} \approx 84\,\%$ pour l'hexane \ce{C6H14}. Le benzène est donc beaucoup plus riche en carbone. Lors de la combustion, le dioxygène de l'air peine à oxyder tout ce carbone : une partie n'est pas transformée en \ce{CO2} et s'échappe sous forme de particules de carbone (suies), qui rendent la flamme jaune éclairante et fuligineuse.
\end{enumerate}
\textbf{Barème indicatif (10 pts) :} équations équilibrées 2 pts ; quantité de benzène 1 pt ; tableau d'avancement 2 pts ; $n(\ce{O2})$ et $V(\ce{O2})$ 2 pts ; $m(\ce{CO2})$ 1,5 pt ; explication des suies 1,5 pt.

\textbf{Points de vigilance :} ne pas oublier d'équilibrer avec des coefficients entiers ou fractionnaires cohérents ; utiliser le volume molaire \emph{uniquement} pour les gaz ; toujours préciser les unités dans chaque calcul.
\end{corrige}

\begin{corrige}[Exercice 10 — Préparation du nitrobenzène]
\begin{enumerate}
\item Il s'agit d'une \cle{substitution électrophile} (nitration) : un atome d'hydrogène du noyau benzénique est remplacé par un groupe nitro \ce{-NO2}. Le noyau aromatique est conservé, ce qui est la signature d'une substitution.
\item L'acide sulfurique concentré joue le rôle de \cle{catalyseur} : il réagit avec l'acide nitrique pour former l'ion nitronium \ce{NO2+}, l'entité électrophile qui attaque le noyau benzénique. Il est régénéré en fin de réaction et absorbe en outre l'eau formée.
\item 
\begin{enumerate}
\item Quantité de benzène :
\[ n(\ce{C6H6}) = \frac{m}{M} = \frac{78}{78} = \SI{1,0}{mol} \]
\item L'acide nitrique étant en excès, le benzène est le réactif limitant. D'après l'équation-bilan, $n(\ce{C6H5NO2}) = n(\ce{C6H6}) = \SI{1,0}{mol}$. D'où la masse théorique :
\[ m_{th} = n \times M = 1,0 \times 123 = \SI{123}{g} \]
\end{enumerate}
\item Rendement de la préparation :
\[ \eta = \frac{m'}{m_{th}} \times 100 = \frac{98,4}{123} \times 100 = 80\,\% \]
Le rendement de la préparation est de \textbf{80 \%}.
\item Précautions (deux parmi les suivantes) :
\begin{itemize}
\item manipuler sous \textbf{hotte ventilée}, car les vapeurs de benzène sont toxiques et cancérigènes ;
\item porter \textbf{gants, blouse et lunettes} de protection (le nitrobenzène pénètre la peau) ;
\item travailler \textbf{loin de toute flamme} (benzène très inflammable) et verser les acides concentrés avec précaution, en refroidissant le mélange réactionnel (réaction exothermique) ;
\item ne manipuler les dérivés nitrés qu'en \textbf{petites quantités}.
\end{itemize}
\end{enumerate}
\textbf{Barème indicatif (10 pts) :} type de réaction justifié 1,5 pt ; rôle de \ce{H2SO4} 1,5 pt ; $n(\ce{C6H6})$ 1 pt ; masse théorique 2 pts ; rendement 2 pts ; sécurité 2 pts.

\textbf{Points de vigilance :} le rendement se calcule toujours par rapport à la masse (ou quantité) \emph{théorique} ; vérifier que le réactif limitant est bien identifié (ici le benzène, car \ce{HNO3} est en excès).
\end{corrige}

\begin{corrige}[Exercice 11 — Du benzène au TNT : le trinitrotoluène]
\begin{enumerate}
\item Formules brutes : toluène \textbf{\ce{C7H8}} (soit \ce{C6H5-CH3}) ; TNT \textbf{\ce{C7H5N3O6}}.
\item La trinitration remplace trois atomes d'hydrogène du noyau par trois groupes \ce{-NO2}, avec formation de trois molécules d'eau :
\[ \ce{C7H8 + 3 HNO3 ->[H2SO4] C7H5N3O6 + 3 H2O} \]
(Vérification : 7 C, 14 H, 3 N et 9 O de chaque côté.)
\item Formule semi-développée du 2,4,6-trinitrotoluène (groupe méthyle en position 1) :
\begin{popfigure}
\begin{tikzpicture}[scale=1]
\draw[thick,popblue] (30:1) -- (90:1) -- (150:1) -- (210:1) -- (270:1) -- (330:1) -- cycle;
\draw[popblue] (0,0) circle (0.62);
\draw[thick,popdark] (90:1) -- (90:1.6) node[above]{\footnotesize \ce{CH3}};
\draw[thick,popink] (30:1) -- (30:1.6) node[above right=-2pt]{\footnotesize \ce{NO2}};
\draw[thick,popink] (330:1) -- (330:1.6) node[below right=-2pt]{\footnotesize \ce{NO2}};
\draw[thick,popink] (270:1) -- (270:1.6) node[below]{\footnotesize \ce{NO2}};
\node[below,popmuted] at (0,-2.05) {\footnotesize positions 2, 4 et 6};
\end{tikzpicture}
\legende{Structure du 2,4,6-trinitrotoluène (TNT). Le groupe méthyle est en position 1, les groupes nitro en positions 2 (ortho), 4 (para) et 6 (ortho).}
\end{popfigure}

\item Par rapport au groupe méthyle (position 1) :
\begin{itemize}
\item les groupes nitro en positions \textbf{2 et 6} sont en position \emph{\textbf{ortho}} (carbones voisins du carbone 1) ;
\item le groupe nitro en position \textbf{4} est en position \emph{\textbf{para}} (carbone opposé au carbone 1).
\end{itemize}
Aucun groupe nitro n'est en position \emph{méta} (positions 3 et 5, restées hydrogénées).
\item 
\begin{enumerate}
\item Quantité de toluène : $m = \SI{46}{kg} = \SI{46000}{g}$.
\[ n(\ce{C7H8}) = \frac{m}{M} = \frac{46000}{92} = \SI{500}{mol} \]
\item D'après l'équation-bilan, $n(\text{TNT}) = n(\ce{C7H8}) = \SI{500}{mol}$. D'où la masse théorique :
\[ m(\text{TNT}) = n \times M = 500 \times 227 = \SI{113500}{g} = \SI{113,5}{kg} \]
\end{enumerate}
\item \textbf{Paragraphe de sécurité (proposition de réponse) :} Les dérivés nitrés aromatiques sont des composés dangereux : beaucoup sont toxiques par contact cutané ou par inhalation, et les polynitrés comme le TNT sont des explosifs puissants, sensibles aux chocs, à la chaleur et aux flammes. Leur manipulation exige donc des mesures strictes. Au laboratoire, on travaillera toujours avec de \emph{très petites quantités} de produit, sous \emph{hotte ventilée}, en portant \emph{gants, blouse et lunettes de protection}. On conservera ces composés \emph{à l'écart de toute source de chaleur, flamme ou étincelle}, et on évitera tout choc mécanique (pas de broyage ni de frottement). Enfin, les réactions de nitration, fortement exothermiques, seront conduites avec un \emph{refroidissement} du milieu réactionnel et une addition lente des réactifs.
\end{enumerate}
\textbf{Barème indicatif (10 pts) :} formules brutes 1 pt ; équation-bilan 2 pts ; formule semi-développée 1,5 pt ; positions relatives 1,5 pt ; $n(\ce{C7H8})$ 1 pt ; masse de TNT 1,5 pt ; paragraphe de sécurité 1,5 pt.

\textbf{Points de vigilance :} bien convertir les kilogrammes en grammes avant d'utiliser la masse molaire en \si{g.mol^{-1}} ; dans l'équation-bilan, ne pas oublier les \textbf{trois} molécules d'eau ; pour la nomenclature, le carbone portant le groupe méthyle est toujours numéroté 1.
\end{corrige}

\newpage

\coursec{Fiche bilan}

\begin{popfigure}
\centering
\begin{tikzpicture}[
  mindmap,
  grow cyclic,
  every node/.style={concept, font=\small},
  root concept/.append style={concept color=popblue, font=\bfseries\small, minimum size=2.6cm},
  level 1/.append style={level distance=3.6cm, sibling angle=51.4},
  level 2/.append style={level distance=2.2cm, sibling angle=40, font=\scriptsize},
  concept connection/.append style={color=popmuted}
]
\node[root concept]{Composés aromatiques}
  child[concept color=poporange]{ node {Structure du benzène}
    child { node {6 C, hexagone plan} }
    child { node {6 électrons $\pi$ délocalisés} }
  }
  child[concept color=popgreen]{ node {Nomenclature}
    child { node {ortho, méta, para} }
    child { node {toluène, phénol\ldots} }
  }
  child[concept color=poppurple]{ node {Combustion}
    child { node {flammes fuligineuses} }
  }
  child[concept color=popgold]{ node {Additions}
    child { node {$+\,3\,\ce{H2}$ (Ni)} }
    child { node {$+\,3\,\ce{Cl2}$ (UV)} }
  }
  child[concept color=poppink]{ node {Substitutions}
    child { node {halogénation} }
    child { node {nitration, sulfonation} }
    child { node {Friedel--Crafts} }
  }
  child[concept color=popblue!60]{ node {Sécurité et usages}
    child { node {toxique, cancérigène} }
    child { node {TNT, insecticides} }
  };
\end{tikzpicture}
\legende{Carte mentale de la leçon : les six grandes idées à retenir sur les composés aromatiques.}
\end{popfigure}

\begin{bilan}[L'essentiel de la leçon]
\textbf{Définitions clés}
\begin{itemize}
  \item \cle{Composé aromatique} : composé possédant au moins un \cle{noyau benzénique} (cycle à 6 atomes de carbone).
  \item \cle{Benzène} \ce{C6H6} : molécule \textbf{plane}, hexagone régulier ; les 6 électrons $\pi$ sont \textbf{délocalisés} sur tout le cycle ; toutes les liaisons C--C sont identiques (longueur intermédiaire entre simple et double liaison).
  \item Représentation symbolique : hexagone avec un \textbf{cercle} (électrons délocalisés) ou trois doubles liaisons alternées (formule de Kekulé).
  \item Isomères de position des dérivés disubstitués : \cle{ortho} (1,2), \cle{méta} (1,3), \cle{para} (1,4).
\end{itemize}

\textbf{Équations-types à connaître par cœur}
\begin{itemize}
  \item Combustion complète : \[ \ce{2 C6H6 + 15 O2 -> 12 CO2 + 6 H2O} \quad \text{(flamme très fuligineuse)} \]
  \item Addition d'hydrogène (cyclohexane) : \[ \ce{C6H6 + 3 H2 ->[Ni] C6H12} \]
  \item Addition de dichlore (sous UV) : \[ \ce{C6H6 + 3 Cl2 ->[UV] C6H6Cl6} \]
  \item Chloration (substitution) : \[ \ce{C6H6 + Cl2 ->[FeCl3] C6H5Cl + HCl} \]
  \item Nitration : \[ \ce{C6H6 + HNO3 ->[\ce{H2SO4}\ conc.] C6H5NO2 + H2O} \]
  \item Sulfonation : \[ \ce{C6H6 + H2SO4 -> C6H5SO3H + H2O} \]
  \item Friedel--Crafts : \[ \ce{C6H6 + CH3Cl ->[AlCl3] C6H5CH3 + HCl} \]
\end{itemize}

\textbf{Quelques composés aromatiques à savoir nommer}
\begin{center}
\begin{tabularx}{\linewidth}{|l|X|X|}
\hline
\rowcolor{popblueL} \textbf{Nom} & \textbf{Formule} & \textbf{Utilisation} \\
\hline
Benzène & \ce{C6H6} & solvant, synthèses industrielles \\
\hline
Toluène & \ce{C6H5-CH3} & solvant, préparation du TNT \\
\hline
Phénol & \ce{C6H5-OH} & désinfectant, résines \\
\hline
Nitrobenzène & \ce{C6H5-NO2} & intermédiaire de synthèse \\
\hline
TNT (trinitrotoluène) & \ce{C6H2(NO2)3CH3} & explosif \\
\hline
Acide benzènesulfonique & \ce{C6H5-SO3H} & détergents \\
\hline
\end{tabularx}
\end{center}

\textbf{Méthodes express}
\begin{itemize}
  \item \textbf{Reconnaître la réaction} : un H du noyau est remplacé $\Rightarrow$ \cle{substitution} (produit + \ce{HCl} ou \ce{H2O}) ; le cycle disparaît ou s'additionne $\Rightarrow$ \cle{addition} (\ce{H2}/Ni, \ce{Cl2}/UV).
  \item \textbf{Pourquoi substitution plutôt qu'addition ?} L'addition détruirait la délocalisation des électrons $\pi$, très stabilisante : le benzène préfère \textbf{conserver son noyau aromatique}.
  \item \textbf{Nomenclature disubstituée} : numéroter le cycle pour obtenir les indices les plus petits ; 1,2 = ortho ; 1,3 = méta ; 1,4 = para.
  \item \textbf{Sécurité} : benzène et dérivés nitrés manipulés sous hotte, avec gants et lunettes, loin de toute flamme (toxiques, cancérigènes, inflammables).
\end{itemize}
\end{bilan}

\begin{aretenir}[Checklist avant le Probatoire]
\begin{itemize}
  \item[$\square$] Je sais dessiner le noyau benzénique (hexagone + cercle) et expliquer la délocalisation des 6 électrons $\pi$.
  \item[$\square$] Je sais que le benzène est plan et que ses 6 liaisons C--C sont identiques.
  \item[$\square$] Je sais écrire et équilibrer la combustion complète du benzène.
  \item[$\square$] Je sais écrire l'addition de \ce{H2} (catalyseur Ni) et de \ce{Cl2} (sous UV) sur le benzène.
  \item[$\square$] Je sais écrire les quatre substitutions : chloration, nitration, sulfonation, Friedel--Crafts, avec leurs catalyseurs.
  \item[$\square$] Je sais expliquer pourquoi le benzène privilégie la substitution à l'addition.
  \item[$\square$] Je sais nommer les trois isomères disubstitués : ortho, méta, para.
  \item[$\square$] Je sais nommer et écrire les formules du toluène, du phénol, du nitrobenzène et du TNT.
  \item[$\square$] Je sais citer les mesures de sécurité pour manipuler le benzène et les dérivés nitrés.
  \item[$\square$] Je sais distinguer, dans un énoncé, une réaction d'addition d'une réaction de substitution.
\end{itemize}
\end{aretenir}

\findecours

\end{document}
